%!TEX program = xelatex
% =============================================================================
%  2026 高教社杯全国大学生数学建模竞赛  C 题
%  《微网与外部电网电力调控策略》论文 —— 主文件
%
%  编译：./build.sh        （或 xelatex main.tex 连编两遍）
%  定稿：把 preamble.tex 中 \drafttrue 改为 \draftfalse，再跑 tools/check_paper.py
%
%  目录结构
%     main.tex                 本文件（只做装配，不写正文；产出 main.pdf = 摘要 + 正文 + 参考文献）
%     full.tex                 另一入口：正文之后再接附录（含源程序清单），产出 full.pdf
%     preamble.tex             导言区：格式合规配置 + 写作宏
%     sections/00_abstract.tex 摘要专用页（自本页起页码为 1）
%     sections/01..10          正文十章
%     sections/refs.tex        参考文献（GB/T 7714）
%     sections/90_appendix.tex 附录（仅由 full.tex 引入，不进 main.pdf）
%     figures/                 图片（fig01/fig02 为 TikZ 源文件，其余为 PDF）
%     code/                    附录所列源程序（由 lstinputlisting 原样引入）
%     tools/check_paper.py     提交前合规自查
%
%  ⚠️ 全文不得出现学校、队号、姓名、赛区信息（格式规范第六条）
%  ⚠️ 不要加入 \tableofcontents（格式规范第四条：不要目录）
% =============================================================================
\documentclass[10pt,a4paper]{ctexart}
%!TEX root = main.tex
% =============================================================================
%  preamble.tex —— 2026 CUMCM C 题论文模板导言区
%  格式依据：《全国大学生数学建模竞赛论文格式规范（2026 年修订稿）》format2026.doc
%   第一条 A4、页边距≥2.5cm、左侧装订        -> geometry
%   第三条 摘要≤1页、页码自摘要页起、页脚中部、阿拉伯数字连续 -> fancyhdr
%   第四条 正文≤30页、不要目录                -> 不加载 \tableofcontents
%   第六条 摘要页/正文/附录不得出现身份信息    -> 不设 \author，PDF 元数据留空
%   第七条 引用需列参考文献并在正文标注        -> thebibliography（GB/T 7714）
% =============================================================================

% ---------------------------------------------------------------------------
% 1. 页面与版式  —— 第一条：A4，上下左右各留 ≥2.5cm，左侧装订
% ---------------------------------------------------------------------------
\usepackage[
  a4paper,
  top=2.5cm, bottom=2.5cm, left=2.5cm, right=2.5cm,
  footskip=1.0cm      % 页脚落在下边距内，正文块本身仍保有 2.5cm 边距
]{geometry}

% 行距：正文约 1.2 倍（格式规范第八条不对字号、行距作统一要求）。
% 正文有 30 页硬上限，此处取值已为「表格与图填入真实结果」留出余量；
% 若确实需要更松的行距，请先确认正文仍不超过 30 页。
\linespread{1.0}
% 允许 TeX 在极窄的表格列中略微拉伸以消除溢出（不影响观感）
\setlength{\emergencystretch}{2em}
% 段间距为 0（中文论文用首行缩进分段）；浮动体与标题间距一并收紧，
% 以便正文容纳更多内容而不牺牲可读性
\setlength{\parskip}{0pt}
\setlength{\textfloatsep}{9pt plus 2pt minus 2pt}
\setlength{\intextsep}{9pt plus 2pt minus 2pt}
\setlength{\floatsep}{9pt plus 2pt minus 2pt}
\setlength{\abovecaptionskip}{4pt}
\setlength{\belowcaptionskip}{0pt}

% ---------------------------------------------------------------------------
% 2. 页码  —— 第三条：自摘要页起、页脚中部、阿拉伯数字从 1 连续编号
%             摘要页即电子版第 1 页（第十条：承诺书与编号页不放电子版）
% ---------------------------------------------------------------------------
\usepackage{fancyhdr}
\pagestyle{fancy}
\fancyhf{}
\fancyfoot[C]{\normalfont\thepage}
\renewcommand{\headrulewidth}{0pt}
\renewcommand{\footrulewidth}{0pt}
\fancypagestyle{plain}{%
  \fancyhf{}\fancyfoot[C]{\normalfont\thepage}%
  \renewcommand{\headrulewidth}{0pt}\renewcommand{\footrulewidth}{0pt}%
}
\pagenumbering{arabic}

% ---------------------------------------------------------------------------
% 3. 数学、表格、图形
% ---------------------------------------------------------------------------
\usepackage{amsmath,amssymb,amsthm,mathrsfs,bm}
\allowdisplaybreaks[2]

% booktabs 的横线间距一并收紧，减少表格占用的纵向空间
\usepackage{booktabs,longtable,tabularx,array,multirow,makecell}
\setlength{\tabcolsep}{4pt}
\setlength{\abovedisplayskip}{4pt plus 2pt minus 2pt}
\setlength{\belowdisplayskip}{4pt plus 2pt minus 2pt}
\setlength{\abovedisplayshortskip}{2pt plus 1pt}
\setlength{\belowdisplayshortskip}{2pt plus 1pt}
\setlength{\aboverulesep}{0.18ex}
\setlength{\belowrulesep}{0.18ex}
\newcolumntype{C}[1]{>{\centering\arraybackslash}p{#1}}
\newcolumntype{L}[1]{>{\raggedright\arraybackslash}p{#1}}
\newcolumntype{Y}{>{\centering\arraybackslash}X}
\renewcommand{\arraystretch}{0.95}
% 浮动体（图/表）排布参数：本文图表较多，放宽参数可显著减少"半空页"
\usepackage{graphicx}
\setcounter{topnumber}{5}
\setcounter{bottomnumber}{3}
\setcounter{totalnumber}{8}
\renewcommand{\topfraction}{0.92}
\renewcommand{\bottomfraction}{0.85}
\renewcommand{\textfraction}{0.03}
\renewcommand{\floatpagefraction}{0.85}
\renewcommand{\dbltopfraction}{0.92}

\usepackage{subcaption}
\usepackage[font=small,labelsep=quad,skip=4pt]{caption}
\usepackage{enumitem}
\setlist{nosep,leftmargin=2.2em}
\usepackage{xcolor}
\usepackage{tikz}
\usetikzlibrary{arrows.meta,positioning,calc,patterns,shapes.geometric,fit,backgrounds,decorations.pathreplacing,decorations.markings}
\usepackage{pgfplots}
\pgfplotsset{compat=1.16}
\pgfplotsset{every axis/.append style={font=\small}}

% 算法
\usepackage{algorithm}
\usepackage{algpseudocode}
\renewcommand{\algorithmicrequire}{\textbf{输入：}}
\renewcommand{\algorithmicensure}{\textbf{输出：}}
\floatname{algorithm}{算法}

% ---------------------------------------------------------------------------
% 4. 定理类环境（中文名）
% ---------------------------------------------------------------------------
\theoremstyle{plain}
\newtheorem{proposition}{命题}
\newtheorem{lemma}{引理}
\newtheorem{corollary}{推论}
\theoremstyle{definition}
\newtheorem{definition}{定义}
\renewcommand{\proofname}{\heiti 证明}
\renewcommand{\qedsymbol}{$\square$}

% ---------------------------------------------------------------------------
% 5. 章节标题样式：一级「一、」三号黑体居中；二级 四号黑体；三级 小四黑体
% ---------------------------------------------------------------------------
\ctexset{
  section = {
    format     = \zihao{3}\heiti\centering,
    name       = {,、},
    number     = \chinese{section},
    aftername  = {},
    beforeskip = 1.2ex plus .4ex minus .2ex,
    afterskip  = 0.8ex plus .2ex,
  },
  subsection = {
    format     = \zihao{4}\heiti\raggedright,
    number     = \arabic{section}.\arabic{subsection},
    aftername  = \quad,
    beforeskip = 1.0ex plus .4ex minus .2ex,
    afterskip  = 0.5ex plus .2ex,
  },
  subsubsection = {
    format     = \zihao{-4}\heiti\raggedright,
    number     = \arabic{section}.\arabic{subsection}.\arabic{subsubsection},
    aftername  = \quad,
    beforeskip = 1.0ex plus .4ex minus .2ex,
    afterskip  = 0.5ex plus .2ex,
  },
  paragraph = {
    format     = \zihao{-4}\heiti\raggedright,
    aftername  = \quad,
  },
}

% ---------------------------------------------------------------------------
% 6. 源码清单（附录）—— 关键：basicstyle 中加入 \songti，
%    使代码中的中文注释用已嵌入的中文字体排版；否则中文会丢失或字体不嵌入。
% ---------------------------------------------------------------------------
\usepackage{listings}
\definecolor{cKeyword}{RGB}{0,80,160}
\definecolor{cComment}{RGB}{0,120,80}
\definecolor{cString} {RGB}{170,60,0}
\definecolor{cGray}   {RGB}{120,120,120}
\lstset{
  basicstyle       = \ttfamily\footnotesize\songti,
  keywordstyle     = \color{cKeyword}\bfseries,
  commentstyle     = \color{cComment}\itshape,
  stringstyle      = \color{cString},
  numberstyle      = \tiny\color{cGray},
  numbers          = left,
  numbersep        = 7pt,
  frame            = single,
  framerule        = 0.4pt,
  rulecolor        = \color{cGray},
  breaklines       = true,
  breakatwhitespace= false,
  columns          = flexible,
  keepspaces       = true,
  showstringspaces = false,
  tabsize          = 4,
  xleftmargin      = 2.2em,
  framexleftmargin = 2.2em,
  aboveskip        = 1.2ex,
  belowskip        = 1.2ex,
}
\lstdefinestyle{py}{language=Python}
\lstdefinestyle{plain}{language=}

% ---------------------------------------------------------------------------
% 7. 超链接与 PDF 元数据（第六条：元数据不得含身份信息）
% ---------------------------------------------------------------------------
% 黑体改用 PingFang SC：STHeiti 在 macOS 下 ToUnicode 生成失败会导致 PDF 文字无法复制
% 正文与黑体均启用伪粗体：macOS 的 Songti SC Bold 不嵌入 ToUnicode，
% 直接用 \textbf 会让加粗中文在复制时变成乱码；伪粗体由常规字重合成，可正常复制。
\setCJKmainfont{Songti SC}[AutoFakeBold=2.2]
\setCJKsansfont{Songti SC}[AutoFakeBold=2.2]
\setCJKmonofont{Songti SC}[AutoFakeBold=2.2]
\setCJKfamilyfont{zhhei}{PingFang SC}[BoldFont=PingFang SC Semibold]
\renewcommand{\heiti}{\CJKfamily{zhhei}}
\usepackage[hidelinks,bookmarks=true,bookmarksnumbered=false]{hyperref}
\hypersetup{
  pdfauthor  = {},
  pdfcreator = {},
  pdfsubject = {},
  pdftitle   = {微网与外部电网电力调控策略},
}

% ---------------------------------------------------------------------------
% 8. 写作期与定稿期开关
%    \drafttrue  ：显示【写作提示】与红色 {待计算} 占位符（写作期默认）
%    \draftfalse ：隐藏全部写作提示；占位符变为醒目红底，便于定稿前扫尾
%    两种进入定稿模式的方式（任选）：
%      (a) 把下面一行改成 \draftfalse
%      (b) 在论文目录下执行  touch final.flag   （便于反复切换与自动检查）
%    定稿前务必运行 tools/check_paper.py
% ---------------------------------------------------------------------------
\newif\ifdraft
\drafttrue
\IfFileExists{final.flag}{\draftfalse}{}

\definecolor{cGuide}{RGB}{90,90,90}
\definecolor{cNote} {RGB}{0,90,180}
\definecolor{cPh}   {RGB}{200,0,0}
\definecolor{cEvid} {RGB}{120,60,150}

\usepackage{environ}

% 带圈数字（①②③…）、罗马数字等符号在默认字体分类下会被交给西文字体而缺字，
% 这里把它们划入 CJK 字符类，交给中文字体排版（仅 XeLaTeX + xeCJK 生效）
\ifdefined\xeCJKsetcharclass
  \xeCJKsetcharclass{"2460}{"24FF}{1}   % ①②③…⒈⒉⒊…
  \xeCJKsetcharclass{"2160}{"217F}{1}   % Ⅰ Ⅱ Ⅲ … ⅰ ⅱ ⅲ
  \xeCJKsetcharclass{"3000}{"303F}{1}   % CJK 标点
\fi

% 【写作提示】块：draft 模式显示，定稿模式整块消失
\NewEnviron{guideblock}{%
  \ifdraft
    \par\medskip\noindent{\color{cGuide}\rule{\linewidth}{0.7pt}}\par\nobreak
    \begingroup\small\color{cGuide}\noindent\textbf{【写作提示】}\ignorespaces\BODY\par\endgroup
    \noindent{\color{cGuide}\rule{\linewidth}{0.7pt}}\par\medskip
  \fi
}

% 行内短注
\newcommand{\note}[1]{\ifdraft{\color{cNote}\small【注：#1】}\fi}

% 未算出的数字：一律用 \ph 占位，严禁编造数字
% 正文中允许在占位符内断行（避免长占位符把行撑出边界），数学模式中则用 \hbox 保持整体
\newcommand{\ph}[1][待计算]{%
  \ifmmode\hbox{\color{cPh}\bfseries 【#1】}%
  \else{\color{cPh}\bfseries 【#1】}\fi}

% 证据编号（证据账本 C-Qn-xx）
\newcommand{\evid}[1]{\ifdraft{\color{cEvid}(#1)}\else（见 #1）\fi}

% 口径编号（口径框 C-01…C-08）
\newcommand{\calref}[1]{口径 #1}

% 口径框（紧凑版：小字号、无首行缩进，正文用「定义/数据/时段/单位/不含/对照/复算」逐项写）
\NewEnviron{caliber}[2]{%
  \par\smallskip\noindent{\color{cEvid}\rule{\linewidth}{0.7pt}}\par\nobreak
  {\noindent\heiti 口径 #1\quad #2\par}\nobreak
  \begingroup\footnotesize\setlength{\parindent}{0pt}\BODY\par\endgroup
  \noindent{\color{cEvid}\rule{\linewidth}{0.7pt}}\par\smallskip
}

% 创新五段式
\NewEnviron{innovation}[2]{%
  \par\medskip\noindent{\heiti 创新 #1\quad #2}\par\nobreak\vspace{0.2ex}
  \begingroup\BODY\par\endgroup\medskip
}
\newcommand{\ip}[1]{\par\noindent\textbf{#1}\quad\ignorespaces}

% 插图：图片不存在时画红框占位，保证任何阶段都能编译
\newcommand{\includefig}[2][width=\linewidth]{%
  \IfFileExists{figures/#2}%
    {\includegraphics[#1]{figures/#2}}%
    {\fbox{\begin{minipage}[c][2.6cm][c]{0.9\linewidth}\centering\color{cPh}\small
      【缺图】\texttt{\detokenize{figures/#2}}\par\vspace{0.6ex}
      待由 \texttt{C题代码/code/plot\_fig*.py} 生成后放入 \texttt{figures/}\end{minipage}}}%
}

% 附录源程序：自动原样引入 code/ 目录下的真实文件，
% 保证"附录代码 = 真正跑出结果的代码"；文件缺失时给出红框提示而不中断编译。
\newcommand{\codefile}[2][py]{%
  \par\medskip\noindent{\heiti 源程序：\texttt{\detokenize{code/#2}}}\par\nobreak
  \IfFileExists{code/#2}%
    {\lstinputlisting[style=#1]{code/#2}}%
    {\begin{center}\fcolorbox{cPh}{white}{\parbox{0.92\linewidth}{\centering\color{cPh}\small
       【缺文件】\texttt{\detokenize{code/#2}}\par\vspace{0.4ex}
       请把源程序放入 \texttt{code/} 目录，附录会\textbf{自动原样引入}（不要用截图）}}\end{center}}%
  \par\medskip
}

% 论文标题（摘要页与 PDF 元数据共用，只改这一处）
\newcommand{\papertitle}{微网与外部电网电力调控策略}

% 单位（正体，与变量斜体区分）
\newcommand{\unit}[1]{\textup{#1}}

% 正文中引用题面表格
\newcommand{\tabref}[1]{表~#1}


\begin{document}

% ---- 摘要专用页（电子版第 1 页，页码 1） ----------------------------------
%!TEX root = ../main.tex
% =============================================================================
%  摘要专用页（格式规范第三条：含标题与关键词，无需英文，原则上不超过一页；
%              论文自本页开始编页码，页脚中部，阿拉伯数字从 1 起）
%  本页全部数字取自 intermediate/paper_numbers.json 与 v2 新增对照实验
%  写法：背景一句 → 逐问"结构理由 + 模型 + 量化结果" → 验证一句 → 结论与边界一句
% =============================================================================
{\linespread{1.35}\selectfont
\begin{center}
  {\zihao{3}\heiti \papertitle\par}
\end{center}

\begin{center}
  {\zihao{4}\heiti 摘\quad 要}
\end{center}

并网微网在分时电价与光伏出力双重不确定性下运行：光伏盈余落在电价低谷的午间，负载高峰落在电价尖峰的
傍晚，而储能容量与充放功率均有限，收益只能来自把电量从低价时段搬到高价时段。本文按信息结构由弱到强给
出四层输出，每层只增加一个结构性对象，并用三组对照实验检验其结论。

\textbf{问题 1} 中电价、负载与光伏预测给定且要求两端储电量相等，目标与约束均为线性，故建立
\textbf{线性规划}模型；由\textbf{对偶变量}给出储电水价值，在储电量界不活跃时由互补松弛得到最优充放电的
\textbf{阈值型刻画}（死区比 \(1/(\eta_c\eta_d)=1.2346\)）。最优全天购电费 \textbf{35\,126.95 元}，
相对无储能基线 \(48\,052.05\) 元节省 \(26.90\%\)，充、放电量 \(20\,740.67\) 与 \(16\,799.94\)~kWh
（比值 \(1.2346\) 与往返效率互为倒数），弃光为 0；容量摆幅是主要瓶颈（解除容量上限可再省
\(22.14\%\)，解除功率上限仅 \(1.66\%\)）。

\textbf{问题 2} 的计划依据为 \(0{:}00\) 光伏预报、结算依据为当日实际值，缺口按 \(5\) 倍电价紧急购电，
故建立\textbf{两阶段随机规划}模型（场景集按日总残差分层取 medoid、一次存盘固定），并用
\textbf{报童型临界分位数}给出安全边际的解析近似（缺货概率两界 \(1/10\)--\(1/5\)）。334 天总费用
\textbf{14\,264\,792.13 元}，紧急购电量 \(172\,015.82\)~kWh、出现在 296 天，节省 \(13.06\%\)；
安全边际的\textbf{边际价值 3\,133\,767.55 元}（消融：移除后费用升至 \(17\,398\,559.68\) 元、
紧急购电量增至 \(6.72\) 倍）。

\textbf{问题 3} 有四个预报时刻，故建立含\textbf{非预期性}的\textbf{多阶段滚动优化}模型（分段线性凸
偏差费用被精确线性化为 LP，无需 0--1 变量）。费用降至 \textbf{13\,932\,402.95 元}，信息价值
\(332\,389.18\) 元（\(2.33\%\)），已实现完全信息上界价值的 \(85.1\%\)；\textbf{再加密最多再省
\(0.42\%\)}，故不建议继续加密预报时刻，应把投入用于提高现有四次预报（尤其上午）的精度。

\textbf{问题 4} 中电价本身随机：先用\textbf{C 藤 pair-Copula}刻画相依结构（电价—负载
\(\tau=+0.698\)；电价—光伏由自由度 \(3.23\) 的 \(t\) Copula 给出 \(\tau=-0.238\)、对称尾相依
\(0.0479\)，Gaussian 族会把它抹平），再用\textbf{分时骨架加随机残差}两层结构建模电价，用保留联合
结构的\textbf{联合情景}与\textbf{条件风险价值}重算问题 2、3：总费用 \textbf{15\,094\,255.84 元}与
\textbf{14\,722\,153.17 元}，节省 \(11.35\%\) 与 \(13.54\%\)；电价波动使费用上升 \(5.81\%/5.67\%\)
且不能由价格水平解释，说明\textbf{高价与高净负荷的共动}才是成本来源。

\textbf{三组对照给出否证性结论。} ① 四种情景生成器在样本外 184 天上的总费用极差仅 \(0.37\%\)，故主
结果对生成范式稳健，但紧急购电量相差 \(48.5\%\)，引用该分项必须带情景口径。② 机会约束最优档
（\(\beta=0.10\)）仍比随机规划贵 \(0.52\%\)，且实测缺电频率约为标称 \(\beta\) 之半，与报童两界互证。
③ 盒式分布鲁棒最优档在\textbf{全部五项指标}上更差（总费用高 \(11.43\%\)、\(\mathrm{CVaR}\) 高
\(8.72\%\)），其价值只在于给出最坏情形上界。分时段风险画像进一步显示风险是结构性的：费用尾部集中在
\(15\)--\(18\) 时，缺电频率峰值出现在 \(11\) 时（\(23.8\%\)）与 \(5\) 时。

\textbf{验证。} 全部结果通过原始可行性（供给缺口残差 \(3.41\times10^{-13}\)~kWh）、最优性（阈值律
一致性，144 槽违反次数为 0）与费用独立复算（相对误差 \(\le10^{-6}\)）三重检验；审计 A1--A10 在
10~min 级解上全部通过，越界、互斥与负购电计数均为 0；另用\textbf{L 形（Benders）分解}这条独立数值
路径复核最优值（334 天最大相对间隙 \(3.76\times10^{-15}\)）。

\textbf{结论与边界。} 在本题参数范围内储能的核心价值是\textbf{把午间光伏盈余与低谷电量搬运到傍晚高价
时段}；当峰谷价比高于盈亏平衡阈值 \(1/(\eta_c\eta_d)\approx1.235\) 时套利有利。结论适用于价格接受型、
不外送、单一储能的并网微网，且节省幅度按电能量购电费用口径计算，是\textbf{保守下界}；允许上网售电、
引入需求响应或计入容量与辅助服务收益时须重新标定。

\noindent{\heiti 关键词：}\quad 微网；储能调度；线性规划；两阶段随机规划；报童模型；多阶段滚动优化；
Copula 联合情景；CVaR；信息价值
}
\newpage


% ---- 正文（≤30 页，无目录） -----------------------------------------------
%!TEX root = ../main.tex
% =============================================================================
%  一、问题重述与分析
%  评审在问：你知道题目在问什么吗？
%  本节含：① 记号约定与决策矛盾（式 (1)）  ② 四问的信息结构与新增数学对象（表 1）
%          ③ 决策时序图（图 1）、依赖图（图 2）  ④ 两条基线的口径分离（表 2、表 3）
%  写法约定：本节起全文用「集合／参数／决策变量／对偶变量」四类记号；
%            完整符号表见 §2.3，公式中每个符号都在其中登记。
% =============================================================================

\section{问题重述与分析}\label{chap:problem}

\subsection{记号与决策矛盾}

全天离散为 \(|\mathcal T|=144\) 个时段，\(t\in\mathcal T\)，步长 \(\Delta t=1/6\)~h；结果窗口为
\(\mathcal D\)，\(|\mathcal D|=334\) 天（2025-02-01 至 12-31）。第 \(t\) 槽的电价、负载功率与光伏功率记为
\(p_t,L_t,PV_t\)。储能充、放电量 \(c_t,d_t\) 与储电量 \(S_t\) 满足
\(0\le c_t,d_t\le E^{\max}=\overline P\Delta t\)、\(S_t\in[\underline S,\overline S]\)，
购电量非负 \(E_t\ge0\)（假设 A1：不外送）。

\textbf{决策矛盾。} 微网的收益只能来自电价的时间价差：把电量从低价时段搬到高价时段。储能搬运
\(1\)~kWh 完成一次循环须克服往返损耗，故套利有利可图的\textbf{盈亏平衡价比}为

\begin{equation}
  \frac{p_{\max}}{p_{\min}}\;>\;\frac{1}{\eta_c\eta_d}=1.2346 ,
  \label{eq:ch1-arb}
\end{equation}

而附件 1 的实测价比为 \(1.3952/0.3713=3.76\)，远高于该阈值——\textbf{价差不是瓶颈}。真正的约束是
容量摆幅、充放功率与两类不确定性（负载与光伏）。四问逐层引入这些约束：

\begin{itemize}
  \item \textbf{问题 1}（确定性计划）：\((p_t,L_t,PV_t)\) 已知且逐日重复，要求 \(S_0=S_{144}\)，
        在 \(0{:}00\) 一次制定全天计划；
  \item \textbf{问题 2}（日前计划）：\(p_t\) 仍固定，\((L_t,PV_t)\) 逐日随机；计划层在 \(0{:}00\) 决策，
        缺口在结算时刻按 \(5p_t\) 补足（\textbf{紧急购电}），构成两阶段决策；
  \item \textbf{问题 3}（滚动调整）：\(m\in\{0,6,12,18\}\) 时可获取未来 24~h 整点光伏预报，
        可在该时刻修正计划；计划高于调整的部分按 \(0.5p_t\) 计价、低于的部分按 \(1.5p_t\) 计价；
        追问是否需要引入其他时刻的预报；
  \item \textbf{问题 4}（波动电价）：\(p_t\) 亦随机，问题由“确定价格 \(+\) 随机出力”升级为
        \textbf{价格与出力的联合不确定性}，且须给出风险度量。
\end{itemize}

\subsection{一个模型族、两个开关}

四问共用同一套变量 \((E_t,c_t,d_t,S_t)\) 与同一组约束：电量平衡式~\eqref{eq:q1-balance}、
储电量递推式~\eqref{eq:q1-soc}、容量与功率界、以及假设 A1。它们仅由两个开关区分：
\textbf{电价向量}（附件 1 日内曲线／附件 4 实际序列）与\textbf{信息层数}
（1 个决策时刻／1 个决策时刻 \(+\) 情景／4 个决策时刻）。表~\ref{tab:fourq} 给出四问在该框架下的位置。

\begin{table}[htbp]
  \centering
  \caption{四问在“一个模型族、两个开关”下的位置（范围：问题 1 为附件 1 典型日，问题 2--4 为
           2025-02-01 至 12-31 共 334 天；单位与量纲见 \S\ref{sec:notation}；样本：4 个子问题）}
  \label{tab:fourq}
  \small
  \setlength{\tabcolsep}{4pt}
  \begin{tabularx}{\linewidth}{C{1.0cm} L{2.0cm} L{2.5cm} L{2.2cm} Y}
    \toprule
    子问题 & 电价向量 & 信息层数 & 目标结构 & 相对上一问\textbf{新增}的数学对象 \\
    \midrule
    1 & 附件 1 日内曲线 & 1 个决策时刻 & 确定性 LP & 对偶水价值 \(\nu_t\) 与阈值律（命题~\ref{prop:threshold}） \\
    2 & 附件 1 日内曲线 & 1 个决策时刻 \(+\) 情景 & 两阶段随机规划 & \textbf{情景维度}与追索量 \(g_{\omega,t}\)；报童临界分位数 \\
    3 & 附件 1 日内曲线 & 4 个决策时刻 & 多阶段（滚动实现） & \textbf{非预期性要求}式~\eqref{eq:q3-nonanticip}；信息价值 \\
    4 & 附件 4 实际序列 & 同上两种 & 同上 \(+\) 风险度量 & 电价随机过程、联合情景、\(\mathrm{CVaR}\) \\
    \bottomrule
  \end{tabularx}
\end{table}

末两列须强调：问题 2 相对问题 1 增加\textbf{情景维度}（计划量共享、追索量随情景变化），问题 3 相对
问题 2 增加\textbf{非预期性要求}与调整定价，二者均为结构性新增而非参数替换。决定建模方式的即三件事：
\textbf{决策时刻个数}、\textbf{每时刻的信息集}、\textbf{偏差的惩罚结构}（图~\ref{fig:timeline}）。
\textbf{一条全局声明}：问题 3（及问题 4-3）给出的多阶段模型\textbf{不展开场景树}，实际求解的是它的
滚动时域（MPC）可行策略，故该问的费用是\textbf{策略评价}而非多阶段最优值；理由、
边界与对信息价值结论的影响集中在 \S\ref{sec:q3-mpc}，并在摘要与\chapref{chap:eval}结论处重复标注。

\begin{figure}[htbp]
  \centering
  %!TEX root = ../main.tex
% =============================================================================
%  图 1  决策时序图（正文 §1.2）
%  规格（四项必须标注，缺一不可）：
%    ① "计划依据"与"结算依据"是两套数据，用不同线型区分
%    ② 18:00 的预报覆盖到次日 18:00，用一条跨越 24:00 的箭头表示
%    ③ 紧急购电发生在实时（结算时），不在决策时刻
%    ④ 四问在图上的差别：Q1 仅 0:00 一个决策点且无惩罚；Q2 有 0:00 + 实时惩罚；
%       Q3 有 4 个决策点；Q4 旁注"电价亦为随机"
% =============================================================================
\begin{tikzpicture}[
  x=0.46cm, y=1cm, font=\small, >=Stealth,
  ev/.style={font=\scriptsize, inner sep=1.5pt},
]

% ---------- 泳道底色 ----------
\fill[gray!8, rounded corners=2pt] (-3.15,-2.15) rectangle (26.2,-0.30);   % 信息
\fill[blue!5, rounded corners=2pt] (-3.15,-3.70) rectangle (26.2,-2.30);   % 决策
\fill[red!4,  rounded corners=2pt] (-3.15,-5.25) rectangle (26.2,-3.85);   % 结算

% ---------- 泳道标签 ----------
\node[rotate=90, font=\footnotesize\heiti, gray!60!black] at (-2.55,-1.22) {信息};
\node[rotate=90, font=\footnotesize\heiti, gray!60!black] at (-2.55,-3.00) {决策};
\node[rotate=90, font=\footnotesize\heiti, gray!60!black] at (-2.55,-4.55) {结算};

% ---------- 时间刻度（横轴 0:00 -> 24:00） ----------
\foreach \h in {0,6,12,18,24}{
  \draw[gray!25] (\h,-5.25) -- (\h,-0.30);
}
\foreach \h/\lab in {0/0:00, 6/6:00, 12/12:00, 18/18:00, 24/24:00}{
  \draw[gray!60] (\h,0.00) -- (\h,-0.16);
  \node[above, font=\scriptsize, gray!70!black] at (\h,0.00) {\lab};
}
\draw[gray!60, ->] (-0.9,0.00) -- (26.2,0.00) node[right, font=\scriptsize, gray!70!black] {时刻};
\draw[gray!45, dashed] (24,0.00) -- (24,-5.25);
\node[font=\scriptsize, gray!70!black, align=center] at (25.45,-5.55) {次日};

% ---------- 泳道 1：信息（附件 3 的四个预报发布时刻） ----------
% 预报覆盖区间的起点用实心圆标注
\foreach \h/\y/\len/\txt in {%
  0/-0.62/24/{\,覆盖至 24:00},
  6/-0.95/18/{\,覆盖至 24:00},
  12/-1.28/12/{\,覆盖至 24:00},
  18/-1.61/8/{~覆盖至次日 18:00（跨日）}}{
  \fill[blue!55!black] (\h,\y) circle (1.5pt);
  \draw[gray!60, line width=0.5pt, -{Stealth[length=4pt]}] (\h,\y) -- (\h+\len,\y);
  \node[ev, anchor=west, gray!25!black] at (\h+0.35,\y) {\h:00 发布\txt};
}

% ---------- 泳道 2：决策 ----------
% Q1：唯一决策时刻 0:00
\fill[blue!65!black] (0,-2.62) circle (2.4pt);
\node[font=\scriptsize, anchor=west, blue!50!black] at (0.35,-2.62)
  {计划购电：Q1\,/\,Q2\,/\,Q3\,/\,Q4（唯一决策时刻，随后不再修正）};
% Q3/Q4：6:00、12:00、18:00 调整
\foreach \h in {6,12,18}{
  \draw[blue!65!black, fill=white, line width=0.7pt] (\h,-3.28) circle (2.4pt);
}
\draw[gray!55, dotted, line width=0.6pt] (0,-2.62) -- (0,-3.28);
\node[font=\scriptsize, anchor=west, blue!50!black] at (6.35,-3.28)
  {调整购电：Q3\,/\,Q4 在 6:00、12:00、18:00 各修正一次};

% ---------- 泳道 3：结算 ----------
\draw[gray!55, densely dashed, line width=0.7pt, -{Stealth[length=5pt]}]
  (0,-4.20) -- (24,-4.20);
\node[font=\scriptsize, anchor=west, gray!25!black] at (0.35,-3.98)
  {结算依据：附件 2 实际负载与实际光伏功率（\textbf{与计划依据是两套数据}）};
% 紧急购电：发生在实时（结算时），不在决策时刻
\foreach \x in {9.4, 15.1, 21.3}{
  \draw[red!75!black, line width=0.9pt] (\x-0.28,-4.52) -- (\x+0.28,-4.52) -- (\x,-4.78) -- cycle;
}
\node[font=\scriptsize, anchor=west, red!65!black] at (-3.05,-4.98)
  {缺口 $\rightarrow$ 紧急购电（按 $5p_t$ 结算），发生在\textbf{实时}，不在任何决策时刻};

% ---------- 图例 ----------
\node[ev, anchor=west] at (-3.05,-5.62) {%
  \tikz\draw[gray!60,line width=0.5pt,-{Stealth[length=4pt]}] (0,0)--(0.55,0);
  ~预报（计划依据）\quad
  \tikz\draw[gray!55,densely dashed,line width=0.7pt,-{Stealth[length=4pt]}] (0,0)--(0.55,0);
  ~实际值（结算依据）\quad
  \tikz\fill[blue!65!black] (0,0) circle (2.2pt);~决策时刻\quad
  \tikz\draw[red!75!black,line width=0.9pt] (0,-0.05)--(0.28,-0.05)--(0.14,-0.26)--cycle;
  ~紧急购电（非决策时刻）\quad
  注：Q4 中电价亦为随机（旁注）};
\end{tikzpicture}

  \caption{微网决策时序与信息结构（上行泳道为信息到达时刻，中行为决策时刻，下行为结算依据；
            \(\Delta t=1/6\)~h，全天 144 个 10~min 时段；样本窗口 2025-02-01 至 12-31 共 334 天）}
  \label{fig:timeline}
\end{figure}

两侧惩罚方向相反，故不能用同一套边际规则处理。问题 3 中追加的边际代价是削减的边际收益的
\(1.5/0.5=3\) 倍，频繁改计划本身即付费；问题 2 与 4-2 则相反——计划量按正常电价全额结算，缺口按
\(5p_t\) 补足，故最优计划\textbf{主动多买}。消融实验（\S\ref{sec:ablation}）给出该结论的量级：
取消安全边际后 334 天总费用由 \(14\,264\,792.13\) 元升至 \(17\,398\,559.68\) 元（比无储能基线高
\(6.04\%\)），紧急购电量由 \(172\,015.82\)~kWh 升至 \(1\,156\,618.02\)~kWh，
边际价值 \(3\,133\,767.55\) 元。

\subsection{依赖关系与误差传播}

图~\ref{fig:dependency} 表明四问共享同一物理骨架（并网点电量平衡、储电量递推、充放电上下界），
逐层增加信息结构模块，使每层新增贡献可被\textbf{单独消融}（\chapref{chap:verify}）。

\begin{figure}[htbp]
  \centering
  %!TEX root = ../main.tex
% =============================================================================
%  图 2  四问的模型依赖关系（正文 §1.3）
%  规格：附件 -> Q1 -> Q2 -> Q3 -> Q4，每个箭头必须标注"提供什么"；
%        本图用于证明四问是一条信息结构递进链，而不是四篇互不相干的报告。
% =============================================================================
\begin{tikzpicture}[
  font=\small, >=Stealth,
  att/.style ={draw=gray!50, fill=gray!8, rounded corners=2pt, align=center,
               font=\scriptsize, text width=2.0cm, minimum height=0.72cm, inner sep=3pt},
  pre/.style ={draw=orange!55!black, fill=orange!8, rounded corners=2pt, align=center,
               font=\scriptsize, text width=6.6cm, minimum height=0.85cm, inner sep=4pt},
  qbox/.style={draw=blue!50, fill=blue!7, rounded corners=2pt, align=center,
               font=\small, text width=4.6cm, minimum height=0.95cm, inner sep=4pt},
  obox/.style={draw=green!45!black, fill=green!7, rounded corners=2pt, align=center,
               font=\scriptsize, text width=2.2cm, minimum height=0.8cm, inner sep=3pt},
  ar/.style  ={->, draw=gray!65, line width=0.6pt},
  lab/.style ={font=\scriptsize, anchor=west, align=left, text width=5.6cm,
               inner sep=1pt, gray!20!black},
]

% ---------------- 附件 ----------------
\node[att] (a1) at (-4.35,0) {附件 1\\电价/负载/光伏};
\node[att] (a2) at (-1.45,0) {附件 2\\负载/光伏实际值};
\node[att] (a3) at ( 1.45,0) {附件 3\\四个时刻光伏预报};
\node[att] (a4) at ( 4.35,0) {附件 4\\波动电价};

% ---------------- 数据预处理 ----------------
\node[pre] (pre) at (0,-1.35)
  {数据预处理：时标对齐（标签 $=$ 时段末端）\ \textbullet\ 附件 1 身份核查\ \textbullet\ 特征与误差分析（第三章）};

\foreach \n in {a1,a2,a3,a4}{\draw[ar] (\n.south) -- (pre.north);}

% ---------------- 四问 ----------------
\node[qbox] (q1) at (0,-2.95) {\textbf{问题 1}\ 确定性线性规划（LP）};
\node[qbox] (q2) at (0,-4.60) {\textbf{问题 2}\ 两阶段随机规划};
\node[qbox] (q3) at (0,-6.25) {\textbf{问题 3}\ 多阶段滚动 $+$ 非预期性};
\node[qbox] (q4) at (0,-7.90) {\textbf{问题 4}\ 波动电价 $+$ 风险度量};

\draw[ar] (pre.south) -- (q1.north);

% 箭头上的"提供什么"标注
\draw[ar] (q1.south) -- (q2.north);
\node[lab] at (0.35,-3.78)
  {\textbf{提供}\ 功率平衡与储电量递推骨架；无储能基线 $C^{base}$ 口径；
   储电水价值与阈值律语言};

\draw[ar] (q2.south) -- (q3.north);
\node[lab] at (0.35,-5.43)
  {\textbf{提供}\ 固定场景集（消融与 VoI 共用）；报童临界分位数给出的安全边际；
   紧急购电的结构性认识};

\draw[ar] (q3.south) -- (q4.north);
\node[lab] at (0.35,-7.08)
  {\textbf{提供}\ 多阶段滚动框架与非预期性实现；终端条件处理；
   预报频次对照（VoI）方法};

% ---------------- 输出 ----------------
\node[obox] (o1) at (-5.55,-2.95) {result1.xlsx\\+ 表 1 / 表 2};
\node[obox] (o2) at (-5.55,-4.60) {result2.xlsx\\+ 表 3};
\node[obox] (o3) at (-5.55,-6.25) {result3.xlsx\\+ VoI 结论};
\node[obox] (o4) at (-5.55,-7.90) {result4-2.xlsx\\result4-3.xlsx};

\foreach \q/\o in {q1/o1,q2/o2,q3/o3,q4/o4}{
  \draw[ar, draw=green!45!black] (\q.west) -- (\o.east);
}

\end{tikzpicture}

  \caption{四问的模型依赖关系与复用内容（箭头标注上游向下游提供的模型要素；四问在 2025-02-01 至
            12-31 共 334 天的同一窗口上报告结果）}
  \label{fig:dependency}
\end{figure}

四条上游链路的误差传播与处理见表~\ref{tab:errorprop}。其中最须防范的是第二节与末节指向的同一风险：
问题 1 的输入是\textbf{全年平均日}，问题 2--4 的窗口只有 334 天，两者是\textbf{两套口径}而非同一量的
两个估计，故分别命名、分别报告（表~\ref{tab:twocaliber}），节省率一律定义为“相对同一窗口、同一电价
口径下的无储能基线“。

\begin{table}[htbp]
  \centering
  \caption{四问之间的误差传播路径与本文的处理方式}
  \label{tab:errorprop}
  \footnotesize
  \setlength{\tabcolsep}{4pt}
  \begin{tabularx}{\linewidth}{L{2.15cm} L{3.05cm} Y}
    \toprule
    上游 \(\rightarrow\) 下游 & 传递的内容 & 潜在的误差传播与本文的处理 \\
    \midrule
    附件 \(\rightarrow\) 问题 1
      & 电价、负载、光伏的时标与单位
      & 时标错位一个槽或丢掉 \(0{:}00{+}1\) 一列会让全部合计系统性偏离，而“全天汇总 \(=\sum\) 144 格”
        的自检对整体错位不可见；本文统一按“标签 \(=\) 时段末端”的\textbf{位置读法}对齐，
        断言每天恰 144 槽并保留末列（\S3.2） \\
    问题 1 \(\rightarrow\) 问题 2
      & 模型骨架、\(C^{base}\) 口径、水价值语言
      & 两套口径量级不同（单日 \(48\,052.0466\) 元对 334 天 \(16\,407\,319.63\) 元），混排一次即失去
        可比性；本文各自命名并在表~\ref{tab:twocaliber} 显式声明 \\
    问题 2 \(\rightarrow\) 问题 3
      & 场景生成与削减方法、安全边际
      & 各问各自生成场景会把“场景噪声”误读成“信息结构改进”；本文把场景集\textbf{一次性固定并存盘}，
        问题 3 与全部消融实验读同一份（\S\ref{sec:ablation}），并以 \(K=60\to120\) 作稳健性检验 \\
    问题 3 \(\rightarrow\) 问题 4
      & 多阶段滚动框架、终端条件
      & 电价随机性与光伏误差叠加，分别抽样会丢失相关结构、低估“高价又缺电”的联合尾部；本文用
        “日净需求残差 \(\times\) 日均电价”二维分层取历史代表日生成\textbf{联合}情景，保留
        \(\rho(p,\text{净负荷})\)（实际 \(0.2620\)、情景 \(0.2597\)） \\
    问题 1--4 \(\rightarrow\) 结论
      & 各问的“总费用”与“节省率”
      & 基线不一致时同一“节省率”指向不同分母；四问统一使用 \calref{C-01} 的无储能基线与
        \calref{C-04} 的节省率，并逐处注明电价口径 \\
    \bottomrule
  \end{tabularx}
\end{table}

\begin{table}[htbp]
  \centering
  \caption{无储能基线的两套口径声明（费用单位：元；\calref{C-01} 定义见\chapref{chap:q1}）}
  \label{tab:twocaliber}
  \footnotesize
  \setlength{\tabcolsep}{4pt}
  \begin{tabularx}{\linewidth}{L{2.6cm} C{1.5cm} L{4.6cm} C{1.95cm} C{2.0cm}}
    \toprule
    口径名称 & 适用问题 & 负载与光伏来源 & 电价来源 & 基线数值 \\
    \midrule
    典型日口径 & 问题 1 & 附件 1 的 144 槽典型日（全年逐时段均值） & 附件 1（固定价）
      & \(48\,052.0466\) \\
    窗口口径（固定价） & 问题 2、3 & 附件 2 的当日实际负载与实际光伏，逐日计 & 附件 1（固定价）
      & \(16\,407\,319.63\) \\
    窗口口径（波动价） & 问题 4 & 同上一行 & 附件 4（波动价） & \(17\,026\,782.86\) \\
    \bottomrule
  \end{tabularx}
\end{table}

\subsection{与既有做法的差异}
各问单独选口径、或用元启发式直接求解，会使节省率的分母不可核对且拿不到对偶信息。本文的三点差异：
① 四问共用\textbf{同一物理骨架与同一基线口径}，逐层增加模块，每层贡献可单独消融；
② 全部模型为\textbf{线性规划}并用 HiGHS 精确求解，可输出对偶量 \(\nu_t\) 与阈值型策略的最优性证据；
③ 对题面“是否需要引入其他时刻的预报”给出\textbf{可判定阈值}：同一场景集与同一实际轨迹下，
预报时刻数 \(n=1,2,4\) 的总费用为 \(1\,426.48\)、\(1\,413.16\)、\(1\,393.24\) 万元，
\(n=4\) 已实现 \(85.1\%\) 的可获得信息价值，继续加密最多再省 \(58\,063\) 元（\(0.42\%\)），
故在本题预报精度与惩罚结构下\textbf{不建议继续加密}（\evid{C-Q3-01}）。

%!TEX root = ../main.tex
% =============================================================================
%  二、模型假设与符号说明        建议 2 页
%  评审在问：你的假设合理吗？符号清楚吗？
%  本节必须出现：① 假设表 A1–A10，每条含「必要性 + 现实/题目依据 + 影响方向 + 验证方式」
%                ② 信息结构 R1/R2/R3 的显式声明与实测对照
%                ③ 完整符号表，含单位，且必须含对偶变量
%                ④ 时标与单位口径声明（位置读法 / 电量口径 / Δt / E^max）
%  判据：一条假设若不能回答“不用它模型会怎样”，它就不该出现在假设表里。
% =============================================================================

\section{模型假设与符号说明}

\subsection{模型假设}

题面已给定储能最大容量 \(12000\)~kWh、最大充放电功率 \(5000\)~kW、2025 年 1 月 1 日
\(0{:}00\) 储电量 \(6000\)~kWh、储电量运行区间 \(1200\)--\(10800\)~kWh 与充放电效率
\(90\)\%；本文取运行区间为额定容量的 \(10\)\%--\(90\)\%（可用摆幅
\(10800-1200=9600\)~kWh），二者与功率上限构成两类硬约束，谁更紧由第四章容量/功率双向扫描
判定。

\begin{table}[htbp]
  \centering
  \caption{模型假设及其必要性、现实/题目依据、影响方向与验证方式（A1--A10）}
  \label{tab:assumptions}
  \footnotesize
  \setlength{\tabcolsep}{3pt}
  \begin{tabularx}{\linewidth}{C{0.65cm} L{2.45cm} L{2.95cm} L{2.6cm} L{1.55cm} Y}
    \toprule
    编号 & 假设 & 必要性（不用会怎样） & 现实/题目依据 & 影响方向 & 验证或灵敏度方式 \\
    \midrule
    A1 & 不允许向外部电网售电
       & 目标函数将多出一项售电收益，盈余有了出路，弃光量恒为 0，模型的可行域与最优值都改变
       & 题面未给出上网电价；就近消纳“优先自用”的政策导向
       & 使费用\textbf{偏高}，结果保守
       & 对结果文件作非负性审计（审计项 A7）：五个结果文件的负购电量数均为 0；
         并对引入上网电价 \(\theta p_t\) 的政策情景作扫描 \\
        A2 & 充电效率与放电效率\textbf{单边各取} \(0.9\)（往返 \(0.81\)）
       & 效率决定套利是否有利可图，是收益的核心参数；不声明读法则同一模型可给出三套结果
       & 题面给出“充放电效率为 \(90\%\)”，\textbf{未区分充、放方向}
       & 效率越低，往返损耗越大，储能套利收益越低
       & 三种主口径读法的节省额对照（单边各 \(0.9\)、仅充电侧 \(0.9\)、仅放电侧 \(0.9\)）
         见本节末段；另报“往返 \(0.9\) 等效单边”口径作条件对照 \\
        A3 & 同一时段不同时充放电
       & 若允许同时充放，须引入 \(0\)--\(1\) 变量，模型由线性规划变为混合整数规划，
         规模与求解时间上升，且对偶信息不再直接可用
       & 效率乘积 \(\eta_c\eta_d<1\) 时同时充放严格劣于只做净动作（命题~\ref{prop:no-simul}）
       & 不影响最优目标值，只影响求解方式
       & 数值断言 \(\min(c_t,d_t)=0\) 全部成立：问题 1 互斥违反数为 0，
         问题 2/3/4 的内部解互斥违反数也均为 0（审计项 A5） \\
        A4 & 忽略自放电与循环衰减
       & 否则目标函数须加入与循环次数（放电量）相关的衰减成本项，并需要衰减参数
       & 题面未给出任何衰减参数；结果窗口仅 334 天
       & 高估长期收益，且放电量越大高估越多
       & 加入循环成本 \(c_{cyc}\)（元/kWh 放电量）作单参数灵敏度，给出“储能节省额归零”的
         临界 \(c_{cyc}\) \\
        A5 & 储电量\textbf{跨日连续递推}，不逐日清零
       & 若每日清零，则与“2025-01-01 \(0{:}00\) 储电量 \(6000\)~kWh”的初值条件矛盾，
         并丢弃跨日套利空间
       & 附录 1 给出的初值 \(+\) 储能设备连续运行的常识；题面只要求 \(0{:}00\) 与
         \(24{:}00\) 储电量相同
       & 影响全窗口费用与每天的储电量初值，跨日套利越大影响越大
       & 与“每日重置初值”变体对照（消融层级 B2/B3）；审计项 A6 逐日核对
         “末值不低于初值”的违反天数（问题 2/3/4 均为 0） \\
        A6 & 每日终端条件 \(S_{144}\ge S_0\)；问题 1 取严格日周期 \(S_0=S_{144}=6000\)~kWh
       & 不设终端条件时，模型会在每天末尾与窗口末尾“掏空”储能以压低当期费用，
         形成短视解，且该解在真实连续运行中不可执行
       & 储能运营常识；问题 1 题面直接要求 \(S_0=S_{144}\)
       & 影响每天末尾与窗口末尾的费用；本数据上日周期口径比自由环状口径贵
         \(8.35\) 元（\(0.02\%\)）
       & 终端处理四组对照：完全自由、\(S_T\ge S_0\)、日周期、水价值终端（§9.4 层级 B5） \\
        A7 & 调整决策“后覆盖前”，调整费用的计价基准始终为 \(0{:}00\) 的\textbf{计划购电量}
       & 若基准随上一次调整漂移，调整费与违约费的口径不可复现，也无法与题面表述对应
       & 题面“计划购电量高于调整购电量的部分”一句明确了基准
       & 影响调整与违约相关费用；四次预报口径下违约费 \(210\,674.71\) 元、
         超额费 \(262\,628.79\) 元
       & 与“相对上一次调整量”口径并列报告并各自命名（第九章） \\
        A8 & 2025 年 1 月作为储电量递推的\textbf{预热期}，其终值即 2025-02-01 \(0{:}00\) 的初值
       & 结果窗口自 2025-02-01 起，而初值自 2025-01-01 起算；若不作预热声明，
         2 月 1 日的储电量没有来源，跨日递推无从开始
       & 附件 2 覆盖 2025-01-01 至 12-31，结果文件要求 2025-02-01 起共 334 天
       & 影响窗口初期的储电量与初期费用；1 月本身不计入任何结果
       & 论文与结果文件中明示 1 月不计入结果，并公开 2 月 1 日初值的来源；
         审计项 A6 逐日核对终端条件 \\
        A9 & \textbf{信息结构取 R3}：问题 2/4-2 的 \(0{:}00\) 计划依据为附件 3 中当日 \(0{:}00\)
         发布的光伏预报，负载曲线按“已知”处理；问题 3/4-3 另有 \(6{:}00/12{:}00/18{:}00\)
         三次滚动的当次预报
       & 题面只说“每天 \(0{:}00\) 制定计划”，未说明 \(0{:}00\) 知道什么；不声明则
         问题 2--4 的任何数值都无法解释（三种解读的窗口总费用相差近千万元）
       & 题面“根据预测数据制定计划”；附件 3 提供 \(0{:}00\) 发布的 24~h 光伏预报；
         题面未给负载预报，故按已知处理
       & 费用落在完美信息（下界）与典型日预报（上界）之间；负载按已知处理会
         \textbf{低估}紧急购电量，使结果偏乐观
       & 三种信息结构 R1/R2/R3 在同一 334 天窗口与同一场景集上的完整对照
         （表~\ref{tab:infostruct}）；\textbf{非预期性置换检验}：置换 \(6{:}00/12{:}00/18{:}00\)
         发布的预报后，其发布时刻之前各槽位的计划量逐位不变（最大差值 \(0.0\)），
         置换 \(0{:}00\) 预报则只改变 \(0{:}00\) 之后的决策（第九章）；
         并对预报误差水平另作灵敏度 \\
        A10 & 波动电价（问题 4）下 \(0{:}00\) 可用的\textbf{电价预报取“分时骨架”}，
         即电价的逐时段日内均值曲线；结算仍按附件 4 的实际电价
       & 附件 4 未给出日前电价预报；若不声明电价预报口径，问题 4 的“计划”将使用
         结算时的实际电价，构成数据泄漏，结果会被系统性压低
       & 分时骨架与附件 4 的逐时段均值最大偏差仅 \(5.15\times10^{-5}\)~元/kWh，
         且等于附件 1 的电价剖面，是可获得的日前信息
       & 忽略电价的残差波动，使问题 4 的费用与节省率\textbf{偏乐观}
       & 结算侧改用附件 4 实际电价，价格随机性只在结算体现；并以“是否保留残差自相关”
         作对照（分块 Bootstrap 条件对照的分位数最大相对偏差达 \(20.44\%\)） \\
    \bottomrule
  \end{tabularx}
\end{table}

假设 A2 影响最大：三种效率读法下，典型日相对无储能基线的节省额为
\(12\,925.10\)~元（单边各 \(0.9\)，\textbf{主口径}）、\(14\,635.52\)~元（仅充电侧 \(0.9\)）、
\(13\,883.44\)~元（仅放电侧 \(0.9\)），极差 \(13.2\%\)；把题面“\(90\%\)”读作往返效率、等效
为单边 \(\sqrt{0.9}\) 时为 \(14\,250.99\)~元。主口径写入全部结果表，其余读法只在灵敏度节出现
（\evid{C-Q1-03}）。

\subsection{信息结构声明}

问题 2--4 中“每天 \(0{:}00\) 制定计划”有三种解读，本文各跑一次全窗口实验，条件完全一致
（同一 334 天窗口、同一场景集、同一终端条件、同一电价口径），只有“\(0{:}00\) 已知什么”
不同，三行可直接相减（表~\ref{tab:infostruct}）。

\begin{table}[htbp]
  \centering
  \caption{信息结构的三种解读及其在 2025-02-01 至 12-31（334 天）上的实测后果
           （费用单位：元，电量单位：kWh，电价均为附件 1 固定价；
           紧急购电量口径见 \calref{C-05}，“出现在多少天”一并列出）}
  \label{tab:infostruct}
  \footnotesize
  \setlength{\tabcolsep}{3pt}
  \begin{tabularx}{\linewidth}{C{0.8cm} L{3.3cm} C{2.55cm} C{1.5cm} C{2.3cm} Y}
    \toprule
    代号 & 解读 & 总费用 & 相对基线节省率 & 紧急购电量 & 本文的使用方式 \\
    \midrule
    R1 & 完美信息：\(0{:}00\) 即知当日负载与光伏真值
       & \(13\,411\,560.90\) & \(18.26\%\) & \(0.00\)（0/334 天）
       & \textbf{不可作为主方案}；用作“预测误差的代价”的\textbf{下界对照} \\
        R2 & 典型日预报：以附件 1 的日均形状作为预报
       & \(23\,484\,646.63\) & \(-43.14\%\) & \(2\,525\,548.76\)（286/334 天）
       & 仅作悲观对照；附件 1 是实际值的全年逐时段均值而\textbf{不是预报}，
         用它当预报须额外论证并做泄漏检验，本文\textbf{不采用} \\
        R3 & 日前预报：\(0{:}00\) 使用附件 3 当日 \(0{:}00\) 发布的光伏预报，
         负载曲线按已知处理
       & \(14\,264\,792.13\) & \(13.06\%\) & \(172\,015.82\)（296/334 天）
       & \textbf{本文主方案}：最符合工程实际与题面“制定计划”的语义 \\
    \bottomrule
  \end{tabularx}
\end{table}

本文采用\textbf{信息结构 R3}：R1 假设 \(0{:}00\) 即知当天真值，紧急购电量恒为 0、问题 2
退化为问题 1；R2 用事后平均值冒充事前信息，总费用比无储能基线还高 \(43.14\%\)，
属口径错误；R3 用附件 3 中\textbf{当天 \(0{:}00\) 发布}的真实预报，全部信息先于决策时刻
可得，无泄漏也不夸大。三点补充：① 计划依据（附件 3 预报 \(+\) 负载曲线）与结算依据
（附件 2 当日实际负载与光伏）是\textbf{两套数据}，差额即紧急购电来源；② “预测误差的代价”
为 \(14\,264\,792.13-13\,411\,560.90=853\,231.23\) 元、占 R1 费用的 \(6.36\%\)
（\evid{C-Q2-01}）；③ 负载按“已知”处理是对 R3 的乐观简化（题面未给负载预报），会低估
紧急购电量，第九章另作预报误差灵敏度。

\subsection{符号说明}

公式符号均在表~\ref{tab:notation} 登记，单位口径见本节末段。

\begin{table}[htbp]
  \centering
  \caption{符号说明（集合、参数、决策变量、对偶变量与指标量；“---”表示无量纲）}
  \label{tab:notation}
  \footnotesize
  \setlength{\tabcolsep}{3pt}
  \begin{tabularx}{\linewidth}{C{1.2cm} C{2.55cm} L{5.1cm} C{1.35cm} Y}
    \toprule
    类别 & 符号 & 含义 & 单位 & 来源/取值 \\
    \midrule
    集合 & \(t\in\mathcal T\) & 时段（10~min 粒度） & --- & \(|\mathcal T|=144\) \\
    集合 & \(d\in\mathcal D\) & 日期 & --- & 334 天（2025-02-01 至 12-31） \\
    集合 & \(m\in\mathcal M\) & 预报发布/决策时刻 & --- & \(\{0,6,12,18\}\) 时 \\
    集合 & \(\omega\in\Omega\) & 情景 & --- & \(|\Omega|=K\) \\
        参数 & \(L_{d,t}\) & 小区负载（功率） & kW & 附件 2 \\
    参数 & \(PV_{d,t}\) & 光伏实际功率 & kW & 附件 2 \\
    参数 & \(\widehat{PV}^{(m)}_{d,h}\) & \(m\) 时发布的光伏预报功率（整点值） & kW & 附件 3 \\
    参数 & \(p_{d,t}\) & 电价 & 元/kWh & 附件 1（固定价）/ 附件 4（波动价） \\
    参数 & \(\bar p_t\) & 电价的日内分时骨架（逐时段均值） & 元/kWh & 附件 1 电价剖面（假设 A10） \\
    参数 & \(\varepsilon_{d,t}\) & 电价相对骨架的残差 & 元/kWh & \(p_{d,t}-\bar p_t\) \\
    参数 & \(\Delta t\) & 时间步长 & h & \(1/6\) \\
    参数 & \(\eta_c,\eta_d\) & 充电效率、放电效率（单边） & --- & \(0.9/0.9\)（假设 A2） \\
    参数 & \(\underline S,\overline S\) & 储电量运行下限、上限 & kWh & \(1200/10800\) \\
    参数 & \(\overline S^{\max}\) & 储能设备最大容量 & kWh & \(12000\) \\
    参数 & \(\overline P\) & 最大充放电功率 & kW & \(5000\) \\
    参数 & \(E^{\max}\) & 单槽充/放电量上限 & kWh & \(\overline P\Delta t=833.3333\) \\
    参数 & \(S^{init}\) & 初始储电量 & kWh & \(6000\) @ 2025-01-01 \(0{:}00\) \\
    参数 & \(\kappa\) & 紧急购电电价倍数 & --- & \(5\) \\
    参数 & \(\kappa^{+},\kappa^{-}\) & 调整超额 / 违约电价倍数 & --- & \(1.5/0.5\) \\
    参数 & \(\beta\) & CVaR 置信水平（预先锁定） & --- & \(0.95\) \\
    参数 & \(\pi_\omega\) & 情景权重 & --- & 分层样本占比，\(\sum_\omega\pi_\omega=1\) \\
    参数 & \(\rho(\cdot,\cdot)\) & 相关系数 & --- & 用于刻画电价与净负荷的共动 \\
        决策 & \(E^{plan}_{d,t}\) & 计划购电量（第一阶段，与 \(\omega\) 无关） & kWh & \(\ge0\) \\
    决策 & \(E^{adj}_{d,t}\) & 调整购电量（最后覆盖该槽位的决策） & kWh & \(\ge0\) \\
    决策 & \(g_{d,t}\) & 紧急购电量（第二阶段追索） & kWh & \(\ge0\) \\
    决策 & \(c_t,d_t\) & 充电量、放电量 & kWh & \(0\le c_t,d_t\le E^{\max}\) \\
    决策 & \(s_t\) & 弃光/盈余处置电量 & kWh & \(\ge0\) \\
    决策 & \(S_t\) & 时段末储电量 & kWh & \(\in[\underline S,\overline S]\) \\
    决策 & \(u_t,v_t\) & 调整量相对计划量的增、减分解 & kWh & \(E^{adj}_t-E^{plan}_t=u_t-v_t\)，均 \(\ge0\) \\
        \textbf{对偶} & \(\alpha_t\) & 电量平衡约束的影子价格（\(E_t>0\) 时等于 \(p_t\)） & 元/kWh & LP 对偶；实测取值 \([0.3713,1.2561]\)，恒等式偏差 \(0.000\) \\
    \textbf{对偶} & \(\nu_t\) & 储电量递推约束的乘子，即储电的边际价值（水价值） & 元/kWh & LP 对偶；实测 \(\nu_t\in[0.4711,1.1305]\) \\
    \textbf{对偶} & \(\mu_t\) & 水价值的等价记法（\(\mu_t\equiv\nu_t\)），第四章阈值律沿用 & 元/kWh & LP 对偶 \\
    \textbf{对偶} & \(\lambda^{\pm}_t\) & 储电量上/下界的影子价格 & 元/kWh & LP 对偶 \\
    \textbf{对偶} & \(\nu^{c}_t,\nu^{d}_t\) & 充/放电功率界的乘子 & 元/kWh & LP 对偶 \\
    \textbf{对偶} & \(z^{c}_t,z^{d}_t\) & 充/放互补松弛量 & 元/kWh & \(z^{c}_t=\alpha_t-\eta_c\nu_t\)；\(z^{d}_t=-\alpha_t+\nu_t/\eta_d\)（阈值律：\(\alpha_t\le\eta_c\nu_t\) 充电、\(\alpha_t\ge\nu_t/\eta_d\) 放电，144 槽违反次数 0） \\
        指标 & \(C^{base}\) & 无储能基线购电费 & 元 & 口径 \calref{C-01} \\
    指标 & \(C^{plan},C^{emg}\) & 计划购电费、紧急购电费 & 元 & 口径 \calref{C-03} \\
    指标 & \(V^{arb}\) & 储能与弃光带来的相对节省 & 元 & \(C^{base}-C^{*}\)，节省率按 \calref{C-04} \\
    \bottomrule
  \end{tabularx}
\end{table}

\paragraph{时标与单位口径声明}
本节符号遵循同一套时标与单位口径（四问一致）：① \textbf{位置读法}——附件 1/2/4 的标签为
\textbf{时段末端}，槽位 \(j\)（\(j=0,1,\dots,143\)）表示区间 \([10j,\ 10(j+1))\) 分钟，
故题面表 1 的“\(10{:}00\)--\(10{:}10\)”对应槽位 60；末列 \(0{:}00{+}1\) 属当日第 144 个槽位
\([23{:}50,24{:}00)\)、不得丢弃（逐格对照见 §3.2 表~\ref{tab:timeruler}）。② \textbf{电量
口径}——决策变量为电量（kWh），\(\Delta t=1/6\)~h，功率上限 \(\overline P=5000\)~kW 对应
单槽电量上限 \(E^{\max}=\overline P\Delta t=833.3333\)~kWh；目标函数的电价系数为
\(p_t\)（元/kWh）、\textbf{不再乘} \(\Delta t\)；等价的功率口径为“变量取功率、目标系数取
\(p_t\Delta t\)、上限取 \(5000\)~kW”，两套口径全文只用一套。③ \textbf{单位与精度}——功率
kW、电量 kWh、电价 元/kWh、费用 元；表格费用保留 2 位小数，内部计算不取整。

%!TEX root = ../main.tex
% =============================================================================
%  三、数据预处理与特征分析      建议 2–3 页
%  评审在问：数据你读懂了吗？
%  本节必须出现：① 时标与对齐结论  ② 附件 1 身份核查（= 全年逐时段均值，非预报）
%                ③ 相关性的尺度拆解  ④ 预报误差分档（全时段 / 白天双口径）
%                ⑤ 数据质量与缺失值结论
%  禁止：相关性写成因果（“导致”“引起”“造成”一律换成“共动”“同向变动”）
% =============================================================================

\section{数据预处理与特征分析}

\subsection{数据来源与基本规模}

四个附件的用途见表~\ref{tab:data}。时间序列均按 \textbf{10~min} 粒度、以\textbf{时段末端}
为标签对齐，全天 144 个时段；负载与光伏为\textbf{功率}（kW）、电价为\textbf{元/kWh}，
不得混用。



\subsection{时标语义与对齐}

附件 1、2、4 的时段标签为 \(0{:}10,0{:}20,\dots,23{:}50,0{:}00{+}1\)（后者为次日
\(0{:}00\)），本文采用\textbf{位置读法}：

\begin{itemize}
  \item 标签表示\textbf{时段末端}：槽位 \(j\)（\(j=0,1,\dots,143\)）表示区间
        \([10j,\ 10(j+1))\) 分钟；第 \(j+1\) 个数据行/列对应槽位 \(j\)，标签只作显示、
        \textbf{不参与重映射}；
  \item 每行 144 个值\textbf{无缝铺满} \(0{:}00\) 至 \(24{:}00\)；末列 \(0{:}00{+}1\) 属当日
        \(23{:}50\)--\(24{:}00\)、\textbf{不得丢弃}，否则合计偏低 \(1/144\approx 0.6944\%\)。
\end{itemize}

附件 3 的列标签为“预报 1 小时”至“预报 24 小时”，约定\textbf{预报 \(k\) 小时对应发布
时刻之后的第 \(k\) 个小时}，即 \([h_0+k-1,\ h_0+k)\)（\(h_0\in\{0,6,12,18\}\)）；若误按
\([h_0+k,\ h_0+k+1)\) 对齐，则\textbf{全部预报整体错 1 小时}。\(18{:}00\) 的预报中
\(k=7,\dots,24\) 落在\textbf{次日} \(0{:}00\)--\(17{:}00\)（18 个整点），跨日用绝对时间轴
映射（目标整点序号 \(d\times 24+h_0+k-1\)），不丢配对样本。表~\ref{tab:timeruler} 给出题面
时段与附件标签的对照。

\begin{table}[htbp]
  \centering
  \caption{题面表 1 的时段写法与附件时段标签对照（示例；时段序号 \(i\) 从 1 起计，
           槽位下标 \(j=i-1\)）}
  \label{tab:timeruler}
  \small
  \begin{tabular}{cccc}
    \toprule
    题面写法 & 附件标签 & 对应区间 & 数据点序号 \(i\) \\
    \midrule
    \(10{:}00\)--\(10{:}10\) & \(10{:}10\) & \([10{:}00,10{:}10)\) & 61 \\
    \(12{:}00\)--\(12{:}10\) & \(12{:}10\) & \([12{:}00,12{:}10)\) & 73 \\
    \(14{:}00\)--\(14{:}10\) & \(14{:}10\) & \([14{:}00,14{:}10)\) & 85 \\
    \(16{:}00\)--\(16{:}10\) & \(16{:}10\) & \([16{:}00,16{:}10)\) & 97 \\
    \(20{:}00\)--\(20{:}10\) & \(20{:}10\) & \([20{:}00,20{:}10)\) & 121 \\
    \(23{:}50\)--\(24{:}00\) & \(0{:}00{+}1\) & \([23{:}50,24{:}00)\) & 144 \\
    \bottomrule
  \end{tabular}
\end{table}

对齐结果由两条独立校验确认：结构断言（附件 1 恰 144 行；附件 3 前向填充后每天恰 4 行）与结果
文件层面的汇总校验——审计项 A10 核对“全天购电量/全天购电费”与逐时段求和，实测最大相对偏差为
\(3.40\times10^{-7}\)（问题 2、3）与 \(4.32\times10^{-7}\)（问题 4-2），远低于
\(10^{-4}\) 的阈值。该汇总校验对“整体错位一个槽”不可见，该风险由位置读法与
表~\ref{tab:timeruler} 逐格对照锁定。

\subsection{附件 1 的身份核查}

附件 1 第 4 列的表头写作“光伏发电\textbf{预测}功率”，但逐项核查表明它并非预报：

其\textbf{光伏}、\textbf{负载}、\textbf{电价}列分别等于附件 2 全年（365 天）逐时段实际光伏功率、
负载与附件 4 逐时段电价的均值，最大偏差分别为 \(0.0431\)~kW、\(5.04\times10^{-5}\)~kW 与
\(5.15\times10^{-5}\)~元/kWh。

三处偏差都在数值精度量级（\(10^{-5}\) 及以下），故附件 1 是\textbf{全年平均日（典型日）}而非
日前预报（电价列的一致同时是假设 A10“分时骨架”的数据依据）：附件 1 \textbf{只能}用于问题 1
与基线口径，\textbf{不可}用作问题 2--4 的预报输入或验证集；问题 1 因此是
\textbf{确定性优化问题}，其最优值是问题 2--4 的对照基准。

换用窗口口径（2025-02-01 至 12-31 共 334 天）的逐时段均值作比较，电价偏差升到
\(3.07\times10^{-2}\)~元/kWh，为\textbf{全年 365 天}口径的
\(3.07\times10^{-2}/5.15\times10^{-5}\approx600\) 倍：\textbf{全年口径}支撑“附件 1 即
附件 2/4 逐时段均值”的身份断言，\textbf{窗口口径}则说明问题 1 的典型日结果\textbf{不可}
直接当作问题 2--4 的窗口基准（§1.3 表~\ref{tab:errorprop} 第二行）。

\subsection{电价与负载的结构特征}

按附件 4 统计，波动电价均值 \(0.7662\)~元/kWh、最小 \(0.0076\)、最大 \(1.7936\)~元/kWh；
按附件 1 统计，固定电价均值同为 \(0.7662\)~元/kWh、最小 \(0.3713\)、最大 \(1.3952\)~元/kWh。
按附件 2 统计，负载均值 \(4626.0337\)~kW、最小 \(1995.7176\)、最大 \(7978.8849\)~kW，日均
用电量 \(110\,802.6332\)~kWh；光伏均值 \(2311.7857\)~kW、峰值 \(10\,216.2000\)~kW，日均
发电量 \(57\,092.4850\)~kWh。

两类结构性事实直接影响建模方式。其一，\textbf{套利是否有利可图}：盈亏平衡价比阈值为
\(1/(\eta_c\eta_d)=1/0.81\approx1.2346\)，而固定电价的极值之比
\(1.3952/0.3713\approx3.76\)（口径见 §1.1）远高于它。其二，\textbf{约束落在哪一侧}：光伏峰值
\(10\,216.2000\)~kW 明显高于充放电功率上限 \(5000\)~kW，故有光伏功率超出充电能力的时段，
盈余只能靠储能吸纳或弃光；容量摆幅与功率谁更紧由第四章的容量/功率双向扫描判定。

电价与负载的关系具有\textbf{尺度依赖性}（表~\ref{tab:corr} 按三个尺度报告）。

\begin{table}[htbp]
  \centering
  \caption{电价与负载、净负荷、光伏的相关系数（三个尺度；池化与去形状口径的样本量为
           \(365\times144=52\,560\) 个时段观测，日级口径为 365 个日观测；
           全样本为 2025 年全年 365 天口径，
           与情景保真度检验（\(0.2620\rightarrow0.2597\)）验证，见 §7.2）}
  \label{tab:corr}
  \small
  \setlength{\tabcolsep}{4pt}
  \begin{tabular}{lccc}
    \toprule
    变量对 & 池化（全部时段） & 日级（逐日先算再合并） & 剔除日内形状后 \\
    \midrule
    电价 --- 负载 & \(+0.5077\) & \(+0.8977\) & \(+0.4554\) \\
    电价 --- 净负荷 \(L-PV\) & \(+0.2602\) & \(+0.9824\) & \(+0.1583\) \\
    电价 --- 光伏 & \(+0.0284\) & \(-0.3431\) & \(+0.0484\) \\
    \bottomrule
  \end{tabular}
\end{table}

读法如下。① \textbf{电价与负载共动}：日级 \(+0.8977\)，剔除日内共同形状后仍保留
\(+0.4554\)。② \textbf{“电价—光伏”的池化值具有误导性}：池化仅 \(+0.0284\)、日级
\(-0.3431\)，方向相反；真实结构是\textbf{光伏出力大的午间对应电价低谷}，独立抽样会生成
“午间光伏大发同时电价高企”的不可能情景并低估尾部费用，故问题 4
用保留相关结构的联合情景（§7.2）。③ “电价—净负荷”随尺度递进：池化 \(+0.2602\)、日级
\(+0.9824\)、剔除日内形状后 \(+0.1583\)；④ “电价—光伏”去形状后为 \(+0.0484\)，接近 0，说明
同向关系几乎全部来自日尺度。以上只描述\textbf{共动/同向变动}，本文不主张其中任何一方为因。

\subsection{光伏预报误差的统计刻画}

本节把附件 3 中\textbf{当天 \(0{:}00\) 发布}的 24~h 预报（即问题 2/4-2 的日前计划依据）与
附件 2 实际功率逐槽配对：整点预报先展开为 6 个 10~min 槽，再逐槽相减得
\(r=\widehat{PV}-PV\)（单位 kW）。夜间光伏接近 0、相对误差无意义，故同时报告两个口径，并统一
以\textbf{目标整点} \(h\ge5\)（\(5{:}00\)--\(24{:}00\)，每天 19 个整点）为白天口径、不按功率
阈值筛选，使两口径分母固定、可比（表~\ref{tab:fcsterr}）。

\begin{table}[htbp]
  \centering
  \caption{光伏预报误差统计（样本为 2025 年每天 \(0{:}00\) 发布的 24~h 预报与附件 2 对应
           实际功率的逐槽配对；MAE、RMSE、bias 单位均为 kW）}
  \label{tab:fcsterr}
  \small
  \begin{tabular}{lcccc}
    \toprule
    口径 & 样本量 \(n\) & MAE & RMSE & bias \\
    \midrule
    全时段（目标整点 \(0\)--\(23\) 时，144 槽/天） & \(52\,560\) & \(361.58\) & \(636.37\) & \(-26.96\) \\
    白天（目标整点 \(\ge5\) 时，114 槽/天） & \(41\,610\) & \(456.72\) & \(715.22\) & \(-34.06\) \\
    \bottomrule
  \end{tabular}
\end{table}

白天口径的 MAE 为全时段口径的 \(456.72/361.58\approx1.26\) 倍（夜间接近 0 的样本摊薄了
全时段指标），建模应以白天口径为依据。更关键的是\emph{bias 的符号}：按目标整点聚合后正负号
系统性变号，见表~\ref{tab:biasprofile}。

表~\ref{tab:biasprofile} 的形状有直接的决策含义：\textbf{上午被系统性高估}（\(7\) 时达
\(+733.00\)~kW）使 \(0{:}00\) 计划的光伏出力偏大、购电量偏少，缺口更易触发 \(5\) 倍紧急购电
（口径 \calref{C-05}）；\textbf{下午被系统性低估}（\(15\) 时达 \(-801.87\)~kW）使计划购电偏多，
多买电量按计划电价白付，或在问题 3 中触发 \(0.5\) 倍违约电价与 \(1.5\) 倍超额电价。全天聚合
bias 只有 \(-26.96\)~kW（白天 \(-34.06\)~kW），是两类相反偏差抵消后的净额，\textbf{不能据此
判断风险方向}；\(19\)--\(23\) 时的平均偏差落在 \([-0.15,0.00]\)~kW 内，出力均接近 0，
不触发任何惩罚。故决策层必须按\textbf{目标整点}设置安全边际而非按全天平均误差一次性纠偏，
场景生成直接用逐槽预报残差的经验分布（§5.2）以保留这种双向结构（\evid{C-Q2-01}）。

\paragraph{数据质量与缺失值}
三条结构断言全部通过：附件 1 恰 144 行；附件 3 日期列前向填充后每天恰有 4 个发布时刻
（合计 1460 行）；结果窗口恰 334 天（2025-02-01 至 12-31），结果表恰 147 列（日期 \(+\) 144
个时段 \(+\) 全天购电量 \(+\) 全天购电费）。断言通过后不再插补：附件 3 稀疏日期列以
\texttt{ffill} 补齐，其余附件按位置读法逐列读入。数据侧唯一实质性质量问题是\textbf{预报质量的
非均匀性}（上午与下午偏差方向相反、傍晚之后无有效预报信息、误差水平白天高于全时段），故本文
不采用“预报误差独立同分布且无偏”的简化，而把逐槽残差的经验分布直接用于场景生成
（§5.2、§7.2），并在决策层按目标整点设安全边际与滚动更新（§5.4、§6.3）。

%!TEX root = ../main.tex
% =============================================================================
%  四、问题 1：确定性条件下的计划购电模型      目标 4 页
%  评审在问：确定性版本做对了吗？
%  本节含：① 电量口径与 LP 全式（目标／平衡含弃光松弛／SOC 递推／SOC 界／功率界／日周期）
%          ② 性质引理（充放互斥、能量恒等式）与证明
%          ③ 对偶分析：水价值 ν_t、并网点边际价值 α_t、阈值律与死区宽度
%          ④ 无储能基线 C-01、成本分解恒等式、容量/功率扫描、效率口径极差
%  图名：fig03_q1_dual_price.pdf / fig10_strategy_mechanism.pdf
%  单位纪律：功率 kW、电量 kWh、电价 元/kWh、费用 元；Δt=1/6 h；P̄Δt=833.3333 kWh
% =============================================================================

\section{问题 1：确定性条件下的计划购电模型}\label{chap:q1}

\subsection{模型}

\textbf{口径。} 采用电量口径：决策变量均以电量（kWh）为单位，电价 \(p_t\)（元/kWh）与电量相乘即
费用，\textbf{不再乘 \(\Delta t\)}；负载与光伏（kW）进模型前乘以 \(\Delta t\) 化为电量。
电价与负载逐日相同、光伏给出一日预测，故只需求解一个典型日（附件 1），信息结构为“当日真值已知”。
储能参数：额定容量 \(12\,000\)~kWh，运行区间 \([\underline S,\overline S]=[1200,10\,800]\)~kWh
（可用摆幅 9600~kWh），额定功率 \(\overline P=5000\)~kW，单边效率 \(\eta_c=\eta_d=0.9\)（假设 A2）。

决策变量为 \((E_t,c_t,d_t,s_t,S_t)\)，\(t\in\mathcal T\)，\(|\mathcal T|=144\)。模型为

\begin{subequations}\label{eq:q1-model}
\begin{align}
  \min_{E,c,d,s,S}\quad & C^{\text{plan}}=\sum_{t}p_tE_t \label{eq:q1-obj}\\
  \text{s.t.}\quad
  & E_t+(PV_t-L_t)\Delta t+d_t-c_t-s_t=0, && \forall t \label{eq:q1-balance}\\
  & S_t=S_{t-1}+\eta_c c_t-d_t/\eta_d, && \forall t \label{eq:q1-soc}\\
  & \underline S\le S_t\le\overline S,\quad 0\le c_t,d_t\le \overline P\Delta t, && \forall t \label{eq:q1-socbound}\\
  & E_t\ge0,\quad s_t\ge0, && \forall t \label{eq:q1-nonneg}\\
  & S_0=S_{144}=S^{init}=6000\ \text{kWh}. \label{eq:q1-periodic}
\end{align}
\end{subequations}

式~\eqref{eq:q1-balance} 即“微网供给不低于小区负载”的代数表达，\((PV_t-L_t)\Delta t\) 为净负荷电量
（盈余为负）。光伏盈余可能超过充电功率与负载之和，故以\textbf{等式加松弛 \(s_t\ge0\)} 写出以保证
可行；典型日的日内盈余低于摆幅 9600~kWh 与单槽吸纳能力 \(\overline P\Delta t=833.3333\)~kWh，
故最优 \(s_t\equiv0\)。同一典型日上“含松弛等式／不含松弛等式／不等式”三种写法的
最优值完全相同（均为 \(35\,126.9486\)~元），保留 \(s_t\) 的意义在于实际日期序列
（窗口内光伏峰值 \(10\,216.2000\)~kW）的可行性。式~\eqref{eq:q1-periodic} 即题面“两端储电量相同”，
本文进一步取两者均等于附件给定的初值 6000~kWh。

模型规模为 \(4\times144=576\) 个变量、288 条约束（显式保留 \(s_t\) 时为 720 个变量），是连续 LP；
充放互斥由模型结构自动成立，\textbf{无需 0--1 变量}。

\subsection{两条结构性质}

\begin{proposition}[充放互斥]\label{prop:no-simul}
若 \(\eta_c\eta_d<1\)，则最优解中不存在 \(c_t,d_t\) 同时为正的时段。
\end{proposition}

\begin{proof}
设某 \(t\) 有 \(c_t>0,d_t>0\)。取 \(\delta=\min\{c_t,\ d_t/(\eta_c\eta_d)\}>0\)，作
\(c_t'=c_t-\delta\)、\(d_t'=d_t-\eta_c\eta_d\delta\ge0\)，其余不变。由式~\eqref{eq:q1-soc}，
储电量变化 \(\eta_c(-\delta)+\eta_c\eta_d\delta/\eta_d=0\)，故储电量轨迹与全部界约束不变，新解可行；
由式~\eqref{eq:q1-balance}，购电量变化
\(\Delta E_t=-\delta+\eta_c\eta_d\delta=-\delta(1-\eta_c\eta_d)<0\)，目标下降
\(p_t\delta(1-\eta_c\eta_d)>0\)，与原解最优矛盾。
\end{proof}

命题~\ref{prop:no-simul} 只用 \(\eta_c\eta_d<1\)、\(p_t>0\) 与平衡式结构，对问题 1--4 的全部变体
一致成立。数值验证：同时充放时段数 0，储电量递推残差 \(9.09\times10^{-13}\)~kWh，
充放比偏差 \(4.44\times10^{-16}\)（审计 A5）。

\begin{lemma}[能量恒等式]\label{lem:energy}
逐槽由式~\eqref{eq:q1-balance} 解出 \(E_t\) 并对 \(t\) 求和，再用式~\eqref{eq:q1-soc} 消去充放电量，得
\begin{equation}
  \sum_{t}E_t=\sum_{t}(L_t-PV_t)\Delta t+\sum_{t}c_t-\sum_{t}d_t
  \ \xrightarrow{\ S_{144}=S_0\ }\
  \sum_{t}(L_t-PV_t)\Delta t+\bigl(1-\eta_c\eta_d\bigr)\sum_{t}c_t .
  \label{eq:q1-energy}
\end{equation}
\end{lemma}

引理~\ref{lem:energy} 有三点用途。① \textbf{它是“恒价无套利”的代数根源}：若日内电价恒为
\(\bar p\)，则目标 \(\bar p\sum_tE_t\) 只随循环量 \(\sum_tc_t\) 变化，与充放时序无关，储能\emph{没有}
时段套利价值；反之，题面背景段所指的“低充高放”只有在日内电价随时刻变化时才有对象，这从代数上支持
假设 A11（附件 1 的 144 点曲线逐日重复）。② \textbf{它是零成本的可行性回归校验}：逐槽核验
式~\eqref{eq:q1-energy}，典型日两侧之差在 \(10^{-4}\)~kWh 量级（即四位小数取整量级）。
③ \textbf{它分离了两问的价值来源}：问题 1 的节省来自改变 \(\sum_tE_t\) 的\emph{时段结构}（套利），
问题 2 的安全边际来自改变其\emph{分布}（降低高价缺口的期望），两者可叠加但不可互替。

\subsection{对偶分析与阈值型最优策略}

记 \(\nu_t\) 为式~\eqref{eq:q1-soc} 的影子价格，即\textbf{单位储电量的边际价值}（水价值，元/kWh）；
\(\alpha_t\) 为式~\eqref{eq:q1-balance} 的影子价格（并网点电量的边际价值，元/kWh，取非负约定）。
在储电量界与功率界均不活跃时，对 \(c_t,d_t\) 的平稳性条件给出互补松弛对

\begin{equation}
  z^{c}_t=\alpha_t-\eta_c\nu_t\ge0,\quad c_tz^{c}_t=0;\qquad
  z^{d}_t=-\alpha_t+\nu_t/\eta_d\ge0,\quad d_tz^{d}_t=0 .
  \label{eq:q1-kkt}
\end{equation}

其经济含义为：多充 1~kWh 的边际代价是 \(\alpha_t\)，换回 \(\eta_c\)~kWh 储电、价值
\(\eta_c\nu_t\)，故仅当 \(\alpha_t\le\eta_c\nu_t\) 时充电；多放 1~kWh 可替代 \(\alpha_t\) 的购电成本，
代价是储电量减少 \(1/\eta_d\)~kWh，故仅当 \(\alpha_t\ge\nu_t/\eta_d\) 时放电。当 \(E_t>0\) 时由互补
松弛得 \(\alpha_t=p_t\)（实测偏差 \(\le10^{-6}\)）；当 \(E_t=0\) 时 \(\alpha_t\le p_t\)。由此得

\begin{proposition}[阈值律]\label{prop:threshold}
在储电量界不活跃（\(\underline S<S_t<\overline S\)）的时段，最优充放策略由两条阈值刻画：
\(\alpha_t\le\eta_c\nu_t\) 时充电，\(\alpha_t\ge\nu_t/\eta_d\) 时放电，而
\begin{equation}
  \eta_c\nu_t\ <\ \alpha_t\ <\ \nu_t/\eta_d
  \label{eq:q1-deadband}
\end{equation}
时 \(c_t=d_t=0\)，构成“死区”。
\end{proposition}

死区宽度随往返效率下降而变宽：两条阈值之比为 \((\nu_t/\eta_d)/(\eta_c\nu_t)=1/(\eta_c\eta_d)=1.2346\)。
即\textbf{放电时段的并网点电量边际价值须高于充电时段的 1.2346 倍，跨时段套利才有利可图}；在购电
时段该判据化为放电电价高于充电电价 1.2346 倍，此即式~\eqref{eq:ch1-arb} 的盈亏平衡比。

\textbf{适用边界。} 命题~\ref{prop:threshold} 在“纯套利子问题（变量仅 \(c_t,d_t,S_t\)）且储电量界
不活跃“时精确成立；界活跃时须回到式~\eqref{eq:q1-kkt} 并计入界乘子；功率界活跃时 \(z^c_t\) 或
\(z^d_t\) 不再为零，简单阈值关系失效。数值验证分两层：144 个时段上按 \(\alpha_t\) 判定充／放／死区的
\textbf{违反次数为 0}；界不活跃的 97 个时段子集（充电槽 42、放电槽 36）上违反次数同样为 0，
两条互补松弛恒等式最大残差为 0。实测 \(\nu_t\in[0.4711,1.1305]\)、\(\alpha_t\in[0.3713,1.2561]\)~元/kWh
均\textbf{时变}，故不存在“全天统一的电价阈值”；正确表述是两级结构——先由 LP 对偶确定时变水价值曲线，
再逐时段比较 \(\alpha_t\) 与两条阈值 \evid{C-Q1-02}。

\begin{figure}[htbp]
  \centering
  \includefig[width=0.60\linewidth]{fig03_q1_dual_price.pdf}
  \caption{最优策略的机制解释（三子图共享时间轴：光伏与负载功率、电价与储电量（右轴标出上下界
            1200/10800~kWh）、充放电量；对象：附件 1 典型日 144 槽）}
  \label{fig:dualprice}
\end{figure}

\subsection{无储能基线与成本分解}

无储能时 \(c_t=d_t=s_t=0\)，由式~\eqref{eq:q1-balance} 得闭式
\(E_t=\Delta t\max(L_t-PV_t,0)\)，据此定义基线

\begin{equation}
  C^{base}=\sum_t p_t\Delta t\max\bigl(L_t-PV_t,0\bigr),\qquad
  \sum_t p_tE_t=C^{base}+C^{batt},
  \label{eq:q1-decomp}
\end{equation}

其中 \(C^{batt}=\sum_t p_t\bigl(\Delta t\min(L_t-PV_t,0)+c_t-d_t+s_t\bigr)\) 为储能与弃光带来的成本
增量（典型日取负值即费用下降）。式~\eqref{eq:q1-decomp} 是\textbf{逐项恒等式}而非两个独立最优值
之和，\(C^{batt}\) 的最小化仍需求解式~\eqref{eq:q1-model}。基线有三个变体，\textbf{不得混用}
（表~\ref{tab:cbases}），否则全文多个“节省率”互不可比。

\begin{table}[htbp]
  \centering
  \caption{无储能基线 \(C^{base}\) 的三个变体及其适用问题（单位：元）}
  \label{tab:cbases}
  \footnotesize
  \begin{tabularx}{\linewidth}{C{1.0cm} L{3.6cm} L{3.0cm} C{1.6cm} Y}
    \toprule
    变体 & 口径 & 电价来源 & 适用问题 & 数值 \\
    \midrule
    V1 & 附件 1 典型日，单日周期 & 附件 1（固定价） & 问题 1 & 48\,052.0466 \\
    V2 & 334 天窗口，逐日独立 & 附件 1（固定价） & 问题 2、3 & 16\,407\,319.63 \\
    V3 & 334 天窗口，逐日独立 & 附件 4（波动价） & 问题 4 & 17\,026\,782.86 \\
    \bottomrule
  \end{tabularx}
\end{table}

\subsection{求解结果与机制}

在附件 1 典型日上求解式~\eqref{eq:q1-model}：全天购电量 \(59\,482.6990\)~kWh（\calref{C-02}）、
购电费 \(35\,126.9486\)~元，相对 \(C^{base}=48\,052.0466\)~元（变体 V1）节省
\(12\,925.0980\)~元，\textbf{节省率 \(26.90\%\)}（\calref{C-04}）\evid{C-Q1-03}；
充电量 \(20\,740.6661\)~kWh、放电量 \(16\,799.9396\)~kWh，二式互为倒数——
\(c/d=1.2346=1/(\eta_c\eta_d)\)，其倒数 \(d/c=0.8100=\eta_c\eta_d\) 即往返效率本身（\calref{C-07}）；
最优弃光 \(0.0000\)~kWh（\calref{C-06}）\evid{C-Q1-01}；储电量全程位于
\([1200.0000,10\,800.0000]\)~kWh，两端均为 \(6000.0000\)~kWh（严格日周期）。由
式~\eqref{eq:q1-decomp} 代入上述数值得 \(C^{batt}=-12\,925.0980\)~元，与目标函数直接相减一致。
表~\ref{tab:q1-buy}、表~\ref{tab:q1-soc} 按题面表 1、表 2 格式给出指定时段与 4~h 块的结果
（与 \texttt{result1.xlsx} 逐格一致）。

\begin{table}[htbp]
  \centering
  \caption{微网在指定时间段的购电量及全天的购电量和购电费（问题 1，单位：\unit{kWh} / \unit{元}）}
  \label{tab:q1-buy}
  \small
  \begin{tabular}{cccccc}
    \toprule
    时间段 & 购电量 & 时间段 & 购电量 & 时间段 & 购电量 \\
    \midrule
    \(10{:}00\)--\(10{:}10\) & 0.0000 & \(12{:}00\)--\(12{:}10\) & 480.4124 & \(14{:}00\)--\(14{:}10\) & 0.0000 \\
    \(16{:}00\)--\(16{:}10\) & 445.4317 & \(18{:}00\)--\(18{:}10\) & 531.8940 & \(20{:}00\)--\(20{:}10\) & 0.0000 \\
    \midrule
    \multicolumn{2}{c}{全天购电量（\unit{kWh}）：59\,482.6990} &
    \multicolumn{4}{c}{全天购电费（\unit{元}）：35\,126.9486} \\
    \bottomrule
  \end{tabular}
  \\[1ex]
  {\footnotesize 表注：时段按题面表述取区间左端时刻；\texttt{result1.xlsx} 沿用模板既有的时段末端
  标签，其数值与本文槽位按\textbf{位置}一一对应（槽 \(j\) 写入第 \(j+1\) 个数据行），标签仅作显示。}
\end{table}

\begin{table}[htbp]
  \centering
  \caption{储能设备在指定时间段的充放电量及 \(0{:}00\) 和 \(24{:}00\) 的储电量（问题 1，单位：\unit{kWh}）}
  \label{tab:q1-soc}
  \small
  \begin{tabular}{cccccc}
    \toprule
    时间段 & 充电量 & 放电量 & 时间段 & 充电量 & 放电量 \\
    \midrule
    \(0{:}00\)--\(4{:}00\)   & 4500.0000 & 0.0000    & \(4{:}00\)--\(8{:}00\)   & 833.3333  & 6365.8412 \\
    \(8{:}00\)--\(12{:}00\)  & 4787.9643 & 1702.9970 & \(12{:}00\)--\(16{:}00\) & 5286.0352 & 91.1014 \\
    \(16{:}00\)--\(20{:}00\) & 0.0000    & 5780.1319 & \(20{:}00\)--\(24{:}00\) & 5333.3333 & 2859.8681 \\
    \midrule
    \multicolumn{2}{c}{\(0{:}00\) 储电量：6000.0000} &
    \multicolumn{4}{c}{\(24{:}00\) 储电量：6000.0000} \\
    \bottomrule
  \end{tabular}
  \\[1ex]
  {\footnotesize 表注：每段含 36 个 10~min 槽；6 段之和与逐槽求和一致——充电 20\,740.6661~kWh、
  放电 16\,799.9396~kWh（\calref{C-07}）。}
\end{table}

\begin{figure}[htbp]
  \centering
  \includefig[width=0.54\linewidth]{fig10_strategy_mechanism.pdf}
  \caption{最优策略的机制解释（三子图共享时间轴：光伏与负载功率（\unit{kW}）、
            电价与储电量（\unit{元/kWh}、\unit{kWh}，储电量右轴标出上下界 1200/10800~\unit{kWh}）、
            充放电量（\unit{kWh}）；对象为附件 1 典型日的 144 个 10~\unit{min} 时段，
            竖直虚线连接光伏盈余时段、充电时刻与低价时段）}
  \label{fig:mechanism}
\end{figure}

图~\ref{fig:mechanism} 显示充电时段与光伏盈余时段、低价时段对齐，放电时段与高价时段对齐，
两组时段只\textbf{部分重合}，说明主机制是面向价格的跨时段套利，而非“把午间盈余搬到傍晚”。
可行性复核：供需平衡残差 \(1.49\times10^{-13}\)~kWh、储电量递推残差 \(9.09\times10^{-13}\)~kWh、
日周期偏差 \(\le10^{-12}\)~kWh；越界、同时充放与负购电的时段数均为 0；费用相对复算误差
\(\le10^{-6}\)（\S\ref{sec:audit}）。

\subsection{容量与功率的边际价值、效率口径}

在保持 10\%--90\% 放电深度（\(\underline S=0.1E_{cap}\)、\(\overline S=0.9E_{cap}\) 同步移动）下扫描
额定容量，并单独扫描功率上限（完整档位见 \appref{tab:scan-app}{E 容量与功率扫描明细}）：

\begin{equation}
  \begin{aligned}
    E_{cap}:12\,000\to28\,800\ \text{kWh} &\Rightarrow \Delta C=-6272.42\ \text{元}\ (-17.86\%),\\
    \overline P:5000\to10\,000\ \text{kW} &\Rightarrow \Delta C=-166.36\ \text{元}\ (-0.47\%),\\
    \text{仅解除 }\overline S &\Rightarrow -22.14\%\ (27\,351.55\ \text{元}),\qquad
    \text{仅解除 }\overline P &\Rightarrow -1.66\%\ (34\,544.43\ \text{元}).
  \end{aligned}
  \label{eq:q1-scan}
\end{equation}

界活跃统计：充电功率界在 \(19/144\) 个时段被打满，放电功率界 \(0/144\)；储电量上界活跃
\(21/144\)、下界活跃 \(8/144\)。故\textbf{容量摆幅是主要瓶颈，充放电功率仅在部分时段同时构成约束}，
两者量级相差一个数量级；与 \chapref{chap:verify} 的 Sobol 全局灵敏度一致（详见
\appref{tab:sobol-app}{E Sobol 灵敏度指数}）。

\textbf{效率口径。} 题面“充放电效率为 90\%”未区分充、放方向，故须显式声明口径
（主口径的选择理由见假设 A2）。以 \(\eta_c,\eta_d\) 的四种读法重解，得节省额

\begin{equation}
  \Delta C:\quad
  \underbrace{12\,925.10}_{0.9/0.9\ (\text{主口径})},\quad
  \underbrace{14\,250.99}_{\text{往返 }0.9},\quad
  \underbrace{14\,635.52}_{\text{仅充 }0.9},\quad
  \underbrace{13\,883.44}_{\text{仅放 }0.9}\quad(\text{元}),
  \label{eq:q1-eff}
\end{equation}

极差 \(1710.42\) 元，占主口径节省额的 \(13.2\%\)。故效率口径必须声明，否则同一模型可报出相差
八分之一的节省率。

%!TEX root = ../main.tex
% =============================================================================
%  五、问题 2：不确定性下的两阶段计划购电模型     建议 4–5 页
%  评审在问：不确定性处理合理吗？
%  本节必须出现：① 信息结构声明（R3：附件 3 的 0:00 预报 + 负载已知；结算用附件 2 实际值）
%                ② 两阶段模型全式（第一阶段 E/c/d 与场景无关；第二阶段追索 g）
%                ③ 场景生成与削减三层（残差定义 / 按日总残差分层 / 层内 medoid / 层权重）
%                ④ 报童型安全边际的推导、两界、四条适用边界与实测覆盖率
%                ⑤ 求解结果：计划费 / 紧急费 / 合计 / 紧急电量与触发天数
%                ⑥ 四个指定日期的表 1/表 2/表 3
%                ⑦ 对照：完美信息 R1 与典型日预报 R2
%                ⑧ 紧急购电的时标解释（发生在结算时刻，不在决策时刻）
%  口径与图名以 notes/写作交底.md §六 为准：C-05 在本节首次定义；图名为 fig07_q2_emergency.pdf。
%  红线：报童模型只作"近似刻画 + 上下界 + 方向性解释"，不得声称它是精确解。
% =============================================================================

\section{问题 2：不确定性下的两阶段计划购电模型}

\subsection{问题分析与决策结构}

问题 2 与问题 1 的差别有三点：电价仍逐日相同（附件 1 固定价，均值 0.7662~\unit{元/kWh}），
但负载与光伏\textbf{逐日逐时段变化}；计划仍在每天 \(0{:}00\) 制定；若微网提供的电能低于负载，
须按\textbf{交易时刻电价的 5 倍}紧急购电。故决策分两个层次：

\begin{itemize}
  \item \textbf{第一阶段（日前 \(0{:}00\)）}：只知预报，一次性确定全天 144 个时段的
        \textbf{计划购电量 \(E_t\)}与\textbf{充放电量 \(c_t,d_t\)}，该决策\textbf{不依赖场景}；
  \item \textbf{第二阶段（实时结算）}：实际值揭示后，计划与实际之差由
        \textbf{紧急购电量 \(g_{t,\omega}\)}（缺口，按 \(5p_t\) 计价）与
        \textbf{光伏盈余及超购部分的处置}吸收。
\end{itemize}

\paragraph{信息结构声明}
本文采用\textbf{信息结构 R3}：计划依据为\textbf{附件 3 中当日 \(0{:}00\) 发布的光伏功率预报}
\(\widehat{PV}^{\text{fcst}}_t=\widehat{PV}^{(0{:}00)}_{t}\)；题面未给负载预报，故计划阶段按
"当日负载曲线已知"处理（等价于负载预报无误差），并在假设节显式声明；\textbf{结算依据为附件 2
的当日实际负载与实际光伏}。两者是\textbf{两套数据}，其差异即紧急购电的来源。作为对照，另设完美信息 R1
（\(0{:}00\) 即知当日真值）与典型日预报 R2（附件 1 的日均形状），结果见 \S\ref{sec:q2-compare}。

\paragraph{时标解释：紧急购电发生在结算时刻}
紧急购电不是在某个决策时刻"决定"的，而是在\textbf{结算时刻}由缺口
\(g_{t,\omega}=\max\{0,(L_{\omega,t}-PV_{\omega,t})\Delta t+c_t-d_t-E_t\}\) 唯一确定的：
计划量照常执行，缺口只能以 \(5p_t\) 的价格即时补足，不能回溯改写计划。
这一时标差异是下文报童模型第三条适用边界的来源（\S\ref{sec:q2-newsvendor}）。

\paragraph{一个必须显式声明的建模选择}
\textbf{充放电计划同样在第一阶段确定，实时不允许重调度}。依据是题面"其他时间段的购电费用均按
计划购电量计算"，以及结果文件要求填报逐日逐时段的充放电量。据此，第一阶段买到的高价电不能因
实际光伏更好而转手，计划购电量\textbf{允许高于}预测口径的需求：超出部分既不产生售电收益、
也不再被后续时段使用（超购即被处置），其成本为 \(p_t\)。这一"超购即弃"的设定使模型天然具有
\textbf{安全边际}，也是报童型权衡在本问成立的结构性原因（经济价值见 \S\ref{sec:q2-newsvendor}）。

\subsection{场景生成与削减}

光伏预报误差是本问唯一的不确定性来源。定义\textbf{逐日逐时段残差}
\begin{equation}
  r_{d,t}=\bigl(\widehat{PV}^{\text{fcst}}_{d,t}-PV^{\text{act}}_{d,t}\bigr)\Delta t,
  \qquad d\in\mathcal D,\ t\in\mathcal T ,
  \label{eq:q2-resid}
\end{equation}
单位为 \unit{kWh}，样本为窗口内全部 \(|\mathcal D|=334\) 天 \(\times\,144\) 槽。场景集分三层生成：

\begin{enumerate}
  \item \textbf{按日总残差分层}：按 \(s_d=\sum_{t}r_{d,t}\) 升序把 334 天\textbf{等频}划分为
        \(K\) 层（主口径 \(K=60\)，每层 5--6 天），保留残差水平分布的\textbf{全部区间}（尾部不被均值化抹平）；
  \item \textbf{层内取 medoid}：每层取"到本层其他样本欧氏距离之和最小"的
        \textbf{真实历史日残差向量}（不做中心化与平滑），逐槽保留残差的\textbf{日内形状}
        与时段间相关结构；
  \item \textbf{权重 = 层样本占比}：\(W_\omega=n_\omega/334\)，\(\sum_\omega W_\omega=1\)。
\end{enumerate}

情景 \(\omega\) 下实际净需求 \(ND_{\omega,t}=ND^{\text{fcst}}_t+r_{\omega,t}\)，
其中 \(ND^{\text{fcst}}_t=(L^{\text{fcst}}_t-\widehat{PV}^{\text{fcst}}_t)\Delta t\)。
场景集一次性生成后\textbf{存盘固定}，问题 3、4 与全部消融实验读取同一份，以保证"改进"来自
模型而非场景噪声。



三种层数下总费用极差 13.65 万元，仅为主口径的 \(0.96\%\)，\textbf{结论对场景层数稳健}；
紧急购电量对 \(K\) 更敏感（\(K=40\) 时为 \(K=120\) 时的 1.68 倍）：层数越少、单层样本越多，代表场景越难贴近尾部。场景集退化为"残差恒为零"的单一场景
（\(K=1\)，即严格按预报口径购电）即下文的\textbf{无安全边际}消融（\S\ref{sec:q2-newsvendor}）。

\subsection{两阶段随机规划模型}

\begin{subequations}\label{eq:q2}
  % 目标：第一阶段购电费 + 第二阶段紧急购电费的期望
  \begin{equation}
    \min_{E,c,d,S}\quad
    \sum_{t=1}^{144}p_t E_t
    \;+\;\sum_{\omega\in\Omega}W_\omega\sum_{t=1}^{144}5\,p_t\,g_{t,\omega}
    \label{eq:q2-obj}
  \end{equation}
  % 第一阶段决策（非预期性）
  \begin{equation}
    E_t,\ c_t,\ d_t,\ S_t\ \text{与}\ \omega\ \text{无关},
    \qquad \forall t\in\mathcal T
    \label{eq:q2-nonanticip}
  \end{equation}
  % 计划可行性（允许超购）
  \begin{equation}
    E_t-c_t+d_t\ \ge\ \bigl(L^{\text{fcst}}_t-\widehat{PV}^{\text{fcst}}_t\bigr)\Delta t,
    \qquad \forall t\in\mathcal T
    \label{eq:q2-plan}
  \end{equation}
  % 第二阶段：缺口由紧急购电补齐
  \begin{equation}
    g_{t,\omega}\ \ge\ \bigl(L_{\omega,t}-PV_{\omega,t}\bigr)\Delta t+c_t-d_t-E_t,
    \qquad g_{t,\omega}\ge 0
    \label{eq:q2-gap}
  \end{equation}
  % 储电量递推、界与终端条件
  \begin{equation}
    S_t=S_{t-1}+\eta_c c_t-\frac{d_t}{\eta_d},\quad
    \underline S\le S_t\le\overline S,\quad
    0\le c_t,d_t\le\overline P\Delta t,\quad
    S_0=S^{init},\quad S_{144}\ge S_0
    \label{eq:q2-soc}
  \end{equation}
  % 非负
  \begin{equation}
    E_t\ge 0,\qquad g_{t,\omega}\ge 0
    \label{eq:q2-nonneg}
  \end{equation}
\end{subequations}

式~\eqref{eq:q2-nonanticip} 是两阶段结构的代数表达：带 \(\omega\) 下标的量为\textbf{事后}量
（结算量），第一阶段变量 \(E_t,c_t,d_t,S_t\) 在全部 60 个场景间共享。式~\eqref{eq:q2-plan} 用
\textbf{不等号}：计划购电量只要求不低于"预测口径净需求 + 充电量 − 放电量"；式~\eqref{eq:q2-gap}
的缺口由紧急购电补齐，\(5p_t>0\) 使最优解中 \(g_{t,\omega}\) 取其下界，即
\begin{equation}
  g_{t,\omega}=\max\Bigl\{0,\ \bigl(L_{\omega,t}-PV_{\omega,t}\bigr)\Delta t+c_t-d_t-E_t\Bigr\},
  \label{eq:q2-gapmax}
\end{equation}
是式~\eqref{eq:q2-gap} 的紧形式；式~\eqref{eq:q2-gap} 用的是\textbf{计划充放电量}
\(c_t,d_t\)（不含 \(\omega\)），与 \S5.1 一致。模型规模（\(K=60\)）为每个子问题 \(576+144\times60=9216\)
个变量、144 个等式与 \(144\times(1+60)\) 个不等式，仍是纯 LP，用 SciPy 的 HiGHS 求解器精确求解；
结果窗口内共求解 334 个子问题（逐日顺序推进，日初储电量取前一日末值，1 月作预热期，假设 A8），
终端条件取 \(S_{144}\ge S_0\)（假设 A6）。

\subsection{报童型安全边际的可解释刻画}\label{sec:q2-newsvendor}

暂不考虑储能的跨期耦合，把每个时段看作独立的"报童问题"：计划购电量 \(E_t\) 相当于订货量，
当日实际需求 \(D_t=(L_t-PV_t)\Delta t+c_t-d_t\) 相当于需求量。缺 1~\unit{kWh} 要多付
\((5-1)p_t=4p_t\)（\textbf{缺货成本} \(c_u=4p_t\)，正常电价 \(p_t\) 已计入计划口径）；
多买 1~\unit{kWh} 的损失记为 \(c_o\)。由报童模型，最优计划量是需求分布的\textbf{临界分位数}：
\begin{equation}
  q^{*}_t=F_{D_t}^{-1}\!\left(\frac{c_u}{c_u+c_o}\right)
  =F_{D_t}^{-1}\!\left(\frac{4p_t}{4p_t+c_o}\right),
  \label{eq:q2-newsvendor}
\end{equation}
与之互补的是\textbf{最优缺货概率} \(1-q^{*}_t\)；两者不可混写。

过剩成本 \(c_o\) 的两种口径构成临界分位数的\textbf{上下界}：\textbf{口径一}（超购全废弃）
\(c_o=p_t\)，临界分位数 \(4/5=0.800\)、缺货概率 \(1/5=0.200\)；\textbf{口径二}（超购被储能吸纳
并在后续时段替代购电，单位回收价值 \(p^{low}\)）\(c_o=p_t-p^{low}\approx0.37\)~\unit{元/kWh}，
取全天电价均值 \(p_t=0.7662\)~\unit{元/kWh} 得临界分位数 \(3.0648/(3.0648+0.37)=0.892\)、
缺货概率 \(0.108\approx1/10\)。故\textbf{缺货概率的两个界是 \(1/10\) 与 \(1/5\)}，
覆盖率对应 \(0.892\) 与 \(0.800\)。

\begin{table}[htbp]
  \centering
  \caption{报童近似的两条界与两阶段模型的实测缺货频率（窗口 2025-02-01 至 2025-12-31）}
  \label{tab:newsvendor}
  \footnotesize
  \begin{tabularx}{\linewidth}{L{3.4cm} C{2.0cm} C{2.4cm} C{1.7cm} Y}
    \toprule
    过剩成本口径 & 临界分位数 & 最优缺货概率 & 样本量 & 与实测的关系 \\
    \midrule
    口径一：超购全废弃 \(c_o=p_t\) & 0.800 & 0.200（\(1/5\)） & --- & 缺货概率上界 \\
    口径二：超购可回收 \(c_o\approx0.37\)~\unit{元/kWh} & 0.892 & 0.108（\(\approx1/10\)） & --- & 缺货概率下界 \\
    \midrule
    两阶段随机规划实测 & 覆盖率 0.9145 & \textbf{0.0855} & 48\,096 槽 & 低于下界（更保守） \\
    \bottomrule
  \end{tabularx}
  \\[1ex]
  {\footnotesize 表注：实测缺货频率 = 紧急购电量大于 \(10^{-6}\)~\unit{kWh} 的槽数占 334 天 \(\times\) 144 槽的比例；
  覆盖率 = \(1-\) 缺货频率；电量满足率为 \(99.04\%\)（口径与 \calref{C-05} 一致）。}
\end{table}

两阶段随机规划的实测行为：\textbf{逐槽缺货频率 8.55\%}（覆盖率 91.45\%）低于下界 \(1/10\)，
\textbf{电量满足率 99.04\%}，全窗口\textbf{安全边际合计 576\,299.67~kWh}（平均 11.98~\unit{kWh}/槽）
\evid{C-Q2-01}。实测更保守的原因在结构而非参数：\textbf{储能同时承担对冲功能}，同一份计划可由
电池放电在缺口时段补足，单期报童模型无法刻画"计划购电量"与"电池吞吐"的替代关系。
故式~\eqref{eq:q2-newsvendor} 与表~\ref{tab:newsvendor} 只作\textbf{方向性近似与上下界}，
\textbf{不是}精确最优解。

安全边际的经济价值由\textbf{无安全边际消融}量化：场景集退化为单一情景（\(K=1\)，残差恒为零，
严格按预报口径购电）后，全年计划购电费 12\,570\,533.91~\unit{元}、紧急购电费
4\,828\,025.76~\unit{元}、合计 17\,398\,559.68~\unit{元}，比无储能基线 16\,407\,319.63~\unit{元}
高 \(6.04\%\)，紧急购电量 1\,156\,618.02~\unit{kWh}、334/334 天全部触发；主口径费用
14\,264\,792.13~\unit{元}、紧急购电量 172\,015.82~\unit{kWh}。分解为：\textbf{安全边际把计划购电费
推高 1\,067\,018.99~\unit{元}（边际电量 576\,299.67~\unit{kWh} 按均价 0.6176~\unit{元/kWh} 计
355\,918.38~\unit{元}，其余 711\,100.60~\unit{元} 来自储能排产与购电时序的同步调整），
却把紧急购电费压低 4\,200\,786.52~\unit{元}，净收益 3\,133\,767.55~\unit{元}
（相对不设边际降低 \(18.0\%\)）}。即边际电量只占全窗口计划购电量 22\,220\,333.14~\unit{kWh} 的
\(2.59\%\)，却消掉 \(87.0\%\) 的紧急购电费。惩罚倍数灵敏度方向一致：同一典型日与同一残差场景集上
把倍数由 \(k=2\) 提高到 \(k=10\)，单日最优安全边际由 671~\unit{kWh} 增至 3\,334~\unit{kWh}，
期望紧急购电量由 1\,246~\unit{kWh} 降至 187~\unit{kWh}，目标值由 39\,224.54~\unit{元} 升至
41\,758.99~\unit{元}（\(k=5\) 时为 40\,747.45~\unit{元}）。

报童模型在本问题中\textbf{不是精确解}，其偏离来源有四条，也就是它的适用边界：

\begin{enumerate}
  \item \textbf{多时段耦合}：144 个时段通过储电量递推耦合，本时段多买的电会占用电池容量、
        影响其后所有时段，不能拆成 144 个独立报童问题；
  \item \textbf{储能缓冲}：过剩电量可被储能吸纳并在后续高价时段替代购电，故 \(c_o\) 严格小于
        \(p_t\)，临界分位数应高于口径一的 0.800；
  \item \textbf{紧急购电的时标在结算时刻}：缺口只能事后按 \(5p_t\) 补足，受当日计划充放电量与
        放电功率限制，不能回溯改写计划；
  \item \textbf{残差分布未知}：本文用 334 天样本的\textbf{样本分位数}替代总体分位数，
        每个时段仅 334 个观测，分位数估计带误差，须报置信区间。
\end{enumerate}



\subsection{对照方案与风险维度}\label{sec:q2-compare}

在同一批实际轨迹上比较三种计划依据：

\begin{itemize}
  \item \textbf{R1 完美信息}：\(0{:}00\) 即知当日实际负载与实际光伏，紧急购电量恒为 0。
        在与主方案相同的日决策结构下逐日求解，总费用 \textbf{13\,411\,560.90~\unit{元}}，
        相对无储能基线节省 \(18.26\%\)，是"预测误差成本"的\textbf{下界}；若放弃日决策解耦、
        把 334 天写成单一 LP（单一储电量轨迹），费用 12\,228\,166.56~\unit{元}，为更松的下界参考；
  \item \textbf{R3 本文主方案}（附件 3 的 \(0{:}00\) 预报）：总费用 14\,264\,792.13~\unit{元}，
        相对基线节省 \(13.06\%\)；
  \item \textbf{R2 典型日预报}：用附件 1 的日均形状作预报，总费用
        \textbf{23\,484\,646.63~\unit{元}}，比无储能基线高 \(43.14\%\)，紧急购电量
        2\,525\,548.76~\unit{kWh}（286/334 天触发）。附件 1 的数值是附件 2 实际功率的
        \textbf{全年逐时段均值}而非日前预报，只用于说明"预报质量决定成本"，\textbf{不可作为主方案}。
\end{itemize}

R3 与 R1 之差 \(14\,264\,792.13-13\,411\,560.90=853\,231.23\)~\unit{元}（相对 R1 高 \(6.36\%\)）
即\textbf{预测误差的成本}，即"5 倍紧急购电"机制对应的金额。

风险维度上，期望费用最小化对极端场景不敏感：在 334 天中每 6 天取 1 天的子样本（56 天）上做
CVaR（置信水平 \(\alpha=0.95\)，\calref{C-11}）与 \(\lambda\) 扫描，\(\lambda\) 由 0 增至 1 时
期望费用由 41\,995.59~\unit{元}/日 升至 42\,411.56~\unit{元}/日（\(+0.99\%\)），CVaR\(_{0.95}\)
由 52\,176.36~\unit{元}/日 降至 50\,510.46~\unit{元}/日（\(-3.19\%\)），日均紧急购电量由
480.2~\unit{kWh} 降至 313.8~\unit{kWh}；完整前沿与回测见 \S7.4，本节只保留与日前决策直接相关的
R1/R2/R3 对照。

\subsection{求解结果}

在 2025-02-01 至 2025-12-31（334 天）窗口内求解式~\eqref{eq:q2}，总费用
\textbf{14\,264\,792.13~\unit{元}}（1\,426.48 万元），其中\textbf{计划购电费
13\,637\,552.90~\unit{元}、紧急购电费 627\,239.24~\unit{元}}；相对无储能基线
\(C^{base}=16\,407\,319.63\)~\unit{元}（\calref{C-01} 变体 V2）节省 2\,142\,527.50~\unit{元}，
节省率 \(13.06\%\)（\calref{C-04}），两项费用占比 \(95.60\%\) 与 \(4.40\%\)（\calref{C-03}）。
全窗口计划购电量 22\,220\,333.14~\unit{kWh}（\calref{C-02}），紧急购电量 172\,015.82~\unit{kWh}，
出现在 296/334 天（\calref{C-05}）；逐槽缺货频率 8.55\%、电量满足率 99.04\%、安全边际合计
576\,299.67~\unit{kWh}，与 \S\ref{sec:q2-newsvendor} 一致。

\begin{caliber}{C-05}{紧急购电量 \(G\)}
\textbf{定义}\quad \(G=\sum_{d\in\mathcal D}\sum_{t=1}^{144}g_{d,t}\)，
其中 \(g_{d,t}\) 由式~\eqref{eq:q2-gapmax} 在\textbf{结算时刻}确定，
是"含电池吞吐项"的完整缺口，而不是 \(\max\{0,\text{实际净负荷}-\text{计划购电量}\}\)；%
\textbf{数据}\quad 结算依据为附件 2 的实际负载与实际光伏，
计划量与充放电量取 \texttt{result2.xlsx} 的对应工作表；%
\textbf{时段}\quad 334 天 \(\times\) 144 槽（2025-02-01 至 2025-12-31）；%
\textbf{单位}\quad \unit{kWh}（表内保留 2 位小数）；%
\textbf{不含}\quad 不含问题 3 调整机制引入的缺口（调整量单独统计，\calref{C-09}）；%
\textbf{对照}\quad 完美信息 R1 下为 0.0000~\unit{kWh}（0/334 天）；
典型日预报 R2 下为 2\,525\,548.76~\unit{kWh}（286/334 天）；%
\textbf{复算}\quad \texttt{code/lp\_engine.py::realized\_slot}，审计项 A1 与 A9。
\end{caliber}

\paragraph{可行性复核}
逐日检查：计划口径残差（供给缺口）最大 \(3.41\times10^{-13}\)~\unit{kWh}，储电量净递推残差最大
\(1.09\times10^{-11}\)~\unit{kWh}，最差单日综合残差 \(1.82\times10^{-12}\)；同时充放、SOC 越界、
功率越界、负购电的时段数均为 0，末值低于初值的天数 0、费用相对复算误差 0.0000；
结果文件回读校验的全天购电量最大偏差 0.0011~\unit{kWh}、全天购电费最大相对偏差
\(3.40\times10^{-7}\)。

\paragraph{题面指定日期的结果（表 1 / 表 2 / 表 3 格式）}
四个指定日期 2025-03-20、2025-06-21、2025-09-23、2025-12-21 的计划购电量
（表~\ref{tab:q2-buy}）、储能充放电量与 \(0{:}00\)/\(24{:}00\) 储电量（表~\ref{tab:q2-soc}）、
紧急购电量（表~\ref{tab:q2-emergency}）见表。四日均满足计划口径供需平衡
（最大残差 \(1.82\times10^{-12}\)）与日周期终端条件，\(0{:}00\) 与 \(24{:}00\) 储电量均为
6000.0000~\unit{kWh}。

\begin{table}[htbp]
  \centering
  \caption{微网在指定时段与指定日期的购电量、全天购电量和购电费（问题 2；电量单位：\unit{kWh}，费用单位：\unit{元}）}
  \label{tab:q2-buy}
  \small
  \begin{tabular}{lrrrr}
    \toprule
    时段／指标 & 2025-03-20 & 2025-06-21 & 2025-09-23 & 2025-12-21 \\
    \midrule
    \(10{:}00\)--\(10{:}10\) & 0.0000 & 0.0000 & 0.0000 & 0.0000 \\
    \(12{:}00\)--\(12{:}10\) & 804.3029 & 0.0000 & 536.2415 & 997.0335 \\
    \(14{:}00\)--\(14{:}10\) & 0.0000 & 0.0000 & 0.0000 & 0.0000 \\
    \(16{:}00\)--\(16{:}10\) & 817.3061 & 255.2254 & 902.7969 & 1001.2240 \\
    \(18{:}00\)--\(18{:}10\) & 697.9973 & 0.0000 & 764.3298 & 715.4144 \\
    \(20{:}00\)--\(20{:}10\) & 0.0000 & 0.0000 & 0.0000 & 0.0000 \\
    \midrule
    全天购电量 & 74\,178.1298 & 35\,701.5391 & 74\,744.7963 & 99\,692.6601 \\
    全天购电费 & 45\,229.90 & 20\,324.26 & 47\,063.39 & 64\,376.98 \\
    紧急购电量 & 1025.8805 & 187.0843 & 193.3669 & 22.5250 \\
    全天总费用 & 49\,654.02 & 21\,255.45 & 47\,717.70 & 64\,432.28 \\
    \bottomrule
  \end{tabular}
  \\[1ex]
  {\footnotesize 表注：时段按题面表述取区间左端时刻；
  全天购电费为计划口径（\(\sum_t p_tE_t\)），全天总费用为计划购电费与紧急购电费之和；
  四个日期的无储能基线分别为 52\,290.65、29\,668.18、54\,741.30、71\,176.92~\unit{元}。}
\end{table}

\begin{table}[htbp]
  \centering
  \caption{储能设备在指定日期的充放电量及 \(0{:}00\) 与 \(24{:}00\) 的储电量（问题 2，单位：\unit{kWh}）}
  \label{tab:q2-soc}
  \footnotesize
  \begin{tabular}{L{2.0cm}L{2.1cm}rr@{\quad}L{2.0cm}L{2.1cm}rr}
    \toprule
    日期 & 时段 & 充电量 & 放电量 & 日期 & 时段 & 充电量 & 放电量 \\
    \midrule
    2025-03-20 & \(0{:}00\)--\(4{:}00\)   & 5333.3333 & 828.0303  & 2025-09-23 & \(0{:}00\)--\(4{:}00\)   & 5333.3333 & 790.6233 \\
    2025-03-20 & \(4{:}00\)--\(8{:}00\)   & 6500.0000 & 7811.9697 & 2025-09-23 & \(4{:}00\)--\(8{:}00\)   & 5833.3333 & 7778.1258 \\
    2025-03-20 & \(8{:}00\)--\(12{:}00\)  & 6905.5931 & 2218.5304 & 2025-09-23 & \(8{:}00\)--\(12{:}00\)  & 6728.4081 & 1606.2615 \\
    2025-03-20 & \(12{:}00\)--\(16{:}00\) & 5333.3333 & 8640.0000 & 2025-09-23 & \(12{:}00\)--\(16{:}00\) & 5333.3333 & 8640.0000 \\
    2025-03-20 & \(16{:}00\)--\(20{:}00\) & 0.0000    & 0.0000    & 2025-09-23 & \(16{:}00\)--\(20{:}00\) & 0.0000    & 0.0000 \\
    2025-03-20 & \(20{:}00\)--\(24{:}00\) & 0.0000    & 0.0000    & 2025-09-23 & \(20{:}00\)--\(24{:}00\) & 0.0000    & 0.0000 \\
    \midrule
    2025-06-21 & \(0{:}00\)--\(4{:}00\)   & 0.0000    & 2723.8246 & 2025-12-21 & \(0{:}00\)--\(4{:}00\)   & 5333.3333 & 0.0000 \\
    2025-06-21 & \(4{:}00\)--\(8{:}00\)   & 5348.6522 & 1710.2113 & 2025-12-21 & \(4{:}00\)--\(8{:}00\)   & 6666.6667 & 9315.0000 \\
    2025-06-21 & \(8{:}00\)--\(12{:}00\)  & 5458.7995 & 0.0000    & 2025-12-21 & \(8{:}00\)--\(12{:}00\)  & 8333.3333 & 2835.0000 \\
    2025-06-21 & \(12{:}00\)--\(16{:}00\) & 5333.3333 & 8640.0000 & 2025-12-21 & \(12{:}00\)--\(16{:}00\) & 5333.3333 & 8640.0000 \\
    2025-06-21 & \(16{:}00\)--\(20{:}00\) & 0.0000    & 0.0000    & 2025-12-21 & \(16{:}00\)--\(20{:}00\) & 0.0000    & 0.0000 \\
    2025-06-21 & \(20{:}00\)--\(24{:}00\) & 0.0000    & 0.0000    & 2025-12-21 & \(20{:}00\)--\(24{:}00\) & 0.0000    & 0.0000 \\
    \bottomrule
  \end{tabular}
  \\[1ex]
  {\footnotesize 表注：四个指定日期的 \(0{:}00\) 与 \(24{:}00\) 储电量均为 6000.0000~\unit{kWh}；
  每段含 36 个 10~\unit{min} 槽，充放电量为交流侧电量（\calref{C-07}）；
  各日充放电量之比均为 1.2346（对应往返效率 0.81）。}
\end{table}

\begin{table}[htbp]
  \centering
  \caption{微网在指定日期的紧急购电量（问题 2；购电量单位：\unit{kWh}，电价单位：\unit{元/kWh}）}
  \label{tab:q2-emergency}
  \footnotesize
  \begin{tabular}{llrr}
    \toprule
    日期 & 购电时间段 & 紧急购电量 & 交易时刻电价（加权平均） \\
    \midrule
    2025-03-20 & \(6{:}00\)--\(8{:}00\)   & 228.9914 & 0.8625 \\
    2025-03-20 & \(8{:}10\)--\(8{:}30\)   & 49.1586  & 0.8625 \\
    2025-03-20 & \(8{:}50\)--\(10{:}30\)  & 531.2739 & 0.8625 \\
    2025-03-20 & \(10{:}40\)--\(11{:}40\) & 216.4567 & 0.8625 \\
    \midrule
    2025-06-21 & \(4{:}00\)--\(4{:}50\)   & 5.2416   & 0.9955 \\
    2025-06-21 & \(14{:}00\)--\(14{:}30\) & 181.8427 & 0.9955 \\
    \midrule
    2025-09-23 & \(6{:}50\)--\(7{:}40\)   & 132.5948 & 0.6768 \\
    2025-09-23 & \(11{:}00\)--\(11{:}40\) & 60.7721  & 0.6768 \\
    \midrule
    2025-12-21 & \(11{:}00\)--\(11{:}30\) & 22.5250 & 0.4910 \\
    \midrule
    \multicolumn{2}{l}{四个日期合计} & 1428.8567 & --- \\
    \bottomrule
  \end{tabular}
  \\[1ex]
  {\footnotesize 表注：本表按"日期 \(\times\) 紧急购电事件"逐行展开，与题面表 3 的横向版式一一对应；
  同一日期可有多行，连续触发的相邻槽按事件聚合，时间段为该事件涉及的槽位区间；
  交易时刻电价为该事件内按电量加权的平均电价（事件电价按 \(5p_t\) 计价，本列为 \(p_t\)）。
  无紧急购电的时段不列入本表（全窗口 334 天中 296 天有触发）。}
\end{table}

%!TEX root = ../main.tex
% =============================================================================
%  六、问题 3：多阶段滚动调整模型      目标 3.5 页
%  评审在问：多阶段做对了吗？
%  本节含：① 后覆盖前的决策结构与非预期性的实现（式 (26) 场景树约束）
%          ② 分段线性凸偏差费用的精确线性化（凸性 → 无需 0-1 变量）
%          ③ 多阶段滚动模型与 MPC 实现（算法 1）
%          ④ 信息价值 VoI：直接回答"是否需要引入其他时刻预报"
%          ⑤ 结果与口径依赖
%  口径：C-09 在本节首次定义；表 tab:q3-epochs 给出四段槽位划分
% =============================================================================

\section{问题 3：多阶段滚动调整模型}\label{chap:q3}

\subsection{后覆盖前的决策结构}

问题 3 在问题 2 上增加两个要素：四个时刻 \(m\in\{0,6,12,18\}\) 均有未来 24~h 整点光伏预报，决策随之
\textbf{滚动修正}；并引入\textbf{两类不对称惩罚}（总费用含计划购电费、紧急购电费与调整相关费用）。
沿用 \chapref{chap:data} 的槽位读法：槽位 \(t=0,\dots,143\) 表示 \([10t,10(t+1))\) 分钟，
\(\Delta t=1/6\)~h。决策量为 \(0{:}00\) 制定的\textbf{计划购电量} \(E^{\text{plan}}_t\)（第一阶段承诺）、
滚动修正后的\textbf{调整购电量} \(E^{\text{adj}}_t\)、由实际负载与光伏决定的\textbf{紧急购电量} \(g_t\)，
以及充放电量 \(c_t,d_t\) 与时段末储电量 \(S_t\)。

四个预报时刻把一天切成四段，槽位 \(t\) 的最终执行量由\textbf{最后一次覆盖它的那次决策}给出
（表~\ref{tab:q3-epochs}）：\(0{:}00\) 用当次预报制定全天 144 槽计划，此后每个决策时刻只用\textbf{当次
预报}对\textbf{其后至 \(24{:}00\) 的剩余时段}重新决策，但\textbf{只执行到下一个决策时刻}。

\begin{table}[htbp]
  \centering
  \caption{四次预报的决策时刻、生效槽位与可用预报（时间范围：2025-02-01 至 12-31 的每一天；
            槽位单位：10 min 槽，共 144 槽；预报来源：附件 3）}
  \label{tab:q3-epochs}
  \footnotesize
  \setlength{\tabcolsep}{4pt}
  \begin{tabularx}{\linewidth}{C{1.6cm} L{3.0cm} C{2.5cm} C{2.4cm} Y}
    \toprule
    决策时刻 & 决策内容 & 生效槽位 & 生效时段 & 重解范围（只执行到下一时刻） \\
    \midrule
    \(0{:}00\) & 计划购电量 \(E^{\text{plan}}\)（144 槽） & \(0\)--\(35\) & \(0{:}00\)--\(6{:}00\) & \(0\)--\(143\) \\
    \(6{:}00\)  & 调整购电量 \(E^{\text{adj}}\) & \(36\)--\(71\) & \(6{:}00\)--\(12{:}00\) & \(36\)--\(143\) \\
    \(12{:}00\) & 调整购电量 \(E^{\text{adj}}\) & \(72\)--\(107\) & \(12{:}00\)--\(18{:}00\) & \(72\)--\(143\) \\
    \(18{:}00\) & 调整购电量 \(E^{\text{adj}}\) & \(108\)--\(143\) & \(18{:}00\)--\(24{:}00\) & \(108\)--\(143\) \\
    \bottomrule
  \end{tabularx}
\end{table}

\textbf{非预期性在实现层面自然满足。} 非预期性要求任一时刻的决策只依赖该时刻及之前已实现的信息；
本文实现有两个环节天然满足，无需显式施加式~\eqref{eq:q3-nonanticip}：其一，取预报值的函数只按
\textbf{提前期}索引——槽位 \(t\) 在时刻 \(m\) 取附件 3 中 \(m\) 时刻发布行的第
\(k=(h_t-m)\bmod24\) 列，故 \(m\) 时刻的子问题\textbf{在数据结构上取不到}其后发布的列；其二，
每个子问题只覆盖 \(t\ge s_m=6m\)，已执行槽位被冻结，决策不可能回头改写 \(t<s_m\) 的历史。
为可证伪地检验这一点，本文做\textbf{预报置换检验}：把某一次预报在 12 天之间随机置换后重解，
记录四个控制区内决策量的最大变化。结果是置换 \(0{:}00\) 预报会显著改变全天决策（说明检验有效力），
而\textbf{任何一次置换都不改变早于该预报发布时刻的槽位决策}（该区变化恰为 \(0.000\)~kWh）；
置换 \(18{:}00\) 预报时四区变化全为 \(0.000\)，这是数据事实而非检验失效——该次预报覆盖的槽
\(108\)--\(143\) 上几乎恒为 0（\(19{:}00\) 以后全年恒为 0），对晚间决策本无新信息
（逐区明细见 \appref{tab:na-app}{E 非预期性置换明细}）。

\subsection{偏差费用的凸性与精确线性化}

偏差费用 \(c^{\text{buy}}_t\) 关于 \(E^{\text{adj}}_t\) 是\textbf{凸的分段线性函数}：调增（\allowbreak
\(E^{\text{adj}}_t\ge E^{\text{plan}}_t\)）斜率为 \(1.5p_t\)，调减斜率为 \(0.5p_t\)，故

\begin{equation}
  c_t^{\text{buy}}
  =p_t\min\bigl(E^{\text{plan}}_t,E^{\text{adj}}_t\bigr)
  +0.5p_t\bigl(E^{\text{plan}}_t-E^{\text{adj}}_t\bigr)^{+}
  +1.5p_t\bigl(E^{\text{adj}}_t-E^{\text{plan}}_t\bigr)^{+}.
  \label{eq:q3-cost}
\end{equation}

凸性意味着\textbf{费用只惩罚偏离的方向}：调减按半价结算差额、调增按 1.5 倍结算。引入
\(u_t,v_t\ge0\) 表示超额量与违约量，

\begin{equation}
  E^{\text{adj}}_t-E^{\text{plan}}_t=u_t-v_t,\qquad
  c_t^{\text{buy}}=p_tE^{\text{plan}}_t+1.5p_tu_t-0.5p_tv_t ,
  \label{eq:q3-lin}
\end{equation}

后式由 \(p_t\min(E^{\text{plan}}_t,E^{\text{adj}}_t)=p_tE^{\text{plan}}_t-p_tv_t\) 得到。该线性化是
\textbf{精确的}，无需 0--1 变量：费用中与 \((u_t,v_t)\) 有关的部分为 \(1.5p_tu_t-0.5p_tv_t\)，
若 \(u_t,v_t\) 同时为正，两边同减 \(\min(u_t,v_t)\) 可使目标下降 \(p_t\min(u_t,v_t)>0\)，
故最优解必有 \(\min(u_t,v_t)=0\)。互斥性的关键是两条分段斜率跨越基准电价 \(1\times p_t\)
（\(0.5<1<1.5\)），使 \(u_t,v_t\) 的有效系数均为正，最优解只能落在折点上。
\textbf{适用边界}：若违约倍数取到 \(1\) 及以上（本题为 \(0.5\)），折点两侧斜率不再跨过基准电价，
互斥性不再自动成立，必须引入 0--1 变量。数值上，\(334\times144=48\,096\) 个“槽-日”中
\(\min(u_t,v_t)>10^{-6}\)~kWh 的违反数为 \textbf{0}。

\subsection{非预期性约束与多阶段滚动模型}\label{sec:q3-mpc}\label{sec:q3-roll}

以四个决策时刻为分支点构造\textbf{场景树}，记 \(\tau(t)\) 为不晚于槽位 \(t\) 的最后一个决策时刻。
\(0{:}00\) 的计划量对所有场景相同；槽位 \(t\) 的调整量只能依赖截至 \(\tau(t)\) 的信息，即

\begin{equation}
  E^{\text{adj}}_t(\omega)=E^{\text{adj}}_t(\omega'),
  \qquad
  \forall\,\omega,\omega'\in\Omega:\ \
  \omega|_{[0,\tau(t)]}=\omega'|_{[0,\tau(t)]},
  \label{eq:q3-nonanticip}
\end{equation}

其中 \(\omega|_{[0,\tau]}\) 表示场景 \(\omega\) 在时刻 \(\tau\) 之前的历史。式~\eqref{eq:q3-nonanticip}
展开成 LP，等价于按场景树中\textbf{深度为 \(\tau(t)\) 的节点}把场景分组、对同组场景施加变量相等的线性
约束。本文不显式展开场景树，而采用与它在确定性等价下一致的滚动实现（\S\ref{sec:q3-mpc}）：
变量规模与场景数解耦，非预期性由“数据可取范围 \(+\) 已执行槽位冻结”自动满足。

记 \(\Omega_m\) 为时刻 \(m\) 用于剩余时段的场景集（权重 \(w_\omega\)，按日总预报残差分 60 层、
层内取 medoid、权重取层样本占比，与 \ref{sec:q2-scen} 同一份存盘场景集），\(s_m=6m\)。综合
式~\eqref{eq:q3-cost}--\eqref{eq:q3-nonanticip}，问题 3 的多阶段模型为

\begin{subequations}\label{eq:q3}
\begin{align}
  \min\quad
  & \sum_{t=0}^{143}p_tE^{\text{plan}}_t
    +\sum_{\omega\in\Omega}w_\omega\sum_{t=0}^{143}
    \Bigl[1.5p_tu_t(\omega)-0.5p_tv_t(\omega)+5p_tg_t(\omega)\Bigr] \label{eq:q3-obj}\\
  \text{s.t.}\quad
  & E^{\text{adj}}_t(\omega)-E^{\text{plan}}_t=u_t(\omega)-v_t(\omega),\quad u_t(\omega),v_t(\omega)\ge0, && \forall t,\omega \label{eq:q3-diff}\\
  & E^{\text{plan}}_t-c_t+d_t\ \ge\ \bigl(L_t-\widehat{PV}^{(m)}_t\bigr)\Delta t, && \forall t \label{eq:q3-planfeas}\\
  & g_t(\omega)\ \ge\ \bigl(L_{\omega,t}-PV_{\omega,t}\bigr)\Delta t+c_t-d_t-E^{\text{adj}}_t(\omega),\quad g_t(\omega)\ge0, && \forall t,\omega \label{eq:q3-gap}\\
  & S_t=S_{t-1}+\eta_cc_t-d_t/\eta_d,\quad \underline S\le S_t\le\overline S,\quad 0\le c_t,d_t\le\overline P\Delta t,\quad S_0=S^{init},\ S_{143}\ge S_0, && \forall t \label{eq:q3-soc}\\
  & E^{\text{plan}}_t\ \text{与}\ \omega\ \text{无关};\qquad E^{\text{adj}}_t\ \text{满足式}~\eqref{eq:q3-nonanticip}. && \label{eq:q3-na-ref}
\end{align}
\end{subequations}

式~\eqref{eq:q3-gap} 的缺口以\textbf{调整后}的 \(E^{\text{adj}}_t(\omega)\) 为基准（假设 A7）；
终端条件 \(S_{143}\ge S_0\) 防止滚动优化在窗口末端放空储能换取短期收益（假设 A6）。参数取
\(\eta_c=\eta_d=0.9\)、\(\overline P\Delta t=833.3333\)~kWh、\([\underline S,\overline S]=[1200,10\,800]\)~kWh、
\(S^{init}=6000\)~kWh。

\textbf{MPC 实现。} 式~\eqref{eq:q3} 的场景树 LP 按决策时刻分解为四个子问题：在时刻 \(m\) 以当前
储电量 \(S_{s_m-1}\) 为初值、当日初值 \(S_0\) 为终端下界，对 \(t\ge s_m\) 求解

\begin{equation}
  \min\ \sum_{t=s_m}^{143}
  \Bigl[(1-\delta_m)p_tE_t+\delta_m\bigl(1.5p_tu_t-0.5p_tv_t\bigr)\Bigr]
  +\sum_{\omega\in\Omega_m}w_\omega\sum_{t=s_m}^{143}5p_tg_{t,\omega},
  \qquad
  \delta_m=\begin{cases}0,&m=0\\ 1,&m>0\end{cases}
  \label{eq:q3-mpc}
\end{equation}

约束为式~\eqref{eq:q3-diff}--\eqref{eq:q3-soc} 在剩余槽位上的限制；\(m>0\) 时计划项已是常数，
故从目标中略去。式~\eqref{eq:q3-mpc} 把“重解范围”与“执行范围”分离：求解后\textbf{只把}
\([s_m,s_{m+1})\) 上的决策写入实际轨迹，其余解丢弃，待下一个预报时刻用更新的预报重优化——充放电量
因此随信息滚动修正，而非冻结在 \(0{:}00\)。全部物理约束由结果文件独立复算：供给缺口残差
\(\le3.5\times10^{-13}\)~kWh、储电量递推残差 \(\le1.1\times10^{-11}\)~kWh、储电量与功率越界数均为 0；
每个结果日重解 4 次子问题（共 1336 次，另加 1 月预热期），每次均为 LP 并由 HiGHS 精确求解。
完整流程见算法~\ref{alg:rolling}。

\begin{algorithm}[htbp]
  \caption{微网日前计划——多阶段滚动调整的一体化求解流程}
  \label{alg:rolling}
  \small
  \begin{algorithmic}[1]
    \Require 电价 \(p_t\)、负载 \(L_t\)、实际光伏 \(P^{act}_t\)、各发布时刻预报 \(P^{fe}_t\)
             (\(e\in\{0,6,12,18\}\))、储能参数 \(\eta_c,\eta_d,\overline P,\underline S,\overline S\)、
             初值 \(S_0=6000\)~kWh
    \Ensure 计划购电量 \(E^{\text{plan}}\)、调整购电量 \(E^{\text{adj}}\)、实际充放电量 \(c,d\)、
            紧急购电量 \(g\)、储电量轨迹 \(S\)
    \For{每一天 \(d\)}
      \State \textbf{计划层}：以 \(0{:}00\) 预报为光伏输入求解式~\eqref{eq:q2-obj}，得
             \(E^{\text{plan}},c,d,S\)（问题 2/4-2 到此为止）
      \For{每个调整时刻 \(e\in\{6,12,18\}\)}
        \State \textbf{调整层}：以该时刻预报重解\textbf{剩余时段} \([6e/6,\,144)\) 的 LP，
               目标含偏差惩罚与 \(5p_t\) 紧急购电，约束加\textbf{非预期性}
               （槽位 \(<6e/6\) 的决策冻结）
        \State 后覆盖前：该次解覆盖其后所有槽位的执行值
      \EndFor
      \State \textbf{结算}：按附件 2 实际负载与光伏计算缺口
             \(g_t=\bigl[(L_t-P^{act}_t)\Delta t+c_t-d_t-E^{\text{adj}}_t\bigr]^{+}\)，
             紧急购电费 \(5\sum_tp_tg_t\)
      \State \textbf{一致性校验}：核验电量平衡、储电量递推、区间与功率界、
             充放互斥与式~\eqref{eq:q1-energy} 的能量恒等式
      \State \textbf{递推}：\(S_0^{(d+1)}\leftarrow S_{144}^{(d)}\)
    \EndFor
    \State \Return 全部策略与费用构成
  \end{algorithmic}
\end{algorithm}

\subsection{信息价值：是否需要引入其他时刻的预报}\label{sec:voi}

题面直接追问该问题，故须给出\textbf{量化答案}与\textbf{成本侧}。设计\textbf{预报时刻数对照实验}：
\(n\in\{1,2,4\}\) 三组使用\textbf{同一场景集、同一实际轨迹回放、同一费用口径}（\calref{C-09}），
只有“计划时能用到几次预报”不同——\(n=1\) 只用 \(0{:}00\) 预报（即问题 2 的结构）、\(n=2\) 用
\(0{:}00\) 与 \(12{:}00\)、\(n=4\) 为题面给定的四次预报。附件 3 只提供这四个发布时刻，故
\(n=8,24\) 在本数据下\textbf{并不存在}：任何加密实验都须自行合成预报，其结论取决于合成规则而非
数据。本文不做这种合成，改用\textbf{完全信息上界}：\(0{:}00\) 计划仍按 \(0{:}00\) 预报制定，
但从 \(6{:}00\) 起的调整假定已知当日实际净需求，其费用\textbf{不高于}任何有限预报时刻数配置，
是“还能再省多少”的\textbf{与数据无关的上界}。记

\begin{equation}
  \mathrm{VoI}(n)=\frac{C(1)-C(n)}{C(1)}\times100\%,
  \label{eq:q3-voi}
\end{equation}

其中 \(C(n)\) 为可用 \(n\) 个预报时刻时、同一场景集同一实际轨迹下的总费用。结果见表~\ref{tab:voi}

\begin{table}[htbp]
  \centering
  \caption{预报时刻数对照实验：总费用、边际节省与信息价值
           （同一场景集、同一实际轨迹、同一偏差量定价\calref{C-09}；
            时间范围：2025-02-01 至 12-31 共 334 天；样本：60 个分层 medoid 代表场景 \ensuremath{\times} 4 组配置；
            单位：费用 元，电量 kWh）}
  \label{tab:voi}
  \footnotesize
  \setlength{\tabcolsep}{4pt}
  \begin{tabularx}{\linewidth}{L{3.0cm} C{2.1cm} C{1.9cm} C{1.6cm} C{1.8cm} Y}
    \toprule
    可用预报时刻 & 总费用/\unit{元} & 边际节省/\unit{元} & 累计 \(\mathrm{VoI}\) & 已实现占比 & 紧急购电量/\unit{kWh} \\
    \midrule
    仅 \(0{:}00\)（\(n=1\)） & 14\,264\,792.13 & --- & 0 & 0 & 172\,015.82 \\
    \(0{:}00\)、\(12{:}00\)（\(n=2\)） & 14\,131\,564.16 & 133\,227.98 & 0.93\% & 34.1\% & 122\,524.83 \\
    四时刻（\(n=4\)，题面给定） & 13\,932\,402.95 & 199\,161.20 & 2.33\% & 85.1\% & 64\,522.09 \\
    完全信息（可获信息上界） & 13\,874\,340.39 & 58\,062.56 & 2.74\% & 100\% & 6\,134.76 \\
    \bottomrule
  \end{tabularx}
  \\[1ex]
  {\footnotesize 表注：“已实现占比”的分母 \(C(1)-C_{\mathrm{PI}}=390\,451.74\) 元是完全信息下最多能省的
  金额。}
\end{table}

结论有三条。① \textbf{滚动调整有效}：总费用降 \(2.33\%\)、紧急购电量降 \(62.49\%\)。
② \textbf{收益边际递减}：\(1\to2\) 次加密省 \(133\,227.98\) 元、\(2\to4\) 次省 \(199\,161.20\) 元，
\(n=4\) 已实现可获得信息价值的 \(\mathbf{85.1\%}\)；\textbf{再增加预报最多只能再省}
\(58\,062.56\) 元（占 \(n=4\) 总费用的 \(0.42\%\)）。③ \textbf{边际增益主要来自 \(6{:}00\)
对上午的更新}：\(199\,161.20\) 元不是三次预报的平均贡献，由置换检验，\(18{:}00\) 预报在其控制的槽上
几乎恒为 0、置换它不改变任何决策，而上午正是 \(0{:}00\) 预报系统性偏离最大（\chapref{chap:data}）
与紧急购电最集中的时段。

\textbf{成本侧与结论。} 把收益上界 \(58\,062.56\) 元摊到 334 天，等价于每天新增预报投入不超过
\(173.84\) 元。故本文的回答是：\textbf{四次预报已经足够，在预报获取成本可以忽略之前不建议继续加密
预报时刻}；若要追加投入，应优先投向\textbf{提高现有四次预报的精度}——相同时刻数下剩余误差仍对应
\(58\,062.56\) 元（\(0.42\%\)），只能靠精度改进消除 \evid{C-Q3-01}。适用边界：① 只在本题惩罚
倍数（违约 \(0.5p_t\)、超额 \(1.5p_t\)、紧急 \(5p_t\)）下成立；② 只在附件 3 的预报质量下成立，
预报误差成倍减小会使边际收益更早饱和；③ \(n=8,24\) 为数据中不存在的配置，以其上界代替，
故“不建议加密”是对\textbf{收益侧上界}的判断。实现正确性自检：四组费用满足不等式链

\begin{equation}
  C_{\mathrm{PI}}\le C_{n=4}\le C_{n=1}:\qquad
  13\,874\,340.39\le13\,932\,402.95\le14\,264\,792.13\quad(\text{元});
  \label{eq:q3-chain}
\end{equation}

若数值违反该链，说明非预期性、信息集或结算口径实现有误。

\subsection{求解结果与口径依赖}

在同一场景集上求解式~\eqref{eq:q3}，四项费用为：计划相关 \(13\,216\,203.48\) 元、违约
\(210\,674.71\) 元、超额 \(262\,628.79\) 元、紧急购电 \(242\,895.98\) 元，合计
\(13\,932\,402.95\) 元（1\,393.24 万元），相对无储能基线 \(16\,407\,319.63\) 元节省 \(15.08\%\)；
紧急购电量 \(64\,522.09\)~kWh，出现在 302/334 天（\calref{C-09}）。

\textbf{调整机制省在哪里。} 与“只按 \(0{:}00\) 计划结算”的问题 2 结构（计划购电费
\(13\,637\,552.90\) 元 \(+\) 紧急购电费 \(627\,239.24\) 元）相比，调整机制把\textbf{紧急购电费压低
\(384\,343.26\) 元}，代价是计划相关费用净增 \(51\,954.08\) 元（计划项减少 \(421\,349.41\) 元、
违约与超额合计增加 \(473\,303.50\) 元），相抵净省 \(332\,389.18\) 元。\(0{:}00\) 计划购电量合计
\(22\,220\,333.20\)~kWh，调整后 \(21\,843\,259.36\)~kWh（净减 \(377\,073.84\)~kWh），其中调减
\(645\,614.42\)~kWh、调增 \(268\,540.58\)~kWh，二者之差恰为净减少量（口径自检）。\(0{:}00\) 计划量
在四种信息结构下\textbf{完全相同}，信息价值的全部来源都在\textbf{调整}环节。

\textbf{与问题 2 的对照与口径依赖。} 两者同为附件 1 固定电价、同一场景集、同一实际轨迹与同一费用
口径，唯一差别是三次滚动预报是否可用；费用差额即信息价值 \(332\,389.18\) 元（\(2.33\%\)），相对
无储能基线的节省率由 \(13.06\%\) 提高到 \(15.08\%\)。该降幅\textbf{依赖于 \calref{C-09} 的偏差量定价
口径}：在“计划购电量全额结算 \(+\) 差额惩罚”的读法下，计划项不再随调整量下降，净收益会减少
\(\sum_tp_tv_t=421\,349.41\) 元并转为负值（\chapref{chap:verify} 给出对照），故该答案与市场规则
（惩罚倍数与结算方式）绑定，不能脱离口径单独引用。按题面表 1／表 2／表 3 格式给出的指定日期结果见
\chapref{chap:results}表~\ref{tab:key4}，其中“调整购电量”一列填的是\textbf{最后一次覆盖该槽位的决策}
给出的购电量（表~\ref{tab:q3-epochs}），与 \(0{:}00\) 计划量之差即该槽的偏差量。

%!TEX root = ../main.tex
% =============================================================================
%  七、问题 4：波动电价下的风险决策模型      建议 3–4 页
%  评审在问：波动电价怎么处理的？
%  本节必须出现：① 电价过程的四项检验  ② 两层建模（分时骨架 + 随机残差）
%                ③ 联合情景（二维分层取代表历史日），必须保留 rho(p, L)
%                ④ 波动电价下的 LP 形式（双重不确定性）
%                ⑤ 风险—费用前沿 + 真实轨迹回测（口径 C-11）
%  数据来源：notes/数值台账.md、intermediate/paper_numbers.json（键 电价检验/Q4_2/Q4_3）、
%            intermediate/q4_price.json、intermediate/q4_cvar.json
%  口径编号：按 notes/写作交底.md §6.1 的统一分配，CVaR 用 C-11
%  图文件名：fig08_q4_price.pdf（电价骨架与残差自相关）、fig09_cvar_frontier.pdf（前沿与回测）
% =============================================================================

\section{问题 4：波动电价下的风险决策模型}

\subsection{电价过程的统计检验}

问题 4 中电价与光伏出力同时随机；若把附件 4 的电价序列直接当作"另一组已知价格"代入模型，等于假设
\(0{:}00\) 就知道全天的实际电价，与题意矛盾。故本文先把电价当作\textbf{随机过程}检验并建模
（表~\ref{tab:price-test}）：检验在 2025 年全年（\(52\,560\) 个观测）上进行，费用结果仍只在 334 天上报告。



四项检验的建模含义见表~\ref{tab:price-test} 末列。其中最需强调的是 T4：电价与光伏的\emph{池化}相关
\(+0.028\) 几乎为 0——这只是"夜间光伏恒为 0"稀释后的假象，日聚合尺度上的真实结构是 \(-0.343\)；若据此
独立抽样，就会生成"午间光伏大发、同时电价高企"这类不可能情景并\textbf{低估尾部风险}（图~\ref{fig:q4-price}）。

\begin{figure}[htbp]
  \centering
  \includefig[width=0.80\linewidth]{fig08_q4_price.pdf}
  \caption{电价过程的骨架与残差结构（上子图：附件 4 电价的逐槽分位带 10\%--90\% 与 25\%--75\%、
           逐槽中位数与附件 1 分时骨架；下子图：日内残差自相关 \(\mathrm{ACF}(1\text{--}6)\)，
           并标注日级残差自相关与峰谷波动比；范围：2025-01-01 至 12-31 共 365 天 \(\times\) 144 槽；
           单位：电价 元/kWh，自相关为无量纲量；样本：\(n=52\,560\) 个电价观测）}
  \label{fig:q4-price}
\end{figure}

\subsection{两层电价模型与联合情景生成}

本文用\textbf{两层结构}刻画电价：第一层是可预测的\textbf{分时骨架} \(\bar p_t\)（取附件 1 的电价剖面，
全年尺度上等于附件 4 的逐时段均值，最大偏差 \(5.15\times10^{-5}\)~元/kWh），第二层是不可预测的
\textbf{随机残差} \(\varepsilon_{d,t}\)：
\begin{equation}
  p_{d,t}=\bar p_{t}+\varepsilon_{d,t},\qquad
  \mathbb{E}[\varepsilon_{d,t}]=0
  \label{eq:q4-price}
\end{equation}
其中 \(d\) 为日期、\(t\) 为槽位，\(p_{d,t}\) 的单位为\unit{元/kWh}。按 \S7.1 的 T3，残差须
\textbf{按时刻对齐}重采样，故本文直接以\textbf{整日}为单位构造情景。

\textbf{联合情景：二维分层取代表历史日。} 设 \(r_{d,t}\) 为第 \(t\) 槽的净需求预报残差（\(0{:}00\) 预报
口径的净需求与实际净需求之差，单位\unit{kWh}），\(p_{d}\) 为该日平均电价。把 334 天按
\((\sum_t r_{d,t},\ p_d)\) 各分 8 层形成 \(8\times8\) 网格，每格取\textbf{medoid}（层内欧氏距离和最小的
历史日）为代表，权重取层内样本占比。所得 64 个代表日（权重 \(0.0030\)--\(0.0299\)，合计 1）保真两个
维度的边缘分布，并因电价剖面与净需求残差来自\textbf{同一个历史日}而\textbf{天然保留联合结构}。

情景保真度见表~\ref{tab:q4-scenfid}。作为对照，本文另用\textbf{分块 Bootstrap}（6~h 块、按时刻对齐）
生成 64 个合成情景：优势是不依赖历史日的可获取性，但本题样本量下边缘保真度明显更差（均值偏差 3.41\%、
1\% 分位相对差 20.4\%），故\textbf{只作为对照}，主口径采用代表历史日。



\textbf{必须保留 \(\rho(p,L)\) 的理由。} 电价与净负荷的正相关（池化 \(+0.260\)）意味着"贵的时候恰好也在
多买电"；若独立抽样，"高价 \(\times\) 高净需求"的联合尾部占比会被人为压低，风险度量随之失去意义。二维
网格取代表日后 \(\rho(p,\text{净负荷})=0.2597\)，与实际 0.2620 相差 0.0023，尾部结构被保真。

\textbf{波动电价下的 LP 形式与双重不确定性。} 设 \(\Omega\) 为联合情景集（权重 \(w_\omega\)），
\(r_{\omega,t}\) 为情景 \(\omega\) 的净需求残差，\(\widehat{PV}^{(0)}_t\) 为 \(0{:}00\) 发布的光伏预报，
\(p_{\omega,t}\) 为情景电价，则重算问题 2 的两阶段 LP 为
\begin{subequations}\label{eq:q4-lp}
  \begin{equation}
    \min\ \ \sum_{t=0}^{143}\bar p^{\text{sc}}_tE_t
    +\sum_{\omega\in\Omega}w_\omega\sum_{t=0}^{143}5\,p_{\omega,t}\,g_{\omega,t},
    \qquad \bar p^{\text{sc}}_t=\sum_{\omega\in\Omega}w_\omega p_{\omega,t}
    \label{eq:q4-lp-obj}
  \end{equation}
  \begin{equation}
    E_t-c_t+d_t\ \ge\ \bigl(L_t-\widehat{PV}^{(0)}_t\bigr)\Delta t
    \label{eq:q4-lp-plan}
  \end{equation}
  \begin{equation}
    g_{\omega,t}\ \ge\ \bigl(L_t-\widehat{PV}^{(0)}_t\bigr)\Delta t+r_{\omega,t}
    +c_t-d_t-E_t,\qquad g_{\omega,t}\ge 0
    \label{eq:q4-lp-gap}
  \end{equation}
\end{subequations}
约束外加式~\eqref{eq:q3-soc} 的储电量递推、界与终端条件（物理参数同 \S6.4）。两个不确定性分别进入
\textbf{目标系数}（电价 \(p_{\omega,t}\)）与\textbf{约束右端项}（净需求残差 \(r_{\omega,t}\)），即问题 4
相对问题 2 的\textbf{双重不确定性}。决策时的电价信息取情景概率加权均值 \(\bar p^{\text{sc}}_t\)：附件 4
未提供日前电价预报，故 \(0{:}00\) 可用的电价预报取\textbf{分时骨架} \(\bar p_t\)（假设 A10），
\(\bar p^{\text{sc}}_t\) 与骨架的最大偏差 0.0462~元/kWh（占均价 6.02\%）。\textbf{结算一律按附件 4 的当日
实际电价}，使"决策用预报、结算用实现"各归其位（\calref{C-03}）。问题 4-3 在式~\eqref{eq:q4-lp} 上再叠加
\(6{:}00/12{:}00/18{:}00\) 三次滚动修正，调整与违约费用按 \calref{C-09} 的偏差量定价结算。

\subsection{阈值律在波动电价下的形式不变性}

问题 1 的阈值律只依赖①目标对购电量线性、②储电量递推两条结构性质，\textbf{完全不依赖电价的分布假设}，
故电价由确定变随机只改变阈值（水价值 \(\nu_t\)）的\textbf{取值}，不改变"电价低于充电阈值则充、高于放电
阈值则放"的\textbf{函数形式}：\(\alpha_t\le\eta_c\nu_t\) 时充电、\(\alpha_t\ge\nu_t/\eta_d\) 时放电、
\(\eta_c\nu_t<\alpha_t<\nu_t/\eta_d\) 时既不充也不放（死区，两侧之比 \(1/(\eta_c\eta_d)=1.2346\)）。
两点边界：①该结论\textbf{在储电量界与充放功率界不活跃时精确成立}，界活跃时须回到互补松弛式~
\eqref{eq:q1-kkt-c}--\eqref{eq:q1-kkt-d} 并计入界乘子；②"形式不变"不等于"阈值是常数"——波动电价下
\(\nu_t\) 仍须\textbf{逐时段由场景对偶确定}，随情景集与风险参数 \(\lambda\) 变化（\evid{C-Q4-01}）。

\subsection{风险—费用有效前沿}

\textbf{CVaR 的线性化。} 记 \(C_\omega\) 为情景 \(\omega\) 的总购电费用、\(\alpha=0.95\) 为锁定置信水平，
条件风险价值（Rockafellar--Uryasev 形式）为
\begin{equation}
  \mathrm{CVaR}_\alpha(C)=\min_{\eta}\Bigl\{\eta+\frac{1}{(1-\alpha)|\Omega|}
  \sum_{\omega\in\Omega}\max\bigl(0,\,C_\omega-\eta\bigr)\Bigr\},
  \label{eq:q4-cvar}
\end{equation}
引入辅助变量 \(z_\omega\ge C_\omega-\eta\)、\(z_\omega\ge0\) 后，式~\eqref{eq:q4-cvar} 的分段线性结构
完全由线性约束表达，\textbf{加入风险度量后模型仍是 LP}，不引入整数变量。

\begin{caliber}{C-11}{条件风险价值 CVaR（\(\alpha=0.95\)）}
\textbf{定义}\quad 式~\eqref{eq:q4-cvar} 的 \(\mathrm{CVaR}_{0.95}\)，
即最差 5\% 情景下费用的条件期望，按 Rockafellar--Uryasev 线性化求解；
目标函数取凸组合 \(\min\ (1-\lambda)\,\mathbb{E}[C]+\lambda\,\mathrm{CVaR}_{0.95}\)，
\(\lambda\in[0,1]\)（\(\lambda=0\) 为纯期望最优、\(\lambda=1\) 为纯风险最优，
与"\(\min\mathbb{E}[C]+\lambda'\mathrm{CVaR}\)"式在 \(\lambda'=\lambda/(1-\lambda)\) 下等价）；
除假设 A10 外，模型、情景集与物理参数在扫描中保持不变；%
\textbf{数据}\quad 联合情景集（\(8\times8\) 二维分层得到的 64 个代表历史日，权重为层占比）；
评估用 2025-02-01 至 12-31 中\textbf{每 6 天取 1 天}的 56 天子样本，逐日独立求解；%
\textbf{时段}\quad 上述 56 个代表日，单日 144 槽；\textbf{单位}\quad 元/日（日费用的样本均值）、\unit{kWh}/日（电量）；%
\textbf{不含}\quad 不含储能投资与运维成本；也不含"最坏情形"（\(\min\max_\omega C_\omega\)）口径，
故 \(\lambda=1\) 只是\textbf{近似最坏}；%
\textbf{对照}\quad 每一点同时给出四种量：情景内期望费用、情景内 \(\mathrm{CVaR}_{0.95}\)、
\textbf{真实轨迹回测}（在同样的代表日上用附件 4 实际电价与附件 2 实际负载、光伏重算的日费用）与紧急购电量；%
\textbf{复算}\quad \texttt{q4.py --mode cvar --K 64 --step 6}，输出 \texttt{intermediate/q4\_cvar.json} 的 frontier 列表。
\end{caliber}

以 \(\lambda\)（\(0,0.1,\dots,1.0\) 共 11 档）为参数扫描得到\textbf{风险—费用有效前沿}（表~
\ref{tab:frontier-backtest}、图~\ref{fig:frontier}），两端极端解为 \(\lambda=0\) 的纯期望最优与
\(\lambda=1\) 的纯风险最优，后者以 0.99\% 的期望费用换取 3.19\% 的 \(\mathrm{CVaR}\) 下降。



\textbf{推荐点的判据事先锁定}：在"情景内期望费用相对 \(\lambda=0\) 上升不超过 0.5\%"（门限
42\,205.57~元/日）下使 \(\mathrm{CVaR}_{0.95}\) 最小。满足该门限的最大档位是 \(\lambda^{*}=0.3\)：
期望费用 42\,168.43~元/日（\(+0.41\%\)）、\(\mathrm{CVaR}_{0.95}\) 50\,814.83~元/日（\(-2.61\%\)）、
\textbf{真实轨迹回测费用 45\,778.03~元/日（仅 \(+0.18\%\)）}，紧急购电量由 480.19 降到
324.60~\unit{kWh}/日（\(-32.40\%\)）。

\textbf{前沿的拐点。} \(\lambda\) 从 0.7 升到 0.9 时 \(\mathrm{CVaR}_{0.95}\) 只再降 11.53~元/日（0.02\%），
期望费用却继续上升 34.27~元/日（0.08\%）；\(\lambda=1\) 的 \(\mathrm{CVaR}_{0.95}\) 与 \(\lambda=0.9\)
基本持平（50\,509.78 对 50\,510.46~元/日，差 0.68~元/日，属噪声量级）。即\textbf{\(\lambda\approx0.7\)
以后风险已基本降不动}。真实轨迹回测：三个代表点在\textbf{实际发生}的 334 天轨迹口径下日费用为
45\,694.27--45\,952.85~元，情景内的排序在真实轨迹上同样成立但幅度小得多（\(+0.57\%\) vs \(+0.99\%\)）。

\begin{figure}[htbp]
  \centering
  \includefig[width=0.80\linewidth]{fig09_cvar_frontier.pdf}
  \caption{风险—费用有效前沿与真实轨迹回测（左子图：情景期望费用、\(\mathrm{CVaR}_{0.95}\)
           与真实轨迹回测费用随风险厌恶系数 \(\lambda\) 的变化；右子图：期望费用—\(\mathrm{CVaR}_{0.95}\)
           前沿散点，颜色表示 \(\lambda\)，星号标出边际风险下降最快的档位；
           对象：\(\lambda\in\{0,0.1,\dots,1.0\}\) 共 11 档的解；范围：2025-02-01 至 12-31 中每 6 天取 1 天的 56 个代表日；
           单位：费用 元/日；样本：64 个代表情景、11 个 \(\lambda\) 档位，\(\alpha=0.95\)）}
  \label{fig:frontier}
\end{figure}

\subsection{求解结果}

在波动电价下重解问题 2 与问题 3（表~\ref{tab:q4-summary}），两问的费用口径与固定电价下\textbf{完全相同}
（\calref{C-03}、\calref{C-09}），差别只在电价序列取附件 4 实际值、决策电价信息取情景加权均值，故两组
结果可直接比较。



\textbf{电价波动把费用抬高了多少。} 同一 334 天窗口、同一实际负载与光伏轨迹下，波动价下两阶段与四次
滚动结构的费用分别高出 5.81\% 与 5.67\%，无储能基线高出 3.78\%，节省率因此由 13.06\%/15.08\% 降到
11.35\%/13.54\%（分别下降 1.71 与 1.55 个百分点；表~\ref{tab:q4-summary}）。波动的成本上升\textbf{不能由
价格水平解释}：窗口内附件 4 电价均值 0.7575~元/kWh 反而比方时骨架的 0.7662~元/kWh 低 1.13\%，基线费用仍
高 3.78\%——差额来自\textbf{高电价与高净负荷同时出现}（池化相关 \(+0.260\)，见 \S7.2），这正是"必须生成
联合情景"的定量证据。

\textbf{波动电价下的滚动调整依然有效。} 四次预报把波动价下的费用降低 372\,102.66~元（2.47\%），
紧急购电量由 182\,759.86~\unit{kWh} 降到 63\,523.11~\unit{kWh}（降 65.24\%），降幅大于固定价下的
2.33\%——即\textbf{电价越不确定，滚动更新的价值越大}。问题 4-3 的四项费用为：计划相关
13\,979\,692.49~元、违约 224\,406.57~元、超额 270\,998.49~元、紧急购电 247\,055.63~元，合计
14\,722\,153.17~元（1\,472.22~万元），其中违约电量 664\,243.01~\unit{kWh}、超额电量
259\,710.27~\unit{kWh}；违约与超额费用同时高于固定价，说明\textbf{价格不确定使"调整到多少"本身也变得
更难判断}。四次滚动下波动价与固定价的紧急购电量几乎相同（反而少 998.98~\unit{kWh}，\(-1.55\%\)），
而两阶段结构下波动价多买 10\,744.04~\unit{kWh}（\(+6.25\%\)），即\textbf{滚动修正几乎完全吸收了电价波动
带来的额外缺口}。四个指定日期的结果见第八章表~\ref{tab:key4} 与表~\ref{tab:emergency}（逐时段完整结果见
\texttt{result4-2.xlsx}、\texttt{result4-3.xlsx}）。

%!TEX root = ../main.tex
% =============================================================================
%  八、结果与分析                建议 4–5 页
%  评审在问：数对不对？有没有解释？
%  本节安排：① 口径框 C-02、C-03、C-04、C-06、C-07、C-08（编号按交底书 §6.1 唯一分配；
%               C-01、C-05、C-09、C-10、C-11 分别首次出现于第四、五、六、七章）；
%            ② 四问结果总表（含无储能基线与节省率）；
%            ③ 成本分解（fig05 为问题 1 日费用瀑布）与策略机制解释（fig03、fig10）；
%            ④ 题面表 1、表 2 格式：问题 2（主模型）在四个指定日期上完整给出；
%            ⑤ 题面表 3 格式：四种方案 × 四个指定日期合并为一张规整表；
%            ⑥ 问题 3、4-2、4-3 的四个指定日期汇总表 + 结果文件指引（附录 C 有对照表）。
%  红线：本节每一个数值都必须与 result*.xlsx 逐格一致（见第九章审计 A9、A10）。
% =============================================================================

\section{结果与分析}

本节给出四问的数值结果与解释；数字取自 \texttt{intermediate/paper\_numbers.json}
（\texttt{code/paper\_numbers.py} 汇总），口径由 \texttt{code/config.py} 唯一确定，
可由 \texttt{code/audit.py} 独立复算（见第九章）。

\subsection{指标口径表}

本节把在本题参数范围内需要逐项说明的六个核心指标口径一次列清，其余引用其定义章：
无储能基线 \calref{C-01}（第四章，V1 固定价窗口／V2 波动价窗口／V3 典型日）、
紧急购电量 \calref{C-05}（第五章）、调整与违约费用 \calref{C-09}、信息价值 \calref{C-10}（第六章）、
条件风险价值 \calref{C-11}（第七章）。

\begin{caliber}{C-02}{全天购电量}
\textbf{定义}\quad \(E^{\text{day}}=\sum_{t=1}^{144}E_t\)，即该日 144 个 10~min 槽的\textbf{计划购电量}之和；%
\textbf{公式}\quad 结果为逐槽求和，无闭式；结果文件每行末列即该日汇总；%
\textbf{数据}\quad 结果文件「计划购电量」工作表的 144 个时段列；%
\textbf{时段}\quad 单日；问题 1 为典型日，问题 2--4 为 2025-02-01 至 12-31 的 334 天；%
\textbf{单位}\quad kWh（结果文件保留 4 位小数，正文保留 2--4 位）；%
\textbf{含}\quad 问题 2、4-2 中即第一阶段计划量；问题 3、4-3 中调整后的量一律称"全天调整购电量"并单独命名；%
\textbf{不含}\quad 不含储能充电量、不含紧急购电量 \calref{C-05}、不含盈余处置量 \calref{C-06}；%
\textbf{对照}\quad 问题 2、3 的计划购电量模型内值为 \(22\,220\,333.2033\)~kWh，结果文件回读值为 \(22\,220\,333.1391\)~kWh，
相对差 \(2.9\times10^{-9}\)，来自写入时的 4 位小数取整；%
\textbf{复算}\quad \texttt{results\_io.py::plan\_sheet\_frame()}（写出行汇总列）与 \texttt{audit.py::audit\_result\_files()}（回读结果文件后重算 \(\sum_{144}\) 槽）
\end{caliber}

\begin{caliber}{C-03}{总购电费用（计划 \(+\) 紧急 \(+\) 调整相关）}
\textbf{定义}\quad \(C^{\text{total}}=C^{\text{plan}}+C^{\text{breach}}+C^{\text{excess}}+C^{\text{emg}}\)；
问题 2、4-2 无调整机制，后两项恒为 0；%
\textbf{公式}\quad 问题 2、4-2：\(C=\sum_tp_tE_t+5\sum_tp_tg_t\)；
问题 3、4-3：\(C^{\text{plan}},C^{\text{breach}},C^{\text{excess}}\) 按 \calref{C-09} 的偏差量定价口径计算，
\(C^{\text{emg}}=\sum_t5p_tg_t\)；%
\textbf{数据}\quad 电价取附件 1（问题 1--3）或附件 4（问题 4）；电量取本队结果文件；%
\textbf{时段}\quad 问题 1 为典型日；问题 2--4 为 334 天；%
\textbf{单位}\quad 元（保留 2 位小数，内部计算不取整）；%
\textbf{不含}\quad 不含储能投资与运维成本、容量补偿、辅助服务收益、绿电环境价值、循环衰减与安全成本；%
\textbf{对照}\quad 若只报计划购电费而不计紧急购电，问题 2 为 \(13\,637\,552.90\)~元，
比 \(C^{\text{total}}\) 低 \(627\,239.24\)~元（\(4.40\%\)）；%
\textbf{复算}\quad \texttt{q2.py::run\_q2()}、\texttt{q3.py::summary()}、\texttt{q4.py::run\_q4\_2()} 与 \texttt{audit.py::audit\_result\_files()}，
四项之和与总费用的偏差为 0
\end{caliber}

\begin{caliber}{C-04}{节省率}
\textbf{定义}\quad \(\eta=\bigl(C^{base}-C^{\text{total}}\bigr)/C^{base}\times100\%\)，分母与分子必须同窗口、同电价口径；%
\textbf{公式}\quad \(C^{base}\) 取 \calref{C-01} 的对应变体（V1 固定价窗口 / V2 波动价窗口 / V3 典型日）；%
\textbf{数据}\quad \(C^{base}\) 与 \(C^{\text{total}}\) 均取自本队结果；%
\textbf{时段}\quad 与 \calref{C-01} 变体一致；%
\textbf{单位}\quad \%（保留 2 位小数）；%
\textbf{不含}\quad 不得相对贪心策略或完美信息下界报同一个数而不加限定；相对其他基线的结果必须\textbf{单独命名}；%
\textbf{对照}\quad 问题 1 相对 V3 为 \(26.90\%\)；问题 2、3 相对 V1 为 \(13.06\%\)、\(15.08\%\)；
问题 4-2、4-3 相对 V2 为 \(11.35\%\)、\(13.54\%\)；%
\textbf{复算}\quad \texttt{paper\_numbers.py::main()}（由 \calref{C-01} 与 \calref{C-03} 导出）
\end{caliber}

\begin{caliber}{C-06}{弃光/盈余处置电量}
\textbf{定义}\quad 光伏盈余与超购部分的处置量：典型日按
\(s_t=PV_t\Delta t+E_t+d_t-L_t\Delta t-c_t\ge0\) 逐槽统计（式~\eqref{eq:q1-balance} 的松弛量）；
窗口口径按实际数据复算为 \(\sum_t\max\{PV_t\Delta t+E_t+d_t+g_t-L_t\Delta t-c_t,0\}\)；%
\textbf{公式}\quad 上式为定义式；该量非负，其余额为超购部分被弃置的电量；%
\textbf{数据}\quad \(L,PV\) 取附件 2 实际值，其余取本队结果文件；%
\textbf{时段}\quad 逐槽统计、按窗口（334 天）汇总；%
\textbf{单位}\quad kWh（按电量统计，\textbf{不得}按功率报）；%
\textbf{不含}\quad 不含"因充电功率受限而未发生的充电"，不含未购电部分；%
\textbf{对照}\quad 问题 1 典型日最优弃光 \(=0.0000\)~kWh（数值残差 \(1.37\times10^{-12}\)~kWh）；
问题 2、3 的窗口盈余处置量为 \(3\,106\,390.84\)~kWh 与 \(2\,638\,963.35\)~kWh；
问题 4-2、4-3 的同一量尚未纳入汇总；%
\textbf{复算}\quad \texttt{config.py::curtailment(price, load, pv, E, c, d)}（典型日）与 \texttt{audit.py::audit\_internal()}（窗口口径）
\end{caliber}

\begin{caliber}{C-07}{充放电量与储电量（交流侧）}
\textbf{定义}\quad 结果文件中的 \(c_t,d_t\) 为\textbf{并网点（AC 侧）}电量；储电量为\textbf{时段末}值 \(S_t\)；%
\textbf{公式}\quad 递推见式~\eqref{eq:q1-soc}，\(\eta_c=\eta_d=0.9\)（单边口径，往返 \(0.81\)）；
充放功率上限 \(5000\)~kW 对应单槽上限 \(\overline P\Delta t=833.3333\)~kWh；%
\textbf{数据}\quad 本队结果文件「充放电量」工作表；%
\textbf{时段}\quad 逐槽；表 2 格式按 6 个 4~h 块（每块 24 槽）汇总，并给出 \(0{:}00\) 与 \(24{:}00\) 两个储电量点；%
\textbf{单位}\quad kWh；运行区间 \([1200,10800]\)~kWh；%
\textbf{不含}\quad 不含电池侧（DC）电量、不含自放电与循环衰减（假设 A4）；%
\textbf{对照}\quad 恒等式校验：问题 1 的充放比 \(20\,740.67/16\,799.94=1.2346=1/(\eta_c\eta_d)\)，
与 \(\eta_c\eta_d=0.81\) 互为倒数，偏差 \(4.44\times10^{-16}\)；问题 2 四个指定日期的合计充放比同为 \(1.2346\)；%
\textbf{复算}\quad \texttt{results\_io.py::soc\_sheet\_frame()}（块级汇总）与 \texttt{audit.py::audit\_result\_files()}（块级净递推）
\end{caliber}

\begin{caliber}{C-08}{储电量递推与终端条件}
\textbf{定义}\quad 储电量\textbf{跨日连续递推}（\(S_{d,0}=S_{d-1,144}\)，假设 A5），
窗口末端施加终端条件 \(S_T\ge S_0\)（假设 A6），问题 1 另加严格日周期 \(S_0=S_{144}\)；%
\textbf{公式}\quad \(S_t=S_{t-1}+\eta_cc_t-d_t/\eta_d\)，\(S_t\in[1200,10800]\)~kWh；
2025 年 1 月为预热期（假设 A8），其终值即 2 月 1 日初值；%
\textbf{数据}\quad 结果文件「充放电量」工作表的储电量列与中间结果 \texttt{q2\_main.json} 等；%
\textbf{时段}\quad 逐槽；窗口 2025-02-01 至 12-31（334 天）；%
\textbf{单位}\quad kWh；%
\textbf{不含}\quad 不含"每日清零""每日重置为 6000"两种读法——前者会引入末日掏空型短视解，后者与 1 月预热期矛盾；%
\textbf{对照}\quad 问题 1：日周期口径 \(35\,126.9486\)~元（\(S_0=S_{144}=6000.00\)~kWh）；
仅要求两端相等的自由环状口径 \(35\,118.5986\)~元（最优 \(S=8550.00\)~kWh），相差 \(8.35\)~元（\(0.02\%\)）；
窗口口径下"末值低于初值的天数"为 0；%
\textbf{复算}\quad \texttt{lp\_engine.py::solve\_day()}（终端约束）与 \texttt{audit.py::audit\_result\_files()}（\texttt{A6\_末值低于初值天数}）
\end{caliber}

\subsection{四问结果的汇总与对照}

\begin{table}[htbp]
  \centering
  \caption{四问主要结果汇总（费用单位：元，电量单位：kWh；问题 1 为典型日，
           其余为 2025-02-01 至 2025-12-31 共 334 天；问题 4 的电价为附件 4 波动价；
           表中带 a 标记者见表注）}
  \label{tab:summary}
  \scriptsize
  \setlength{\tabcolsep}{3pt}
  \begin{tabularx}{\linewidth}{L{2.35cm} C{1.60cm} C{1.95cm} C{1.95cm} C{1.95cm} Y}
    \toprule
    指标 & 问题 1 & 问题 2 & 问题 3 & 问题 4-2 & 问题 4-3 \\
    \midrule
    计划购电量 & 59\,482.70 & 22\,220\,333.14 & 22\,220\,333.14 & 22\,198\,236.57 & 22\,255\,478.92 \\
    调整购电量 & --- & --- & 21\,843\,259.36 & --- & 21\,850\,946.24 \\
    充电量\textsuperscript{a} & 20\,740.67 & 89\,108.12 & 86\,965.89 & 88\,470.42 & 86\,423.35 \\
    放电量\textsuperscript{a} & 16\,799.94 & 72\,177.58 & 70\,442.37 & 71\,661.04 & 70\,002.91 \\
    紧急购电量 \calref{C-05} & --- & 172\,015.82 & 64\,522.09 & 182\,759.86 & 63\,523.11 \\
    盈余处置量 \calref{C-06} & 0.0000 & 3\,106\,390.84 & 2\,638\,963.35 & 3\,100\,463.79 & 2\,645\,427.81 \\
    计划购电费 & 35\,126.95 & 13\,637\,552.90 & 13\,216\,203.48 & 14\,397\,542.37 & 13\,979\,692.49 \\
    违约相关费用 \calref{C-09} & --- & --- & 210\,674.71 & --- & 224\,406.57 \\
    超额相关费用 \calref{C-09} & --- & --- & 262\,628.79 & --- & 270\,998.49 \\
    紧急购电费 & --- & 627\,239.24 & 242\,895.98 & 696\,713.47 & 247\,055.63 \\
    \textbf{总费用} & \textbf{35\,126.95} & \textbf{14\,264\,792.13} & \textbf{13\,932\,402.95} & \textbf{15\,094\,255.84} & \textbf{14\,722\,153.17} \\
    无储能基线 \(C^{base}\)（\calref{C-01}） & 48\,052.05（V3） & 16\,407\,319.63（V1） & 16\,407\,319.63（V1） & 17\,026\,782.86（V2） & 17\,026\,782.86（V2） \\
    相对 \(C^{base}\) 节省率（\calref{C-04}） & 26.90\% & 13.06\% & 15.08\% & 11.35\% & 13.54\% \\
    \bottomrule
  \end{tabularx}
  \normalsize
\end{table}

注 a：问题 1 为典型日全天值；问题 2--4-3 为题面指定四个日期（2025-03-20、06-21、
09-23、12-21）的合计值，逐日值见表~\ref{tab:q2-t2}、表~\ref{tab:key4}；334 天全窗口合计未纳入台账，
可由 \texttt{result*.xlsx} 的「充放电量」工作表逐日求和复算。

表中问题 4 的窗口级数值按 334 天求和：问题 4-3 的全天调整购电量、问题 4-2 与 4-3 的
窗口盈余处置量（\calref{C-06}），可由 \texttt{result4-2.xlsx}、\texttt{result4-3.xlsx} 与附件 2
按 \calref{C-02}、\calref{C-06} 的定义复算后补入。

表~\ref{tab:summary} 的验收：四项费用之和与总费用偏差为 0（\calref{C-03}），每列均可追溯到对应的
\texttt{result*.xlsx}。紧急购电量 \calref{C-05} 分别在问题 2、3、4-2 的 296、302、294 天触发；
问题 2 的 \(172\,015.82\)~kWh 占该窗口计划购电量的 \(0.77\%\)，对应紧急购电费占总费用 \(4.40\%\)。
固定价口径下问题 2 日均节省 \(6\,414.75\)~元（问题 1 典型日 \(12\,925.10\)~元），差额主要来自
预报误差的紧急购电费（日均 \(1\,877.96\)~元），问题 3 四次滚动后升至 \(7\,409.93\)~元；
波动价口径下同一套模型的费用整体上移，但"四次滚动优于一次计划"的方向不变。\evid{C-Q1-03}\evid{C-Q3-01}

\subsection{成本分解与基线对照}

图~\ref{fig:costwaterfall} 给出问题 1 的日费用分解：底柱为无储能基线（\calref{C-01} 变体 V3）、
中间柱为储能削峰填谷收益、顶柱为最优购电费，三者同口径、同单位（元/日）。

\begin{figure}[htbp]
  \centering
  \includefig[width=0.80\linewidth]{fig05_cost_waterfall.pdf}
  \caption{问题 1 日费用分解瀑布图（对象：微网典型日购电费用；范围：附件 1 全年平均日、
           单日 144 个 10~min 槽；基线：无储能基线 \(C^{base}=48\,052.05\)~元/日（\calref{C-01} 变体 V3）；
           中间柱：储能削峰填谷收益 \(12\,925.10\)~元/日；终点：最优购电费 \(35\,126.95\)~元/日；
           单位：元/日）}
  \label{fig:costwaterfall}
\end{figure}

334 天窗口上的同一套分解由表~\ref{tab:summary} 承载：问题 2 的计划购电费与紧急购电费分别占总费用的
\(95.60\%\) 与 \(4.40\%\)——储能的价值在压低计划购电费，紧急购电费是预报误差的代价，两者不可相抵
后只报净数。四次滚动（\calref{C-09}）把四项占比改为 \(94.86\%\)、\(1.51\%\)、\(1.89\%\)、\(1.74\%\)；
波动价口径下问题 4-2、4-3 相对基线分别节省 \(1\,932\,527.02\)~元与 \(2\,304\,629.69\)~元，
两套基线不得混用。

\textbf{策略机制。}问题 1 的最优解（图~\ref{fig:mechanism}）中，充电量 \(20\,740.67\)~kWh 落于低谷
\(0{:}00\)--\(4{:}00\)（\(4\,500.00\)~kWh）、\(20{:}00\)--\(24{:}00\)（\(5\,333.33\)~kWh）与午间
\(12{:}00\)--\(16{:}00\)（\(5\,286.04\)~kWh），放电量 \(16\,799.94\)~kWh 集中于 \(4{:}00\)--\(8{:}00\)
（\(6\,365.84\)~kWh）与 \(16{:}00\)--\(20{:}00\)（\(5\,780.13\)~kWh）；两段低谷占全日充电量的
\(47.4\%\)、午间盈余段仅占 \(25.5\%\)，放电在两段高峰合计占 \(72.3\%\)——故主机制是
\textbf{面向电价的跨时段搬运}，而非"把午间盈余存到傍晚"。\evid{C-Q1-02}

对偶量给出与求解器无关的策略形式（图~\ref{fig:dualprice}）：最优性条件为
"\(\alpha_t\le\eta_c\nu_t\) 时充电、\(\alpha_t\ge\nu_t/\eta_d\) 时放电"（\(\nu_t,\alpha_t\) 分别为储电水价值
与电量平衡影子价格），实测 \(\nu_t\in[0.4711,1.1305]\)、\(\alpha_t\in[0.3713,1.2561]\)~元/kWh，
\(E_t>0\) 时 \(\alpha_t=p_t\)（偏差 \(0.000\)）；144 个槽（充电 45、放电 37）的阈值判定
\textbf{违反次数为 0}，互补松弛最大残差 \(0.000\)，死区比 \(1/(\eta_c\eta_d)=1.2346\) 与实测充放比一致。
\evid{C-Q1-02}

问题 1 的储电量日内触及下限 \(1200.00\)~kWh 与上限 \(10\,800.00\)~kWh，\(0{:}00\) 与 \(24{:}00\)
均为 \(6000.00\)~kWh。以典型日为基准算例（容量 \(12\,000\)~kWh、功率 \(5000\)~kW、费用 \(35\,126.95\)~元），
仅解除容量上限费用降至 \(27\,351.55\)~元（再省 \(22.14\%\)），仅解除充放功率上限降至 \(34\,544.43\)~元
（再省 \(1.66\%\)）：\textbf{容量摆幅是主要瓶颈，充放电功率在部分时段同时构成约束，两者量级不同}。

问题 2 在 \(2025\text{-}03\text{-}20\) 的解复现同一结构：当日充电 \(24\,072.26\)~kWh、放电
\(19\,498.53\)~kWh，比值同为 \(1.2346\)，逐块充放电量与储电量见表~\ref{tab:q2-t2}（其中
\(8{:}00\)--\(16{:}00\) 两段充电合计 \(13\,405.59\)~kWh）；被 \(5\) 倍紧急购电覆盖的缺口集中在上午
（表~\ref{tab:emergency} 的 \(6{:}00\)--\(12{:}00\) 段），四种方案下规律一致，故安全边际与预报纠偏
应加在上午时段，而不是全天均匀加码。

\subsection{题面指定时段与日期的结果表}

按题面要求，问题 2、3、4 需对 \(2025\text{-}03\text{-}20\)、\(06\text{-}21\)、\(09\text{-}23\)、\(12\text{-}21\)
四个日期各给出表 1 格式（指定 6 个 10~min 时段购电量、全天购电量与全天购电费）、表 2 格式
（6 个 4~h 块充放电量与 \(0{:}00\)、\(24{:}00\) 储电量）与表 3 格式（紧急购电时间段与购电量）。
本文以\textbf{问题 2 为主模型}，其表 1、表 2 格式按四个日期完整给出（表~\ref{tab:q2-t1}、
表~\ref{tab:q2-t2}）；问题 3 与 4-2、4-3 不逐日铺开，由表~\ref{tab:key4} 汇总承载（逐槽数值见
\texttt{result3.xlsx}、\texttt{result4-2.xlsx}、\texttt{result4-3.xlsx}，对应关系见附录 C），
紧急购电结果按题面表 3 格式合并于表~\ref{tab:emergency}。

\begin{table}[htbp]
  \centering
  \caption{问题 2 在四个指定日期的购电量（表 1 格式；单位：购电量 kWh、购电费 元；
           时段为 10~min，按附录 5 模板的位置读法，\(10{:}00\text{-}10{:}10\) 对应附件标签 \(10{:}10\)；
           全天购电量与全天购电费为对应日期 144 个槽的合计）}
  \label{tab:q2-t1}
  \footnotesize
  \setlength{\tabcolsep}{5pt}
  \begin{tabular}{lcccc}
    \toprule
    时间段 & 2025-03-20 & 2025-06-21 & 2025-09-23 & 2025-12-21 \\
    \midrule
    \(10{:}00\)--\(10{:}10\) & 0.0000 & 0.0000 & 0.0000 & 0.0000 \\
    \(12{:}00\)--\(12{:}10\) & 804.3029 & 0.0000 & 536.2415 & 997.0335 \\
    \(14{:}00\)--\(14{:}10\) & 0.0000 & 0.0000 & 0.0000 & 0.0000 \\
    \(16{:}00\)--\(16{:}10\) & 817.3061 & 255.2254 & 902.7969 & 1\,001.2240 \\
    \(18{:}00\)--\(18{:}10\) & 697.9973 & 0.0000 & 764.3298 & 715.4144 \\
    \(20{:}00\)--\(20{:}10\) & 0.0000 & 0.0000 & 0.0000 & 0.0000 \\
    \midrule
    全天购电量（kWh） & 74\,178.1298 & 35\,701.5391 & 74\,744.7963 & 99\,692.6601 \\
    全天购电费（元） & 45\,229.90 & 20\,324.26 & 47\,063.39 & 64\,376.98 \\
    \bottomrule
  \end{tabular}
\end{table}

\begin{table}[htbp]
  \centering
  \caption{问题 2 在四个指定日期的储能充放电量与储电量（表 2 格式；单位：kWh；
           6 个 4~h 块各含 24 个 10~min 槽；\(0{:}00\) 与 \(24{:}00\) 储电量对应题面表 2 的"时刻/储电量"列，
           四个日期均为 \(6000.00\)~kWh）}
  \label{tab:q2-t2}
  \footnotesize
  \setlength{\tabcolsep}{4pt}
  \begin{tabular}{l*{4}{cc}}
    \toprule
    \multirow{2}{*}{时间段} & \multicolumn{2}{c}{2025-03-20} & \multicolumn{2}{c}{2025-06-21}
                            & \multicolumn{2}{c}{2025-09-23} & \multicolumn{2}{c}{2025-12-21} \\
    \cmidrule(lr){2-3}\cmidrule(lr){4-5}\cmidrule(lr){6-7}\cmidrule(lr){8-9}
     & 充电量 & 放电量 & 充电量 & 放电量 & 充电量 & 放电量 & 充电量 & 放电量 \\
    \midrule
    \(0{:}00\)--\(4{:}00\)   & 4\,500.0000 & 0.0000 & 0.0000 & 0.0000 & 4\,500.0000 & 0.0000 & 4\,500.0000 & 0.0000 \\
    \(4{:}00\)--\(8{:}00\)   & 833.3333 & 6\,224.1108 & 140.7851 & 4\,303.4721 & 833.3333 & 5\,547.8264 & 1\,666.6667 & 1\,508.3333 \\
    \(8{:}00\)--\(12{:}00\)  & 6\,500.0000 & 2\,415.8892 & 5\,207.8672 & 130.5638 & 5\,833.3333 & 3\,020.9227 & 5\,833.3333 & 7\,806.6667 \\
    \(12{:}00\)--\(16{:}00\) & 6\,905.5931 & 2\,218.5304 & 5\,458.7995 & 0.0000 & 6\,728.4081 & 1\,606.2615 & 8\,333.3333 & 2\,835.0000 \\
    \(16{:}00\)--\(20{:}00\) & 0.0000 & 5\,709.7835 & 0.0000 & 6\,154.1006 & 0.0000 & 5\,602.9091 & 0.0000 & 5\,797.4742 \\
    \(20{:}00\)--\(24{:}00\) & 5\,333.3333 & 2\,930.2165 & 5\,333.3333 & 2\,485.8994 & 5\,333.3333 & 3\,037.0909 & 5\,333.3333 & 2\,842.5258 \\
    \midrule
    \(0{:}00\) 储电量 & \multicolumn{2}{c}{6\,000.00} & \multicolumn{2}{c}{6\,000.00}
                      & \multicolumn{2}{c}{6\,000.00} & \multicolumn{2}{c}{6\,000.00} \\
    \(24{:}00\) 储电量 & \multicolumn{2}{c}{6\,000.00} & \multicolumn{2}{c}{6\,000.00}
                       & \multicolumn{2}{c}{6\,000.00} & \multicolumn{2}{c}{6\,000.00} \\
    \bottomrule
  \end{tabular}
\end{table}

\begin{table}[htbp]
  \centering
  \caption{四个指定日期的紧急购电时间段与购电量（表 3 格式；单位：kWh；
           行 = 日期 \(\times\) 方案，括号内为该时段购电量，末列为该日合计；
           时段按连续触发的 10~min 槽合并、左闭右开；
           问题 2、4-2 无调整机制，问题 4-2 在 \(2025\text{-}12\text{-}21\) 未触发紧急购电）}
  \label{tab:emergency}
  \footnotesize
  \setlength{\tabcolsep}{4pt}
  \begin{tabular}{C{1.35cm} C{1.85cm} L{8.10cm} C{1.85cm}}
    \toprule
    方案 & 日期 & 紧急购电时间段（括号内购电量，kWh） & 合计（kWh） \\
    \midrule
    问题 2 & 2025-03-20 & \(6{:}00\)--\(8{:}00\)（228.99）；\(8{:}10\)--\(8{:}30\)（49.16）；\(8{:}50\)--\(10{:}30\)（531.27）；\(10{:}40\)--\(11{:}40\)（216.46） & 1\,025.88 \\
    问题 2 & 2025-06-21 & \(4{:}00\)--\(4{:}50\)（5.24）；\(14{:}00\)--\(14{:}30\)（181.84） & 187.08 \\
    问题 2 & 2025-09-23 & \(6{:}50\)--\(7{:}40\)（132.59）；\(11{:}00\)--\(11{:}40\)（60.77） & 193.37 \\
    问题 2 & 2025-12-21 & \(11{:}00\)--\(11{:}30\)（22.53） & 22.53 \\
    \midrule
    问题 3 & 2025-03-20 & \(6{:}00\)--\(7{:}00\)（138.11）；\(7{:}40\)--\(9{:}20\)（206.13）；\(9{:}30\)--\(10{:}00\)（88.68）；\(10{:}50\)--\(12{:}00\)（344.46）；\(12{:}30\)--\(12{:}50\)（102.31）；\(13{:}20\)--\(13{:}50\)（94.60） & 974.29 \\
    问题 3 & 2025-06-21 & \(4{:}00\)--\(4{:}50\)（5.24） & 5.24 \\
    问题 3 & 2025-09-23 & \(7{:}00\)--\(7{:}40\)（52.76）；\(8{:}30\)--\(8{:}40\)（4.26）；\(11{:}00\)--\(11{:}10\)（40.45） & 97.47 \\
    问题 3 & 2025-12-21 & \(9{:}50\)--\(10{:}00\)（2.12）；\(11{:}00\)--\(12{:}00\)（358.56） & 360.69 \\
    \midrule
    问题 4-2 & 2025-03-20 & \(6{:}00\)--\(8{:}00\)（225.30）；\(8{:}10\)--\(8{:}40\)（55.07）；\(8{:}50\)--\(10{:}30\)（523.99）；\(10{:}40\)--\(11{:}20\)（109.82） & 914.18 \\
    问题 4-2 & 2025-06-21 & \(4{:}00\)--\(4{:}50\)（5.24）；\(5{:}00\)--\(5{:}30\)（82.79）；\(14{:}00\)--\(14{:}30\)（164.79） & 252.82 \\
    问题 4-2 & 2025-09-23 & \(6{:}20\)--\(6{:}30\)（1.30）；\(6{:}50\)--\(7{:}40\)（125.27）；\(11{:}00\)--\(11{:}10\)（23.07） & 149.64 \\
    问题 4-2 & 2025-12-21 & 无 & 0.00 \\
    \midrule
    问题 4-3 & 2025-03-20 & \(6{:}00\)--\(7{:}00\)（125.36）；\(7{:}40\)--\(9{:}20\)（181.60）；\(9{:}30\)--\(10{:}00\)（87.91）；\(10{:}50\)--\(12{:}00\)（325.44）；\(12{:}30\)--\(12{:}50\)（86.88）；\(13{:}20\)--\(13{:}40\)（87.36） & 894.55 \\
    问题 4-3 & 2025-06-21 & \(4{:}00\)--\(4{:}50\)（5.24） & 5.24 \\
    问题 4-3 & 2025-09-23 & \(7{:}00\)--\(7{:}40\)（61.40）；\(11{:}00\)--\(11{:}10\)（30.51） & 91.92 \\
    问题 4-3 & 2025-12-21 & \(11{:}00\)--\(12{:}00\)（345.33） & 345.33 \\
    \bottomrule
  \end{tabular}
\end{table}

\begin{table}[htbp]
  \centering
  \caption{问题 3、4-2、4-3 在四个指定日期的汇总结果（单位：电量 kWh、费用 元；
           全部为当日 144 个槽的合计；"全天购电费"为 \calref{C-09} 偏差量定价口径下的 \(C^{\text{plan}}\)；
           该日总费用 \(=\) 全天购电费 \(+\) 违约与超额相关费用 \(+\) 紧急购电费）}
  \label{tab:key4}
  \footnotesize
  \setlength{\tabcolsep}{3.5pt}
  \begin{tabular}{C{1.35cm} C{1.75cm} C{2.10cm} C{2.10cm} C{1.85cm} C{1.75cm} C{1.75cm}}
    \toprule
    方案 & 日期 & 全天计划购电量 & 全天调整购电量 & 全天购电费 & 紧急购电量 & 该日总费用 \\
    \midrule
    问题 3 & 2025-03-20 & 74\,178.13 & 69\,212.46 & 41\,926.08 & 974.29 & 47\,482.28 \\
    问题 3 & 2025-06-21 & 35\,701.54 & 35\,432.45 & 20\,103.36 & 5.24 & 20\,260.50 \\
    问题 3 & 2025-09-23 & 74\,744.80 & 72\,213.13 & 45\,415.68 & 97.47 & 46\,558.94 \\
    问题 3 & 2025-12-21 & 99\,692.66 & 98\,984.11 & 63\,729.13 & 360.69 & 65\,154.58 \\
    \midrule
    问题 4-2 & 2025-03-20 & 74\,196.33 & --- & 46\,622.31 & 914.18 & 50\,923.43 \\
    问题 4-2 & 2025-06-21 & 35\,543.98 & --- & 17\,764.30 & 252.82 & 18\,587.09 \\
    问题 4-2 & 2025-09-23 & 74\,693.23 & --- & 49\,288.50 & 149.64 & 49\,839.31 \\
    问题 4-2 & 2025-12-21 & 99\,538.07 & --- & 76\,013.33 & 0.00 & 76\,013.33 \\
    \midrule
    问题 4-3 & 2025-03-20 & 74\,321.58 & 69\,229.87 & 43\,088.31 & 894.55 & 48\,781.27 \\
    问题 4-3 & 2025-06-21 & 35\,725.56 & 35\,440.19 & 17\,651.83 & 5.24 & 17\,769.08 \\
    问题 4-3 & 2025-09-23 & 74\,869.82 & 72\,305.05 & 47\,748.10 & 91.92 & 48\,876.69 \\
    问题 4-3 & 2025-12-21 & 99\,720.87 & 98\,953.92 & 75\,242.01 & 345.33 & 77\,242.42 \\
    \bottomrule
  \end{tabular}
\end{table}

四个日期中，\(2025\text{-}03\text{-}20\) 与 \(09\text{-}23\) 的负载形状接近而紧急购电量相差
\(1\,025.88\) 对 \(193.37\)~kWh，差异来自当日光伏预报偏差；\(2025\text{-}12\text{-}21\) 购电量最大
（\(99\,692.66\)~kWh）而紧急购电量最小，说明该日预报偏差最小、计划量本身已足够。四次滚动把
\(2025\text{-}03\text{-}20\) 的紧急购电量由 \(1\,025.88\)~kWh 降到 \(974.29\)~kWh，代价是
\(1\,949.15\)~元的违约与超额相关费用（\(1\,651.91+297.24\)），该日总费用净降 \(2\,171.74\)~元。
全窗口上四次滚动的净收益为 \(332\,389.18\)~元，而继续加密预报的剩余空间只有 \(58\,062.56\)~元
（第九章表~\ref{tab:ablation}）：滚动决策的收益集中在前几次预报，边际收益快速衰减。

回读校验：结果文件的工作表名、列序与模板一致；\texttt{result2.xlsx} 的「计划购电量」为
335 行 \(\times\) 147 列（334 个日期行），「充放电量」按 \(334\times6=2004\) 行重建，
「紧急购电量」按事件逐行生成（问题 2、3、4-2、4-3 分别为 1156、1190、1193、1187 行）；
全天汇总列与逐时段求和的最大偏差 \(0.0012\)~kWh，全天购电费最大相对偏差 \(4.32\times10^{-7}\)
（第九章审计 A9、A10）。

%!TEX root = ../main.tex
% =============================================================================
%  九、模型检验、灵敏度与稳健性      目标 2.5 页
%  评审在问：你怎么证明是对的？
%  本节含：① 三层验证与 L1--L5 链条（逐项表指向附录 E）
%          ② 非预期性置换检验  ③ 独立审计结论  ④ 基线与消融矩阵（表 27）
%          ⑤ 假设 A9 的量化  ⑥ 灵敏度与结论反转点（表 28）  ⑦ 样本外与伪验证自查
%          ⑧ 口径自洽性检验（三条机器可执行判据）
%  已移入附录 E：L1--L5 五级链条表、A1--A10 逐项审计表、非预期性逐区明细、
%                Sobol 完整指数、容量/功率双向扫描明细、L 形分解对拍、电价过程四项检验
%  红线：求解器返回 optimal 只说明"这个模型被解出来了"，不等于"模型建对了"。
% =============================================================================

\section{模型检验、灵敏度与稳健性}\label{chap:verify}

\subsection{三层验证与审计结论}\label{sec:verify-three}\label{sec:audit}

验证分三层。\textbf{思路层}指假设、信息结构与决策时序，不能靠数值反推，只能由假设与口径声明回答
（\chapref{chap:assump} 的 A1--A11、R1/R2/R3 与 \chapref{chap:q3} 的非预期性约束）。\textbf{模型层}指
变量、约束、目标与单位，“求解器返回最优”\emph{不构成}该层的证明：平衡式漏项或效率方向写反时 LP 同样
返回最优，故本文改用退化恒等式与独立复算替代。\textbf{结果层}指数值、口径与稳健性。三层按 L1--L5
落地：L1 退化与恒等式单测、L2 交叉模型对拍、L3 真实数据独立审计、L4 基线与消融、L5 统计与稳健性；
各级的手段、通过标准与实测值见 \appref{tab:audit-appendix}{E 独立审计报告}（L1--L3）、
表~\ref{tab:ablation}（L4）与表~\ref{tab:sens}（L5）。

\textbf{审计结论。} 审计由独立脚本执行：\textbf{不调用建模的目标函数}，只读原始附件与结果文件重算，
A1--A7 任一失败禁止进入写作、A8--A10 任一失败禁止提交。A1--A10 在内部解（10~min 级）上
\textbf{全部通过}：供给缺口与储电量递推的最大残差分别为 \(3.41\times10^{-13}\) 与
\(1.09\times10^{-11}\)~kWh（阈值 \(10^{-6}\)~kWh），储电量越界数、充放功率越界数、充放互斥违反数与
负购电量数\textbf{均为 0}，费用独立复算与优化目标值的相对偏差 \(\le10^{-6}\)（内部解）、结果文件
回读后 \(4.32\times10^{-7}\)（阈值 \(10^{-4}\)）。交付物层 A3、A4、A6、A7、A9、A10 通过；A1、A5 因
交付物只含逐日 4~h 聚合量而\textbf{不可核验}；A2 的 \(10^{-4}\)~kWh 残差可由“6 个块级充放电量与两个
储电量按 4 位小数取整“的误差上界 \(4\times10^{-4}\)~kWh 完全解释，不是物理越界。覆盖缺口有两处：
内部解审计覆盖问题 1--3 而问题 4 只做交付物层审计；交付物审计按块级口径执行，10~min 级结论只在模型
自身解上成立。

\textbf{最优性证书。} 本文\textbf{不报告“对偶相对间隙”}——多阶段模型的对偶量不构成可行域的最优性
证书。问题 1 改用与求解器无关的证书，即阈值律一致性检验：144 个槽的违反次数为 0、互补松弛最大残差
\(\le10^{-12}\)、\(E_t>0\) 时 \(\alpha_t=p_t\) 的偏差 \(\le10^{-6}\)（适用条件见
命题~\ref{prop:threshold}）。问题 2--4 另用 L 形（Benders）分解这条完全不同的数值路径复核最优值，
结果窗口 334 天上逐日最大相对间隙 \(3.76\times10^{-15}\)、全部收敛（明细见
\appref{sec:lshape}{E L 形分解对拍}）。\evid{C-Q1-02}

\textbf{非预期性置换检验。} 把某一次发布的预报在 12 天之间随机置换后重解，观察四个控制区
（槽 \(0\)--\(36\)、\(36\)--\(72\)、\(72\)--\(108\)、\(108\)--\(144\) 分别由 \(0{:}00\)、\(6{:}00\)、
\(12{:}00\)、\(18{:}00\) 预报控制）的最大变化量。\textbf{任何一次预报被置换，都不改变早于其发布时刻
的决策（变化恰为 0）}，即不存在用未来信息决策的泄漏；而“置换 \(0{:}00\) 预报会改变全天决策”说明检验
本身有效力。\(18{:}00\) 预报的零影响源于其覆盖的傍晚时段光伏实际与预报都接近 0、信息量极低，与
\chapref{chap:q3}“再增加预报时刻最多只能再省 \(58\,062.56\) 元（\(0.42\%\)）”相互印证。
逐区变化量见 \appref{tab:na-app}{E 非预期性置换明细}。\evid{C-Q3-01}

\subsection{基线与消融}\label{sec:ablation}

消融的公平性取三条约束：同一场景集（先固定存盘、所有变体读同一份）、同一指标定义与口径、同一预算。
表~\ref{tab:ablation} 的十组对照除末行外均用同一固定场景集（\(K=60\)、分层 medoid、层样本占比权重）
与同一实际轨迹。

\begin{table}[htbp]
  \centering
  \caption{基线与消融对照矩阵（费用单位：元，电量单位：kWh；全部为 2025-02-01 至 12-31 共 334 天、
           附件 1 固定电价口径；“相对 \(C^{base}\)”以无储能基线 \(16\,407\,319.63\) 元
           （\calref{C-01} 变体 V2）为分母）}
  \label{tab:ablation}
  \scriptsize
  \setlength{\tabcolsep}{3pt}
  \begin{tabularx}{\linewidth}{C{0.72cm} L{3.15cm} C{2.05cm} C{1.55cm} C{1.65cm} Y}
    \toprule
    编号 & 配置（新增或移除的模块） & 总费用 & 相对 \(C^{base}\) & 紧急购电量 & 结论 \\
    \midrule
    B0 & 无储能：净负荷全额购电（闭式 \calref{C-01}）
      & 16\,407\,319.63 & 0 & --- & 基线 \\
    R2 & 典型日预报：计划依据改为附件 1 的日均形状
      & 23\,484\,646.63 & \(-43.14\%\) & 2\,525\,548.76 & 比无储能基线还贵 \(7\,077\,327.00\) 元，不可用 \\
    E0 & 无安全边际：确定性等价（\(K=1\)、零残差情景）
      & 17\,398\,559.68 & \(-6.04\%\) & 1\,156\,618.02 & 移除后反而高于无储能基线；
        \textbf{安全边际的边际价值 \(3\,133\,767.55\) 元} \\
    R1 & 完全信息：\(0{:}00\) 知当日真值
      & 13\,411\,560.90 & 18.26\% & 0.00 & 预测误差成本 \(853\,231.23\) 元（下界） \\
    M & 主模型：\(0{:}00\) 一次预报 \(+\) 两阶段随机规划（\(K=60\)）
      & \textbf{14\,264\,792.13} & 13.06\% & 172\,015.82 & 主方案（基准行） \\
    N2 & 两次预报：\(0{:}00+12{:}00\)
      & 14\,131\,564.16 & 13.87\% & 122\,524.83 & 比主模型省 \(133\,227.98\) 元 \\
    N4 & 四次滚动：\(0/6/12/18\) \(+\) 非预期性
      & 13\,932\,402.95 & 15.08\% & 64\,522.09 & 比主模型省 \(332\,389.18\) 元 \\
    PI & 完全信息上界：四个时刻均知真值
      & 13\,874\,340.39 & 15.44\% & 6\,134.76 & 可获得信息价值的上界；四次预报已实现 \(85.1\%\) \\
    FY & 全年单一 LP：完全信息、单一储电量轨迹、无日决策解耦（192\,384 个变量）
      & 12\,228\,166.56 & 25.47\% & --- & 全局下界；\textbf{日决策解耦代价
        \(1\,183\,394.34\) 元（\(9.68\%\)）} \\
    K120 & 场景数加倍（\(K=120\)，其余同主模型）
      & 14\,235\,163.95 & 13.24\% & 141\,707.94 & 费用仅降 \(29\,628.19\) 元（\(-0.21\%\)），结论稳健 \\
    \bottomrule
  \end{tabularx}
\end{table}

表~\ref{tab:ablation} 中三行需单独说明。① \textbf{安全边际不是装饰}：整行移除后费用升至
\(17\,398\,559.68\) 元、比无储能基线高 \(6.04\%\)，紧急购电量由 \(172\,015.82\) 增至
\(1\,156\,618.02\)~kWh（6.72 倍），故“按情景分位数超购”是问题 2 的核心模块；与实测对照，
逐槽缺货频率 \(8.55\%\)、电量满足率 \(99.04\%\)，落在报童两界 \(1/10\)--\(1/5\) 的\textbf{下侧}
（储能除承担安全边际外还起对冲作用），故报童结论只作方向性近似与上下界 \evid{C-Q2-01}。
② \textbf{场景数不是敏感维度}：\(K\) 由 60 加倍到 120 时费用只降 \(0.21\%\)，但紧急购电量降
\(17.62\%\)、触发天数由 296 降到 277——费用结论稳健，而该分项指标对场景分辨率敏感，引用时
必须带 \(K\)。③ \textbf{日决策解耦的代价}（FY 行）：相差 \(1\,183\,394.34\) 元（\(9.68\%\)），
差额来自“每天固定初态、逐日终端条件”这一工程必需假设，方向保守而非乐观。

\textbf{假设 A9 的量化。} 题面只给光伏预报，未给负载预报，故 A9 把当日负载按已知处理，
此前的声明只是“会低估紧急购电量”。补上量级：只改\textbf{计划阶段}的负载信息，令
\(\widehat L_t=L_t+\varepsilon_t\)，\(\varepsilon\) 为按相对 MAE 标定的零均值误差
（\(\mathrm{MAE}_{\text{rel}}=\mathbb{E}|\varepsilon_t|/\overline L\)，\(\overline L=4626.03\)~kW），
模型、情景机制与结算口径全部不变，重跑 334 天（表~\ref{tab:loaderr}）。

\begin{table}[htbp]
  \centering
  \caption{负载预报误差对问题 2 的影响（只改计划期的负载信息；范围：2025-02-01 至 12-31 共 334 天；
           费用单位元、电量单位 kWh；样本：3 档相对 MAE；\(0\%\) 档与主结果逐位一致）}
  \label{tab:loaderr}
  \small
  \setlength{\tabcolsep}{5pt}
  \begin{tabular}{crrrrc}
    \toprule
    负载预报相对 MAE & 总费用 & 紧急购电量 & 出现天数 & 相对 \(C^{base}\) 节省率 & 相对 \(0\%\) 档费用 \\
    \midrule
    \(0\%\)（主口径） & \(14\,264\,792.13\) & \(172\,015.82\) & \(296\) & \(13.06\%\) & --- \\
    \(5\%\)           & \(15\,611\,126.88\) & \(293\,748.90\) & \(239\) & \(4.85\%\)  & \(+9.44\%\) \\
    \(10\%\)          & \(17\,472\,694.03\) & \(479\,838.20\) & \(202\) & \(-6.49\%\) & \(+22.49\%\) \\
    \bottomrule
  \end{tabular}
\end{table}

两条读数。① \textbf{量级}：相对 MAE 达 \(10\%\) 时问题 2 的总费用上升 \(22.49\%\)，且\textbf{相对无
储能基线的节省率转为负值}（\(-6.49\%\)）——在预报很差的微网里，两阶段计划加紧急购电的组合可能不如
不投运储能。这把 A9 的方向性声明变成可引用的区间：\(13.06\%\) 只适用于负载预报误差可忽略（或与光伏
残差同量级）的场景。② \textbf{一个反直觉但可解释的读数}：误差变大时紧急购电总量上升
（\(172\,015.82\to479\,838.20\)~kWh）而\textbf{触发天数反而下降}（\(296\to202\)）。原因不是计数口径，
而是模型的保守化：零均值误差使情景集变宽，计划层按更宽的分位数多购，“少而大”的缺口取代了“多而小”的
缺口。这也从反面说明主结果里 296 天触发、日均仅 \(581\)~kWh 的形态，是 \(0{:}00\) 光伏预报
\textbf{存在系统性时段偏差}（上午高估、下午低估，见 \ref{sec:sens}）的结果，而非随机噪声的结果。
\evid{C-Q2-01}

\subsection{灵敏度与结论反转点}\label{sec:sens}

灵敏度分两类：\textbf{单参数扫描}回答“参数变多少结论会反转”，\textbf{Sobol 全局灵敏度}回答“哪个参数
决定哪个输出量“。前者见表~\ref{tab:sens}，后者完整指数见 \appref{tab:sobol-app}{E Sobol 灵敏度指数}。

\begin{table}[htbp]
  \centering
  \caption{单参数灵敏度：扰动档位与“结论是否改变”的判断（费用单位 元，电量单位 kWh；
           单参数扫描除终端条件外均在同一组典型日参数下求解；完整档位实测值见
           \appref{tab:sens-app}{E 灵敏度明细}）}
  \label{tab:sens}
  \small
  \setlength{\tabcolsep}{4pt}
  \begin{tabularx}{\linewidth}{L{2.05cm} L{3.15cm} Y}
    \toprule
    参数 & 扰动档位 & 结论是否改变（含反转点） \\
    \midrule
    紧急购电倍数 \(\kappa\) & 2、3、5（基准）、8、10
      & 方向不变：\(\kappa\) 越大安全边际越大、缺口越小。反转点在 \(\kappa\to1\)——
        紧急购电与正常电价同价时缺货成本消失，“需要安全边际”的结论不再成立 \\
    效率口径 \(\eta_c,\eta_d\) & 单边各 0.9（\textbf{主口径}）、往返 0.9 等效单边、
        仅充电侧 0.9、仅放电侧 0.9
      & 区间内结论不变（四种读法下储能均有利可图），但节省额极差 \textbf{13.2\%}
        （式~\eqref{eq:q1-eff}）——同一模型可因效率口径不同报出相差八分之一的数 \\
    储电量运行区间 & 保持 10\%--90\% DOD：容量 7\,200--48\,000~kWh 共 7 档
      & 区间内结论不变：容量摆幅始终是主要瓶颈，仅解除容量上限可再省 22.14\% \\
    充放功率上限 \(\overline P\) & 1\,250、2\,500、5\,000（基准）、7\,500、10\,000~kW
      & 区间内结论不变：功率在 5000~kW 以上边际很小（\(\le0.47\%\)），
        但收紧到 2\,500~kW 时节省率由 26.90\% 降到 24.27\%、1\,250~kW 时降到 16.41\% \\
    终端条件 \calref{C-08} & 窗口：\(S_T\ge S_0\)（主口径）／自由终端；典型日：日周期／自由环状
      & 三组差异均不超过 0.12\%，结论不变；窗口口径下自由终端反而略贵（\(+0.08\%\)），
        说明终端约束并非靠“掏空”储能省钱 \\
    场景数 \(K\) & 60（主口径）、120
      & 费用结论不变（仅降 \(0.21\%\)）；但紧急购电量降 \(17.62\%\)、触发天数由 296 降到 277——
        \textbf{分项指标对场景分辨率敏感，引用时必须带 \(K\)} \\
    \bottomrule
  \end{tabularx}
\end{table}

Sobol 全局灵敏度须\textbf{区分输出量}：对\textbf{典型日全天购电费}起主导作用的是储电量上限
（一阶指数 \(0.2510\)）与下限（\(0.1099\)），其次是往返效率（\(0.0695\)）与充放功率上限
（\(0.0263\)）；对\textbf{节省率}起主导作用的却是电价水平（\(0.8464\)）与负载水平（\(0.4711\)），
储能参数均不超过 \(0.0033\)。原因是问题 1 的确定性模型中，电价、光伏与负载的\emph{水平}对全天购电费
只是整体缩放，不改变调度结构，故一阶指数为 0。结论是“哪一个参数决定哪一个输出量”，而非“所有参数都
重要“。

\subsection{样本外回测与伪验证自查}\label{sec:oos}

场景集与参数在结果窗口的 334 天残差上一次性估计并在同一窗口回放，属\textbf{样本内评估}。故本文
\textbf{不声称样本外有效性}，而以四条代偿性证据约束外推范围：场景保真度（均值相对差 \(0.082\%\)、
关键分位最大相对差 \(6.99\%\)、电价—净需求相关 \(0.2597\) 对 \(0.2620\)）、场景数稳健性
（\(K:60\to120\) 时费用只变 \(-0.21\%\)）、交叉模型对拍（差距方向保守，\(9.68\%\)）与非预期性置换
检验。此外另做两组真正的样本外实验：① 场景集只用 2025-02-01 至 06-30 的残差构造，在测试窗
07-01 至 12-31（184 天）上回放得总费用 \(8\,328\,592.63\) 元，而改用测试窗自身残差构造（样本内上界）
为 \(8\,314\,304.58\) 元，样本外劣化仅 \(14\,288.05\) 元（\(0.172\%\)）；② 四种情景生成器在样本外的
总费用极差仅 \(0.37\%\)（\appref{sec:q4-gen}{F 情景生成对照}），即结论不依赖生成范式，但紧急购电量
的分项差异达 \(48.5\%\)，该分项必须与情景口径同时引用。2025 年 1 月为储电量递推的预热期，不计入结果
（假设 A8）。

\textbf{伪验证自查。} 逐条排除八类伪验证：\textbf{自证循环}（审计脚本不共用建模费用函数）；
\textbf{以求解器状态代替物理校验}（同时报告缺口、递推、越界与互斥残差）；\textbf{比较不同样本}
（四种信息结构、四种预报时刻数、两个场景数与四种生成器均在\emph{同一场景集或同一实际轨迹}上配对
比较）；\textbf{训练集当测试集}（明确声明样本内，另做两组样本外回测）；\textbf{只报点估计}
（关键量均给多档：效率 4 档、容量 7 档、功率 5 档、\(\kappa\) 5 档、\(\lambda\) 11 档、\(\beta\) 6 档、
\(k\) 7 档）；\textbf{事后挑选指标}（全部口径在建模前固定）；\textbf{把“结果好看”当“模型正确”}
（问题 2 的 \(13.06\%\) 小于问题 1 的 \(26.90\%\)，差额可由预测误差成本 \(853\,231.23\) 元解释）；
\textbf{扰动无效}（灵敏度扰动作用于运行区间与功率上限本身，且每档重新求解）。逐条证据见
\appref{sec:fakeverify-app}{J 伪验证自查清单}。

\subsection{口径自洽性检验}\label{sec:caliber-check}

前两节回答“解是否可行、是否最优、是否稳健”，还剩一问：\textbf{同一个数字被报出来时属于哪个口径}。
本题共 11 个口径，而“每天的电价相同”“购电费用”“弃光”在四问之间含义并不相同，口径与数字一旦对不上就
无法复算，故再加三道\textbf{机器可执行}的自洽性判据。

\textbf{检验一（加权均价必须落在电价区间内）。} 对任一含购电费与购电量的结果，令
\(\bar p=\text{报告费用}/\text{报告电量}\)。报告电量是电价向量的逐时段加权平均，故必然满足

\begin{equation}
  \min_t p_t\ \le\ \bar p\ \le\ \max_t p_t ,
  \label{eq:caliber-price-bound}
\end{equation}

否则该费用与该电量不属于同一电价口径。五个主口径的实测值全部通过：问题 1、2、3 的电价向量同为附件 1
（\([0.3713,1.3952]\)~元/kWh），\(\bar p\) 分别为 \(0.5905\)、\(0.6137\)、\(0.5948\)；问题 4-2、4-3
为附件 4（\([0.0076,1.7936]\)），\(\bar p\) 为 \(0.6486\)、\(0.6281\)~元/kWh。\textbf{若某处误用了
全天单一均价}（本题均值 \(0.7662\)），则该口径下所有可行解的 \(\bar p\) 都只能等于 \(0.7662\)，
与上述任何一个值都不符——式~\eqref{eq:caliber-price-bound} 因此能把“报了一个价、算的是另一个价”这类
错配直接暴露。该判据另有一条副产品：问题 2 与问题 3 的\textbf{计划电量完全相同}而计划购电费相差
\(421\,349.41\) 元，恰等于 \calref{C-09} 另一读法下的 \(\sum_tp_tv_t=421\,349.41\) 元
（\S\ref{sec:ablation}），两个口径互相印证。

\textbf{检验二（同一量的三处取值必须一致）。} 每个进入正文的数字都有三个落点：
\texttt{notes/数值台账.md}、\texttt{intermediate/*.json} 与 \texttt{results/result*.xlsx}。做法是从
\textbf{结果文件回读}并重算、再与中间结果对拍：全天购电量的块级最大残差
\(2.96\times10^{-4}\)~kWh（四位小数取整上界 \(7.0\times10^{-4}\)~kWh）、全天购电费最大相对偏差
\(4.32\times10^{-7}\)、日期首末严格为 2025-02-01 与 2025-12-31。题面表 1／表 2 的“指定时段”与槽位的
对应由 \texttt{config.py} 单点定义，正文表格与结果文件因此不可能各写各的。

\textbf{检验三（跨问题比较必须同费用结构）。} 最容易被误用的比较是把“只有计划购电费”与
“计划 \(+\) 紧急 \(+\) 调整相关”两种口径直接相减。本文将信息价值 \calref{C-10} 定义为\textbf{同一
场景集、同一实际轨迹}下仅改变预报时刻数的费用差（\(n=1\) 的 \(14\,264\,792.13\) 元 \(\to\) \(n=4\) 的
\(13\,932\,402.95\) 元），两口径费用项构成相同；问题 4-2 与 4-3 的费用项不同，故只作分项并列、不相减
得节省率。三条判据的合取构成本文的结论防线：\textbf{可行、最优、稳健、口径可分、三处一致、可比者
才比}。\evid{C-Q5-01}

%!TEX root = ../main.tex
% =============================================================================
%  十、模型评价与推广            建议 2 页
%  评审在问：你诚实吗？能推广吗？
%  本节必须出现：① 创新五段式（缺陷—改动—证据—代价—边界，每条都要有消融/对照数字）
%                ② 局限四要素（指名假设编号 + 影响方向 + 所需数据 + 校准动作）
%                ③ 行业口径声明（电能量收益的保守下界）
%                ④ 推广的临界条件（价格接受型、不外送、单一储能）
%                ⑤ 结论与建议（不过度声称）
%  禁止空话（如"数据有限""模型有待改进"）；一切改进都必须给出对照数字。
% =============================================================================

\section{模型评价与推广}

\subsection{主要创新点}

每条创新按五段式写：\textbf{缺陷}（残差、约束紧性或对照实验）\(\rightarrow\)\textbf{改动}（可写成公式）
\(\rightarrow\)\textbf{证据}（同一场景集、同一实际轨迹的消融或对照）\(\rightarrow\)\textbf{代价}
\(\rightarrow\)\textbf{边界}（失效条件）。

\begin{innovation}{1}{报童型安全边际：把"该多买多少"写成临界分位数（S 级）}
\ip{缺陷}问题 2 的计划只能依据 \(0{:}00\) 的光伏预报制定，而结算时缺口的代价是正常电价的 \(5\) 倍、
超购只白付电费（两侧成本不对称 \(5:1\)）；按预测中心值制定计划等价于假设预报无误差，但实测逐槽缺货
频率为 \(8.55\%\)、全窗口紧急购电量 \(172\,015.82\)~kWh；确定性阈值律只回答"何时充放电"，
回答不了"计划量该留多少余量"。
\ip{改动}对单槽的超购—缺货权衡建立报童近似：超购成本 \(c_o=p_t\)（超购部分即被弃置），
缺货成本 \(c_u=4p_t\)（紧急购电 \(5p_t\) 中本应付的 \(p_t\) 已计入计划量），故临界分位数
\(F^*(m)=c_u/(c_u+c_o)=4/5\)，即最优缺货概率 \(1/5\)；若超购电量可转存电池并在后续高价时段使用，
取 \(c_o\approx p-p_{\min}=0.7662-0.3713=0.3949\)~元/kWh、\(c_u=4p=3.0648\)~元/kWh 时
\(F^*\approx0.886\)，缺货概率降至约 \(1/10\)。实现上把安全边际写成"按情景缺口分位数超购"的第一阶段量，
第二阶段只对残余缺口紧急购电。
\ip{证据}同一场景集、同一实际轨迹的消融（表~\ref{tab:ablation} 的 E0 行）：移除安全边际
（确定性等价、\(K=1\)）后总费用由 \(14\,264\,792.13\)~元升至 \(17\,398\,559.68\)~元（\(+22.0\%\)），
紧急购电量由 \(172\,015.82\)~kWh 增至 \(1\,156\,618.02\)~kWh（\(6.72\) 倍），且高于无储能基线
\(16\,407\,319.63\)~元（节省率 \(-6.04\%\)）——\textbf{安全边际的边际价值为 \(3\,133\,767.54\)~元}。
方向性检验：\(\kappa\) 由 \(2\) 增到 \(10\) 时安全边际单调升、期望紧急购电单调降（分档数值见
附录表~\ref{tab:sens}），与报童结论方向一致。\evid{C-Q2-01}
\ip{代价}每个情景需解一次第二阶段 LP（单日 \(576+144K\) 个变量，主口径 \(K=60\)），且场景集须先
固定存盘以保证所有变体可比；安全边际使计划购电量偏高、超购部分被弃置（问题 2 窗口盈余处置量
\(3\,106\,390.84\)~kWh）。
\ip{边界}四条失效条件：① 单槽独立近似忽略跨槽储电量耦合与终端条件；② 储能本身即对冲工具，
会替代一部分安全边际（实测缺口频率落在两界下侧）；③ 依赖"紧急购电在结算时刻发生、计划量不可回溯修改"
的时标设定；④ 残差分布未知，只能以样本分位数替代（本文用 \(K=60/120\) 两档量化）。故报童结果只作
\textbf{方向性近似与上下界}，精确解仍由两阶段 LP 给出。
\end{innovation}

\begin{innovation}{2}{预报更新价值的量化：把"要不要加密预报"变成可判定的阈值（S 级）}
\ip{缺陷}题面直接追问"是否需要引入其他时刻的预报"：只答"建议引入更多预报"既不可证伪也无法决策，
只画收益曲线又必然得出"越多越好"。
\ip{改动}在同一固定场景集、同一实际轨迹上\textbf{只改变信息结构}，定义
\(\mathrm{VoI}(n)=\bigl(C(1)-C(n)\bigr)/C(1)\)（\calref{C-10}），并同时给出完全信息上界与"已实现的
可获得信息价值比例"。
\ip{证据}\(C(1)=14\,264\,792.13\)~元 \(\rightarrow\) \(C(2)=14\,131\,564.16\)~元 \(\rightarrow\)
\(C(4)=13\,932\,402.95\)~元，完全信息上界 \(13\,874\,340.39\)~元；\(\mathrm{VoI}(2)=0.93\%\)、
\(\mathrm{VoI}(4)=2.33\%\)，四次预报已实现可获得信息价值的 \(85.13\%\)，再增加预报最多只能再省
\(58\,062.56\)~元（占 \(C(1)\) 的 \(0.42\%\)）；紧急购电量由 \(172\,015.82\)~kWh 降到 \(64\,522.09\)~kWh
（上界 \(6\,134.76\)~kWh）。机理见非预期性置换检验：置换 \(18{:}00\) 预报对全部控制区的影响为 0
（该次预报信息量极低），缺口集中在 \(6{:}00\)--\(12{:}00\)，即若要加密应加在上午而非均匀加密。\evid{C-Q3-01}
\ip{代价}需额外求解两组模型（\(n=2\) 与 \(n=4\)，问题 3 共 \(334\times4=1336\) 个子问题），每天四次
决策带来四次通信、上报与结算复杂度，而全部收益为 \(332\,389.18\)~元（占 \(C(1)\) 的 \(2.33\%\)）。
\ip{边界}结论依赖本题的预报误差结构与惩罚倍数：① 预报误差水平或紧急购电倍数放大时收益平稳带
右移，须重算；② \(18{:}00\) 的预报在本题数据上信息量近于 0，换一套数据后不成立；③ 成本侧未实测——
判据是"单次加密的年化净收益是否为正"，即预报获取的年化成本须低于 \(58\,062.56\)~元；本文未取得气象服务报价，故只给出该阈值而不下"是否值得采购"的结论。
\end{innovation}

\begin{innovation}{3}{波动电价下的联合情景与 CVaR 风险—费用前沿（B 级）}
\ip{缺陷}问题 4 的电价本身随机：实测电价区间 \(0.0076\)--\(1.7936\)~元/kWh、残差组内标准差
\(0.1238\)~元/kWh、一阶自相关 \(0.8589\)，且与净需求正相关（池化 \(0.2602\)）。若对电价、负载、光伏各自
独立抽样，会生成"午间光伏大发同时电价高企"这类不可能情景，从而低估尾部风险；只报期望费用则没有回答"波动"。
\ip{改动}用两层结构建模电价：分时骨架取附件 4 的逐时段均值（与附件 1 的最大偏差
\(5.15\times10^{-5}\)~元/kWh），叠加按日内自相关刻画的随机残差；情景按"日净需求残差 \(\times\) 日均电价"
二维分层取代表历史日（层内 medoid、权重为层样本占比），天然保留电价—净需求相关；风险侧用
Rockafellar--Uryasev 线性化 \(\min\ \mathbb{E}[\text{Cost}]+\lambda\,\mathrm{CVaR}_{0.95}\)
（\calref{C-11}）并对 \(\lambda\) 扫描。
\ip{证据}情景保真度：情景的电价—净需求相关 \(0.2597\) 对历史 \(0.2620\)（差 \(0.0023\)，判据 \(0.05\)），
均值相对差 \(0.082\%\)、关键分位最大相对差 \(6.99\%\)（分块 Bootstrap 对照 \(3.41\%\) 与 \(20.44\%\)）。
前沿（\(56\) 天子样本、日均值）与\textbf{真实轨迹回测}同向：\(\lambda=0\) 时期望 \(41\,995.59\)、
\(\mathrm{CVaR}_{0.95}=52\,176.36\)、回测 \(45\,694.27\) 元/日；\(\lambda=0.5\) 时期望 \(42\,280.55\)
（\(+0.68\%\)）、CVaR \(50\,570.72\)（\(-3.08\%\)）、回测 \(45\,863.46\)（\(+0.37\%\)）、紧急购电量
\(309.62\)~kWh/日（\(-35.52\%\)）；\(\lambda=1\) 时期望 \(42\,411.56\)（\(+0.99\%\)）、CVaR
\(50\,510.46\)（\(-3.19\%\)）、回测 \(45\,952.85\)（\(+0.57\%\)）。即约 \(1\%\) 的期望费用溢价可换取
约 \(3.2\%\) 的尾部风险下降，前沿在 \(\lambda\approx0.6\)--\(0.9\) 后进入平台（\(\lambda=1\) 的 CVaR
反而略高于 \(\lambda=0.9\)）。
\ip{代价}每个 \(\lambda\) 点需重解子样本上的随机模型（共 11 个点），且前沿是"事前"结果、须配真实
轨迹回测才能用于决策；联合情景把电价与负荷抽样耦合成二维分层，情景构造与保真度检验的实现成本上升。
\ip{边界}① 情景集来自 334 天历史，不能外推政策或市场机制突变；
② CVaR 的置信水平固定为 \(0.95\)，未做 \(\alpha\) 扫描；
③ 情景电价截断在实测范围内，出现负电价或极端尖峰时须重新标定；
④ \(\lambda\) 的推荐点取决于决策者风险偏好，本文只给风险—费用前沿，不给唯一答案（\calref{C-11}）。
\end{innovation}

\subsection{模型优点}

\begin{itemize}
  \item \textbf{结构统一}：四问共用同一套物理模型骨架（式~\eqref{eq:q1-balance}--\eqref{eq:q1-nonneg}），
        差别只在信息结构与风险处理，每层新增贡献都可单独消融（表~\ref{tab:ablation} 十组对照）；
  \item \textbf{可解释}：对偶水价值与阈值律给出与求解器无关的策略形式（\(\alpha_t\le\eta_c\nu_t\) 充电、
        \(\alpha_t\ge\nu_t/\eta_d\) 放电，全部 144 个槽违反次数 0），使充放电时段可解释、可质疑；
  \item \textbf{可复算}：全部数值由脚本从原始附件重算，并通过原始可行性（供给缺口 \(3.41\times10^{-13}\)~kWh）、
        最优性（阈值律一致性）与费用独立复算（相对误差 0）三重检验（第九章）。
\end{itemize}

\subsection{模型局限}



\textbf{本文目标函数只含电能量购电费用，未计入容量补偿、辅助服务收益与绿电环境价值，
也未计入电池循环衰减与安全成本；因此本文所得节约幅度应视为储能电能量价值的保守下界。}
用户到户电价还含输配电费、系统运行费与政府性基金及附加，本文只优化第一项：附加电费按电量比例收取
等价于对全部电价加同一常数，最优调度会因往返效率小于 1 而改变；按需量（容量）收取则更会实质性改变
策略。两种情形都须重解，故本文的节省幅度既不是储能项目投资可行性的结论，也不是全成本口径的收益。
按电量比例计收的附加费等价于给全部电价加上同一常数，由于往返效率小于 1，最优充放电量会相应减少（本文未重解，故不给出该口径下的数值）；按最大需量或容量计收的附加费则会实质性改变策略，须在目标函数中显式建模后重解。

\subsection{模型推广与临界条件}

本文模型在以下条件下成立，偏离时须重新校准而非直接套用：

\begin{itemize}
  \item \textbf{价格接受型}：电价外生给定（附件 1 与附件 4）。若储能或微网规模足以影响出清价格，
        目标函数须引入价格对自身报价的响应，模型由 LP 变为带均衡约束的问题，本文的阈值律与安全边际
        结论都不能直接沿用；
  \item \textbf{不外送（假设 A1）}：允许按同一电价售电时，目标函数需增加售电收益项，盈余处置量
        \calref{C-06} 由"必要"变为"可选"，收益结构整体改变（售电可能成为主要收益来源），须重解；
  \item \textbf{单一储能}：推广到多储能或多微网协同后状态维度随设备数增长，"逐日解耦"的代价须
        重新评估——本文数据上为 \(1\,183\,394.34\)~元（\(9.68\%\)，表~\ref{tab:ablation} 的 FY 行）；
  \item \textbf{价差水平足够大}：套利有利的盈亏平衡条件是谷峰价比高于 \(1/(\eta_c\eta_d)=1.2346\)
        （主口径单边各 \(0.9\)）；价差压到该阈值附近时纯套利不再有利，模型退化为"以规避紧急购电为主"
        的保险型策略。容量与功率扫描给出同类临界：容量降到 \(7\,200\)~kWh 时节省率由 \(26.90\%\) 降到
        \(19.59\%\)，功率降到 \(1\,250\)~kW 时降到 \(16.41\%\)；
  \item \textbf{负载刚性}：若引入需求响应，式~\eqref{eq:q1-balance} 中的 \(L_t\) 变为决策量，可行域
        变为含柔性负荷约束的多面体，阈值律须重新推导（水价值与负荷转移的边际价值共同决定阈值），
        但"两阶段 + 安全边际"框架可沿用；
  \item \textbf{单一电能量结算}：容量补偿、辅助服务与绿电环境价值在多数地区顺序出清、难以同时兑现，
        收益\textbf{不可简单叠加}；计入后本文结论方向不变（当前口径为保守下界），但幅度须重算。
\end{itemize}

\textbf{结论与建议。}在本题参数范围内（容量 \(12\,000\)~kWh、充放功率 \(5000\)~kW、单边效率 \(0.9\)、
紧急购电 \(5\) 倍价、不外送），储能的价值是把低谷电量与午间光伏盈余搬运到高价时段：典型日节省
\(26.90\%\)，334 天窗口在预报不确定下降到 \(13.06\%\)，波动电价下为 \(11.35\%\)；容量摆幅是主要瓶颈
（仅解除容量上限可再省 \(22.14\%\)，仅解除功率上限再省 \(1.66\%\)）；预报价值已接近饱和（四次预报实现
可获得信息价值的 \(85.13\%\)，再加密的收益上限为 \(58\,062.56\)~元）；风险偏好可通过 \(\lambda\) 调节，
约 \(1\%\) 的期望费用溢价可换约 \(3.2\%\) 的尾部风险下降。据此给出三条建议：① 预算有限时优先扩容储能
容量而不是加密预报（容量档位的收益区间为 \(19.59\%\)--\(42.56\%\)，远大于加密预报的上限 \(0.42\%\)）；
② 安全边际必须保留（边际价值 \(3\,133\,767.54\)~元），并应加在上午的缺口高发时段；③ 部署前按
表~\ref{tab:limitation} 逐项校准效率、衰减、负载预测与终端条件。上述结论只对"价格接受型、不外送、
单一储能、以电能量购电费用最小"这一设定成立，不构成储能项目的投资可行性结论。


% ---- 参考文献 --------------------------------------------------------------
%!TEX root = ../main.tex
% =============================================================================
%  参考文献（格式规范第七条：引用他人或公开资料成果必须列出参考文献，
%            并在正文引用处予以标注；格式按 GB/T 7714）
%
%  ⚠️ 三条引用纪律
%     1 宁可少引三条，不可引错一条——引用前必须核对官方原文
%       （国家发展改革委 / 国家能源局官网、\url{https://openstd.samr.gov.cn/）}
%     2 不得编造政策文号、标准号、企业数据、参考文献（属学术不端）
%     3 行业数字只能用于背景与合理性检验，不得当作本题计算结果，
%       不得出现在结果表或「计算结果」章节
%
%  本节只保留三条经官网原文核对通过的政策文件；文号、发文机关与成文日期逐条核对，
%  无法核实的条目一律删除，不保留占位。
% =============================================================================

% thebibliography 环境自带标题（ctexart 的 \refname 为「参考文献」），
% 因此不再手写 \section*{参考文献}，避免出现两个标题。
\renewcommand{\refname}{参考文献}

\begin{thebibliography}{99}
  \setlength{\itemsep}{2pt}

  \bibitem{ref136}
  国家发展改革委, 国家能源局.
  关于深化新能源上网电价市场化改革~促进新能源高质量发展的通知:
  发改价格〔2025〕136号[Z]. 2025-01-27.
  \url{https://www.gov.cn/zhengce/zhengceku/202502/content_7002959.htm}
  \note{正文第一章引用；发文机关、文号与成文日期按中国政府网「国务院部门文件」栏目核对（2026-09-10）}

  \bibitem{ref1093}
  国家发展改革委.
  关于进一步完善分时电价机制的通知:
  发改价格〔2021〕1093号[Z]. 2021-07-26.
  \url{https://www.gov.cn/zhengce/zhengceku/2021-07/29/content_5628297.htm}
  \note{分时电价的时段划分与峰谷价差机制出处；本文用于分时电价骨架与峰谷价比的背景说明}

  \bibitem{ref1217}
  国家发展改革委, 国家能源局.
  电力现货市场基本规则（试行）:
  发改能源规〔2023〕1217号[Z]. 2023-09-07.
  \url{https://www.gov.cn/zhengce/zhengceku/202309/content_6904881.htm}
  \note{「日前计划与实时偏差」结算机制的政策原型；本文问题 2、3 的信息结构与偏差定价口径参考此机制}

\end{thebibliography}


% 正文结束页标记：tools/check_paper.py 用它核算正文页数是否 ≤30
\label{pg:bodyend}

% ---- 附录不进 main.pdf ------------------------------------------------------
% 交付形态（用户指定）：
%   main.pdf      只含 摘要页 + 正文 + 参考文献（正文 ≤30 页，不含任何源程序）
%   appendix.pdf 只含 附录 A--O（页数不限，含全部源程序，需与正文一起打印装订）
%   full.pdf     正文 + 附录合并版（备用）
% 三者共用同一套 sections/；指向附录的引用统一用 \appref，附录未载入时自动降级为
% 备用文字（如「附录 F」），因此三个入口都不会出现 ??。
% 请勿在此处再加 %!TEX root = ../main.tex
% =============================================================================
%  附录（格式规范第四条：页数不限，必须打印并与正文装订在一起提交）
%  附录必备八项（依据 format2026.doc 第五条、第十一条）：
%    1 支撑材料文件清单        -> 附录 A  ← 必需
%    2 全部完整、可运行的源程序 -> 附录 E  ← 缺少可能被取消评奖资格
%    3 运行环境与依赖版本      -> 附录 B
%    4 运行说明               -> 附录 B
%    5 随机种子清单            -> 附录 H
%    6 结果文件与论文表格对应表 -> 附录 C
%    7 数据字典               -> 附录 F
%    8 较长的推导              -> 附录 G
%  ⚠️ 若确实没有用到程序，须写明「本论文没有用到程序」（第五条）
%  ⚠️ 若确实没有支撑材料，须注明「本论文没有支撑材料」（第十一条）
%  ⚠️ 附录中同样不得出现任何身份信息（第六条）：校名、队伍编号、姓名、所在区域一律不写
%  ⚠️ 附录 A 的文件清单与附录 E 的 \codefile 清单都按工程**实际文件**逐项写，
%     新增脚本或结果文件必须同步补入这两处（附录 E 缺文件会被判为红线）。
% =============================================================================

\section{支撑材料文件清单}

支撑材料压缩包 \texttt{support.zip} 打包本工程「可复现必需」的文件：论文源文件、全部源程序、
五个结果文件、八张数据图的矢量文件，以及作为论文数字来源的汇总类中间结果。
下表与本工程下的实际文件逐项对应；附录~\ref{sec:app-code} 所列源程序与本表的 \texttt{code/} 完全一致。

\begin{lstlisting}[style=plain,basicstyle=\ttfamily\footnotesize\songti]
support.zip
+-- main.tex                        论文装配（宏包输入、附录起点、正文结束标记）
+-- preamble.tex                    导言区：格式合规配置与写作宏（口径框/证据号/占位符）
+-- build.sh                        编译入口（自动探测 latexmk+xelatex / xelatex / tectonic）
+-- sections/                       正文与附录源文件，13 个
|   +-- 00_abstract.tex             摘要专用页
|   +-- 01_problem.tex              一、问题重述与分析
|   +-- 02_assumptions.tex          二、模型假设与符号说明
|   +-- 03_data.tex                 三、数据预处理与特征分析
|   +-- 04_q1.tex                   四、问题 1 模型
|   +-- 05_q2.tex                   五、问题 2 模型
|   +-- 06_q3.tex                   六、问题 3 模型
|   +-- 07_q4.tex                   七、问题 4 模型
|   +-- 08_results.tex              八、结果与分析
|   +-- 09_validation.tex           九、模型检验与灵敏度
|   +-- 10_evaluation.tex           十、模型评价与推广
|   +-- refs.tex                    参考文献（GB/T 7714）
|   +-- 90_appendix.tex             本附录
|
+-- tools/
|   +-- check_paper.py              提交前合规自查（--strict 为定稿模式）
|   +-- check_numbers.py            论文数字与数值台账的机械对拍
|   +-- verify_ref.py               对照稿红队审查的只读反算脚本（复现 notes/F 的全部断言）
|   +-- package_support.sh          本压缩包的打包脚本（保证本清单与实际文件一致）
|
+-- code/                           全部源程序，与附录 E 逐项一致
|   +-- config.py                   口径、参数与附件读取（全文唯一参数来源）
|   +-- lp_engine.py                单日两阶段 LP、滚动子问题 LP、场景削减
|   +-- q1.py  q2.py  q3.py  q4.py  四问求解与结果文件写出
|   +-- results_io.py               结果工作表写出器（按附件 5 模板的列序）
|   +-- exp_scan.py                 容量/功率/效率扫描与报童经验覆盖率
|   +-- exp_sobol.py                全局灵敏度（Sobol 指数）
|   +-- exp_na.py                   非预期性置换检验
|   +-- exp_oos.py                  样本外回测（训练窗构造情景集、测试窗回放）
|   +-- exp_benchmark.py            全年单一 LP 基准实验（完全信息参照）
|   +-- audit.py                    独立审计 A1--A10（不共用建模费用函数）
|   +-- paper_numbers.py            汇总数值台账 paper_numbers.json
|   +-- plot_figures.py             八张数据图的绘图脚本
|   +-- plot_figures_v2.py          v2 新增五张数据图（相依结构/生成器对照/风险画像/鲁棒前沿/分解）
|   +-- copula.py                  二元 Copula 工具库与 C 藤 pair-Copula（密度/h 函数/拟合/抽样）
|   +-- scen_copula.py             KDE + Copula 联合情景生成器与保真度度量
|   +-- exp_copula.py              相依结构拟合、四种情景生成器样本外对照、分时段 CVaR 风险画像
|   +-- exp_chance.py              机会约束对照（安全边际分位数扫描与报童互证）
|   +-- exp_dro.py                 盒式分布鲁棒对照（保守度—费用曲线与最坏情形评估）
|   +-- exp_loaderr.py             假设 A9 的量化：负载预报误差 0%/5%/10% 三档重跑问题 2
|   +-- decomp.py                  L 形（Benders）分解对照（解析多割、与直接求解对拍）
|   +-- run_all.py                  一键复现
|   +-- requirements.txt            依赖清单（版本锁定）
|
+-- results/                        五个结果文件（工作表名与列序按附件 5 模板）
|   +-- result1.xlsx                问题 1（典型日）：计划购电量、充放电量
|   +-- result2.xlsx                问题 2：计划购电量、充放电量、紧急购电量
|   +-- result3.xlsx                问题 3：计划购电量、调整购电量、充放电量、紧急购电量
|   +-- result4-2.xlsx              问题 4-2：与 result2 同结构
|   +-- result4-3.xlsx              问题 4-3：与 result3 同结构
|
+-- figures/                        图 1 至图 10
|   +-- fig01_timeline.tex          图 1 决策时序图（TikZ 源文件）
|   +-- fig02_dependency.tex        图 2 问题依赖图（TikZ 源文件）
|   +-- fig03_q1_dual_price.pdf     电价、水价值与阈值型策略（问题 1）
|   +-- fig04_q1_capacity_power.pdf 容量与充放功率上限双向扫描（问题 1）
|   +-- fig05_cost_waterfall.pdf    问题 1 成本分解瀑布
|   +-- fig06_voi.pdf               信息价值曲线（问题 3）
|   +-- fig07_q2_emergency.pdf      逐日紧急购电与报童两界（问题 2）
|   +-- fig08_q4_price.pdf          电价过程骨架与残差结构（问题 4）
|   +-- fig09_cvar_frontier.pdf     CVaR 风险--费用前沿与真实回测（问题 4）
|   +-- fig10_strategy_mechanism.pdf 最优策略机制（问题 1）
|   +-- README.md                   图规格、配色语义与出图约定
|
+-- intermediate/                   中间结果（论文数字的来源），分两组
|   +-- 随包提交（汇总与审计类，单个不超过 0.21 MB）
|   |   +-- audit.json  audit_report.md          审计 A1--A10（附录 D 的机器可读与可读版）
|   |   +-- paper_numbers.json                   数值台账（对拍基准，含四问、对偶、消融、Sobol、CVaR）
|   |   +-- q1.json  q1_dual.json  q1_dual2.json 问题 1 解与最优性证书
|   |   +-- q2_margin.json  exp_benchmark.json  安全边际分解、全年单一 LP 基准
|   |   +-- q4_price.json  q4_cvar.json          电价过程检验与情景保真度、CVaR 前沿
|   |   +-- exp_scan.json  exp_sobol.json  exp_na.json  扫描、灵敏度、非预期性
|   |   +-- full_year_lp.txt                     全年单一 LP（完全信息参照）的目标值
|   +-- 不随包（逐日解转储，单个 2.7--4.4 MB，可由 run_all.py 一键重建）
|       +-- q2_main.json  q2_noMargin.json  q2_cmp_R1.json  q2_cmp_R2.json  q2_cmp_R3.json
|       +-- q2_K120.json  q2_full_strat.json  q2_term_R1_eq.json  q2_term_R1_free.json
|       +-- q2_term_R3_free.json  q3_n1.json  q3_n2.json  q3_n4.json  q3_perf.json
|       +-- q4_2.json  q4_3.json
|
+-- notes/                          口径与台账文档，7 个
    +-- 写作交底.md                 本队实际模型口径（唯一口径说明）
    +-- 数值台账.md                 论文数字的表格版台账
    +-- A_坑点与验证.md  B_模型与算法.md  C_写作规范.md
    +-- D_行业与引用.md  E_知识点与模板.md
\end{lstlisting}

三点说明：

\begin{enumerate}
  \item \textbf{体积}：上述随包内容合计约 3.4~MB（其中五个结果文件约 2.1~MB、
        八张数据图约 0.3~MB），低于支撑材料压缩包 20~MB 的上限（格式规范第十一条）。
  \item \textbf{中间结果分级打包}：\texttt{intermediate/} 下第二组\textbf{逐日解转储}
        （\texttt{q2\_main.json}、\texttt{q2\_noMargin.json}、\texttt{q3\_n4.json}、\texttt{q4\_2.json} 等，
        单个 2.7--4.4~MB）可由 \texttt{python3 code/run\_all.py} 一键重建，
        且其中每个进入论文的数字都已汇总进 \texttt{paper\_numbers.json}，故不随压缩包提交；
        包内保留体积小、且为论文数字与审计直接依据的汇总类文件（上表已逐一列出）。
  \item \textbf{不随包提交的中间产物}：LaTeX 编译产物、Python 字节码缓存、以下划线开头的临时脚本与
        \texttt{.npy} 临时数组，均为可再生的工作文件，不属于支撑材料。
\end{enumerate}

\section{运行环境与运行说明}

\paragraph{运行环境}

\begin{table}[htbp]
  \centering
  \caption{本文的运行环境与依赖版本清单（对象：操作系统、解释器、数值与绘图库、求解器接口；
            范围：全部求解、审计与出图的实际执行环境；单位与量纲不适用；样本量为 1 套环境配置）}
  \label{tab:env}
  \footnotesize
  \begin{tabularx}{\linewidth}{L{3.4cm} L{4.6cm} Y}
    \toprule
    项目 & 版本 & 说明 \\
    \midrule
    操作系统 & macOS 15.7.5（arm64，Build 24G624） & 全部求解、审计与出图单机完成，未使用集群 \\
    Python & 3.13.12 & 唯一解释器 \\
    NumPy / pandas & 2.4.3 / 3.0.1 & 数组运算与附件 \texttt{.xlsx} 读写 \\
    SciPy & 1.18.1 & 线性规划 \texttt{linprog} 与统计检验 \\
    求解器 & SciPy \texttt{linprog(method=“highs”)}，HiGHS 后端（随 SciPy 1.18.1 分发）
      & 四问模型在场景离散后均为 LP，由 HiGHS 精确求解，未使用商业求解器 \\
    Matplotlib & 3.11.1 & 图 3--图 10 的数据图；图 1、图 2 为 TikZ 源文件，由 XeLaTeX 排版 \\
    scikit-learn & 1.9.0 & 仅用于 $k$-means 场景削减的对照路径 \\
    openpyxl & 3.1.5 & 结果文件 \texttt{.xlsx} 的读写 \\
    \bottomrule
  \end{tabularx}
\end{table}

依赖清单锁定在 \texttt{code/requirements.txt}（\texttt{numpy==2.4.3}、\texttt{pandas==3.0.1}、
\texttt{scipy==1.18.1}、\texttt{openpyxl==3.1.5}、\texttt{matplotlib==3.11.1}、
\texttt{scikit-learn==1.9.0}），版本号写确切值、不写区间。

\paragraph{运行说明}

\begin{enumerate}
  \item \textbf{依赖安装}：\texttt{python3 -m pip install -r code/requirements.txt}；
  \item \textbf{数据就位}：把附件 1--4 与附件 5 模板放在工作区的
        \texttt{CUMCM2026Problems(1)/C题/附件/}（模板置于其 \texttt{附件5/} 子文件夹）。
        该路径由 \texttt{code/config.py} 的 \texttt{ATTACH\_DIR} 给出，并由脚本自身位置反推，
        整个工程文件夹移动后无需改动代码；
  \item \textbf{依次运行四问}（在 \texttt{code/} 下执行，命令参数即论文主口径）：

\begin{lstlisting}[style=plain,basicstyle=\ttfamily\footnotesize\songti]
python3 q1.py                                                      # 问题 1：典型日 LP
python3 q2.py --tag main --K 60 --scen strat --info R3 --write    # 问题 2：334 天
python3 q3.py --epochs 0,6,12,18 --tag n4 --K 60 --write          # 问题 3：四次预报
python3 q4.py --mode q42 --K 64 --write                           # 问题 4-2：波动电价
python3 q4.py --mode q43 --K 64 --write                           # 问题 4-3：滚动
\end{lstlisting}

        \texttt{--write} 表示按附件 5 模板写出对应结果文件；
        条件项（论文中作为对照或检验给出）另有：\texttt{q4.py --mode tests --K 64}（电价过程检验与
        情景保真度）、\texttt{q4.py --mode cvar --K 64 --step 6}（CVaR 前沿）、
        \texttt{exp\_scan.py}（容量、功率与效率扫描、报童经验覆盖率）、
        \texttt{exp\_sobol.py}（Sobol 全局灵敏度）、\texttt{exp\_na.py}（非预期性置换检验）、
        \texttt{exp\_benchmark.py}（全年单一 LP 的完全信息基准）、
        \texttt{plot\_figures.py}（八张数据图）；
  \item \textbf{独立审计}：\texttt{python3 audit.py}。脚本只读原始附件与五个结果文件，
        独立复算 A1--A10，写出 \texttt{intermediate/audit.json} 与 \texttt{audit\_report.md}；
  \item \textbf{数值台账}：\texttt{python3 paper\_numbers.py}，把全部中间结果汇总为
        \texttt{intermediate/paper\_numbers.json}（论文数字的对拍基准）；
  \item \textbf{对拍与自查}：\texttt{python3 ../tools/check\_numbers.py} 列出论文中找不到
        台账出处的数字；\texttt{python3 ../tools/check\_paper.py --strict} 做定稿模式合规自查
        （占位符、禁令词、身份信息、缺图与缺源码、页数与摘要页）。
\end{enumerate}

\noindent\textbf{一键复现}：\texttt{python3 run\_all.py} 按上述顺序自动执行问题 1--4、
扫描、审计与台账；加 \texttt{--fast} 参数则跳过信息结构、场景数、VoI 与 CVaR 等对照实验。

\paragraph{实测耗时}

\begin{table}[htbp]
  \centering
  \caption{单机实测耗时（macOS 15.7.5 / arm64 / Python 3.13.12，与正文同一份代码与数据）}
  \label{tab:runtime}
  \footnotesize
  \begin{tabularx}{\linewidth}{L{4.0cm} L{2.0cm} Y}
    \toprule
    步骤 & 实测耗时 & 规模与说明 \\
    \midrule
    问题 1（典型日，确定性 LP） & 不足 0.1~s & 单日 144 槽；两次实测 0.014~s 与 0.008~s（纯求解时间） \\
    问题 2（固定价，两阶段） & 约 30~s & 334 天逐日两阶段子问题（$K=60$ 个场景） \\
    问题 3（固定价，四次预报） & 约 72~s & 四个预报时刻 $\times$ 334 天的滚动子问题 \\
    问题 4-2（波动电价，两阶段） & 约 50~s & 334 天逐日两阶段子问题（$K=64$ 个联合情景） \\
    问题 4-3（波动电价，四次预报） & 约 101~s & 四个预报时刻 $\times$ 334 天的滚动子问题 \\
    独立审计（A1--A10） & 约 13~s & 两次实测 12.9~s 与 16.7~s；其中交付物复算 3.3~s、逐日解复算 10.1~s \\
    \midrule
    合计（四问主链 $+$ 审计） & 约 4.4~min & 单机完成，无需集群或并行 \\
    \bottomrule
  \end{tabularx}
\end{table}

\noindent\textbf{输出文件}：\texttt{results/} 下的五个结果文件（\texttt{result1.xlsx}、
\texttt{result2.xlsx}、\texttt{result3.xlsx}、\texttt{result4-2.xlsx}、\texttt{result4-3.xlsx}）；
\texttt{figures/} 下的十三张矢量图（图 3--图 10 与 v2 新增的图 11--图 15）与两张 TikZ 源文件（图 1、图 2）；
\texttt{intermediate/} 下的审计报告、数值台账与各问中间结果。

\section{结果文件、论文表格与生成函数}\label{sec:app-filemap}

五个结果文件均按附件 5 模板的工作表名、列名与列序写出：\texttt{result1.xlsx} 为竖表
（一个时段一行，对应表~\bref{tab:q1-buy} 与表~\bref{tab:q1-soc}）；其余四个文件为横表
（一天一行，147 列 $=$ 日期 $+$ 144 个时段 $+$ 全天购电量 $+$ 全天购电费）。

\begin{table}[htbp]
  \centering
  \caption{结果文件 $\leftrightarrow$ 论文表格 $\leftrightarrow$ 生成函数 对应表}
  \label{tab:filemap}
  \footnotesize
  \setlength{\tabcolsep}{4pt}
  \begin{tabularx}{\linewidth}{L{3.4cm} L{2.4cm} L{4.0cm} Y}
    \toprule
    结果文件（工作表） & 论文位置 & 生成函数 & 行数与中间结果 \\
    \midrule
    result1.xlsx（计划购电量） & 表~\bref{tab:q1-buy}
      & \texttt{q1.solve\_q1} $\rightarrow$ \texttt{q1.write\_result1}
      & 144 行（1 行 1 个 10~min 槽）；\texttt{q1.json} \\
    result1.xlsx（充放电量） & 表~\bref{tab:q1-soc}
      & \texttt{q1.write\_result1}
      & 6 行（4~h 块）$+$ $0{:}00$ 与 $24{:}00$ 储电量；\texttt{q1.json} \\
    result2.xlsx（计划购电量） & 表~\ref{tab:q2-buy}、表~\bref{tab:q2-t1}、表~\bref{tab:summary}
      & \texttt{q2.run\_q2} $\rightarrow$ \texttt{results\_io.write\_result}
      & 334 行 $\times$ 147 列；\texttt{q2\_main.json} \\
    result2.xlsx（充放电量） & 表~\ref{tab:q2-soc}、表~\bref{tab:q2-t2} & 同上
      & 2\,004 行（334 天 $\times$ 6 个 4~h 块） \\
    result2.xlsx（紧急购电量） & 表~\ref{tab:q2-emergency}、表~\bref{tab:emergency} & 同上
      & 1\,156 行（缺口时段合并，同一日期可多行） \\
    result3.xlsx（计划购电量、调整购电量） & 表~\bref{tab:q3-epochs}（决策时刻与生效槽位）、表~\bref{tab:key4}
      & \texttt{q3.run\_q3} $\rightarrow$ \texttt{results\_io.write\_result}
      & 各 334 行 $\times$ 147 列；\texttt{q3\_n4.json} \\
    result3.xlsx（充放电量、紧急购电量） & 表~\bref{tab:key4}、表~\bref{tab:summary} & 同上
      & 2\,004 行 / 1\,190 行 \\
    result4-2.xlsx（计划购电量、充放电量、紧急购电量） & 表~\bref{tab:q4-summary}、表~\bref{tab:key4}、表~\bref{tab:summary}
      & \texttt{q4.run\_q4\_2} $\rightarrow$ \texttt{results\_io.write\_result}
      & 334 $\times$ 147、2\,004、1\,193 行；\texttt{q4\_2.json} \\
    result4-3.xlsx（计划购电量、调整购电量、充放电量、紧急购电量） & 表~\bref{tab:q4-summary}、表~\bref{tab:key4}
      & \texttt{q4.run\_q4\_3} $\rightarrow$ \texttt{results\_io.write\_result}
      & 334 $\times$ 147（两次）、2\,004、1\,187 行；\texttt{q4\_3.json} \\
    \bottomrule
  \end{tabularx}
\end{table}

表中「生成函数」为模块内的真实函数名：\texttt{q2.run\_q2}、\texttt{q3.run\_q3}、
\texttt{q4.run\_q4\_2}、\texttt{q4.run\_q4\_3} 完成求解，由各脚本的 \texttt{main()} 装配后
统一调用 \texttt{results\_io.write\_result} 写出工作簿；问题 1 的竖表由 \texttt{q1.solve\_q1} 求解、
\texttt{q1.write\_result1} 写出。论文其余图表不来自结果文件，其来源为：最优性证书取自
\texttt{q1\_dual.json} 与 \texttt{q1\_dual2.json}（\texttt{q1.py} 的对偶校验），
CVaR 前沿取自 \texttt{q4\_cvar.json}（\texttt{q4.py --mode cvar}，底层函数 \texttt{q4.solve\_cvar\_day}），
全局灵敏度取自 \texttt{exp\_sobol.json}，非预期性置换检验取自 \texttt{exp\_na.json}。

\section{独立审计报告}

审计由\textbf{非建模实现者}独立完成：脚本 \texttt{code/audit.py} 只读原始附件与本队写出的五个结果文件，
\textbf{不共用建模代码的费用函数与约束构造函数}，以保证「自证循环」不会发生。
机器可读结果在 \texttt{intermediate/audit.json}，可读报告在 \texttt{intermediate/audit\_report.md}。

\paragraph{审计口径与阈值}

\begin{itemize}
  \item \textbf{残差}：模型内部解（10~min 槽级、未取整）阈值 $10^{-6}$~kWh；
        交付物（结果文件只给 4~h 块级充放电量与 $0{:}00$、$24{:}00$ 两个储电量）的电量按
        $10^{-4}$~kWh 取整后回读，块级递推残差的上界为
        $0.5\times10^{-4}+0.9\times6\times0.5\times10^{-4}+6\times0.5\times10^{-4}/0.9+0.5\times10^{-4}=7.0\times10^{-4}$~kWh，
        故交付物的残差阈值取 $10^{-3}$~kWh
        （贴合交付精度的判定，不放松内部解的口径）；
  \item \textbf{越界数与违反数}：储电量越界、充放电功率越界、充放互斥违反、负电量，全部要求等于 $0$；
  \item \textbf{费用复算相对误差}：$\le10^{-4}$（费用只用结果文件与附件逐时段重算，不调用建模费用函数）。
\end{itemize}

\paragraph{逐项结果（正文口径）}

下表为正文第九章原先给出的\textbf{双列口径}审计表：左列为 10~min 级模型内部解，右列为结果文件
回读后的块级复算，便于对照“内部解”与“交付物”两层判定的差别。

\begin{table}[htbp]
  \centering
  \caption{独立审计项、阈值与实测结果（内部解：问题 1、2、3 的 10~min 级解；交付物：result1、result2、
           result3、result4-2、result4-3 共五个结果文件；窗口 2025-02-01 至 12-31）}
  \label{tab:audit}
  \scriptsize
  \setlength{\tabcolsep}{3pt}
  \begin{tabularx}{\linewidth}{C{0.62cm} L{3.95cm} L{3.35cm} L{3.35cm} C{0.95cm}}
    \toprule
    项 & 检查内容与阈值 & 内部解实测 & 交付物实测 & 判定 \\
    \midrule
    A1 & 功率/电量平衡残差 \(\le10^{-6}\)~kWh
      & 问题 1 \(1.49\times10^{-13}\)；问题 2、3 \(3.41\times10^{-13}\)~kWh
      & 交付物无逐槽 \(L,PV\) 列，平衡只能在内部解上核验（汇总列对拍见 A10） & PASS \\
    A2 & 储电量递推残差 \(\le10^{-6}\)~kWh
      & 问题 1 \(9.09\times10^{-13}\)；问题 2、3 \(1.09\times10^{-11}\)~kWh
      & 问题 1 \(3.64\times10^{-12}\)；问题 2、3、4-2、4-3 为
        \(2.96/2.72/2.11/2.61\times10^{-4}\)~kWh，为写入取整所致 & PASS\textsuperscript{注} \\
    A3 & 储电量边界：越界数 \(=0\)（区间 \([1200,10\,800]\)~kWh） & 0 & 0 & PASS \\
    A4 & 充放功率边界：内部解按单槽上限 \(833.3333\)~kWh；交付物按块级上限 \(20\,000\)~kWh & 0 & 0 & PASS \\
    A5 & 充放互斥：\(\min(c_t,d_t)\le10^{-6}\)
      & 违反槽数 0（充放同时槽数 0）
      & 交付物只有 4~h 聚合量，块内可同时有充放电，故块级不可检验 & PASS \\
    A6 & 周期/终端：问题 1 要求 \(S_0-S_{144}=0\)；问题 2--4 要求跨日衔接且 \(S_T\ge S_0\)
      & 问题 1 周期偏差 \(0.000\)；问题 2、3 末值低于初值的天数为 0
      & 五个结果文件末值低于初值的天数均为 0 & PASS \\
    A7 & 非负性：所有电量与费用 \(\ge-10^{-6}\) & 负购电槽数 0 & 负购电量数 0 & PASS \\
    A8 & 费用复算（计划 \(+\) 紧急 \(+\) 调整相关）相对误差 \(\le10^{-4}\)
      & 问题 1、2、3 的相对误差均为 0
      & 交付物层只能核验计划购电费（最大相对偏差 \(4.32\times10^{-7}\)，问题 4-2） & PASS \\
    A9 & 日期与结构：334 天、工作表名/列名/列序与模板一致
      & 334 天
      & 五个文件均为 334 行 \(\times\)147 列，日期范围 2025-02-01 至 12-31；紧急购电行数
        1156/1190/1193/1187 & PASS \\
    A10 & 全天汇总列与逐时段求和一致：相对误差 \(\le10^{-4}\) & ---
      & 全天购电量最大偏差 \(0.0012\)~kWh（问题 4-3）；全天购电费最大相对偏差
        \(4.32\times10^{-7}\)（问题 4-2） & PASS \\
    \bottomrule
  \end{tabularx}
  \\[1ex]
  {\footnotesize 表注：交付物 A2 的 \(10^{-4}\)~kWh 量级残差可由写入精度完全解释——结果写入对 6 个块级
  充放电量与 \(0{:}00\)/\(24{:}00\) 两个储电量按 4 位小数取整，取整误差上界为
  \(6\times5\times10^{-5}+2\times5\times10^{-5}=4\times10^{-4}\)~kWh，与实测值一致；该残差不是物理
  越界（A3、A4 越界数均为 0）。}
\end{table}

\paragraph{逐项结果}

\begin{table}[htbp]
  \centering
  \caption{独立审计 A1--A10（同时承载正文的 L1--L5 验证链条）：阈值、实测值与判定
           （「内」为 10~min 级模型内部解，「交」为结果文件块级）}
  \label{tab:audit-appendix}
  \footnotesize
  \setlength{\tabcolsep}{4pt}
  \begin{tabularx}{\linewidth}{C{0.9cm} L{3.4cm} L{2.9cm} Y C{1.2cm}}
    \toprule
    项 & 检查内容 & 阈值 & 实测值 & 判定 \\
    \midrule
    A1 & 电量平衡残差 & $\max\lvert E_t+PV_t\Delta t+d_t+g_t-L_t\Delta t-c_t-s_t\rvert\le10^{-6}$~kWh
       & 内：问题 1 为 $1.49\times10^{-13}$，问题 2、3 均为 $3.41\times10^{-13}$ & PASS \\
    A2 & 储电量递推残差 & 内 $\le10^{-6}$~kWh；交 $\le10^{-3}$~kWh
       & 内：$9.09\times10^{-13}$（问题 1）、$1.09\times10^{-11}$（问题 2、3）；交（问题 1、2、3、4-2、4-3）：$3.64\times10^{-12}$、$2.96\times10^{-4}$、$2.72\times10^{-4}$、$2.11\times10^{-4}$、$2.61\times10^{-4}$ & PASS \\
    A3 & 储电量越界数 & 等于 $0$（运行区间 $[1200,\,10800]$~kWh）
       & 交：五个结果文件均为 $0$；内：问题 1--3 均为 $0$ & PASS \\
    A4 & 充放电功率越界数 & 等于 $0$（单槽上限 $\overline P\Delta t=833.3333$~kWh）
       & 交：块级上限 $4\overline P\Delta t=3333.3333$~kWh，四个横表均为 $0$；内：问题 1--3 的 10~min 级均为 $0$ & PASS \\
    A5 & 充放互斥违反数 & 等于 $0$（$\eta_c\eta_d=0.81<1$ 时最优解不同时充放）
       & 内：问题 1--3 均为 $0$，并在 97 个储电量双界均不活跃的槽位复核 & PASS \\
    A6 & 周期与终端条件 & 问题 1 的 $\lvert S_0-S_{144}\rvert=0$；问题 2--4 每日 $S_{144}\ge S_0$
       & 问题 1 的日周期偏差 $0.0$~kWh；四个横表「末值低于初值的天数」均为 $0$ & PASS \\
    A7 & 非负性违反数 & 等于 $0$
       & 交：负购电量 $0$ 个（四个横表各 $334\times144$ 个时段值，加 result1 的 144 个）；内：问题 1--3 均为 $0$ & PASS \\
    A8 & 费用复算相对误差 & $\le10^{-4}$
       & 交：问题 1 为 $0.0$；内：问题 1--3 均为 $0.0$ & PASS \\
    A9 & 日期与结构 & 334 天；工作表名、列数与列序逐项一致
       & 四个横表均为 $334\times147$，日期 2025-02-01 至 2025-12-31；result1 为 144 行计划表加 6 行充放电表；工作表名 result1 为计划购电量与充放电量，result2、result4-2 为计划购电量、充放电量、紧急购电量，result3、result4-3 为计划购电量、调整购电量、充放电量、紧急购电量 & PASS \\
    A10 & 全天汇总列 & 购电量偏差 $\le10^{-2}$~kWh（144 槽取整上界 $7.2\times10^{-3}$）；费用相对偏差 $\le10^{-4}$
       & 购电量最大偏差 $1.10\times10^{-3}$、$1.10\times10^{-3}$、$1.00\times10^{-3}$、$1.20\times10^{-3}$~kWh（对应问题 2、3、4-2、4-3）；费用最大相对偏差 $3.40\times10^{-7}$、$3.40\times10^{-7}$、$4.32\times10^{-7}$、$4.05\times10^{-7}$ & PASS \\
    \bottomrule
  \end{tabularx}
\end{table}

\paragraph{反面证据：结算平衡不得漏掉紧急购电项}

附件 5 把「计划购电量」与「紧急购电量」分成\textbf{两张表}，因此计划量不等于实际供给量。
若 A1 的结算平衡漏掉紧急购电项 $g_t$（即把计划购电量当作全部供给），残差恰好等于 $g_t$：
在问题 2 的 334 天逐槽解上 $\max_t g_t=326.18$~kWh，在问题 3 上为 $204.87$~kWh，
分别比 $10^{-6}$~kWh 的容差大 $3.3\times10^{8}$ 倍与 $2.0\times10^{8}$ 倍；
补上 $+g_t$ 后，问题 2、3 的平衡残差降为 $3.41\times10^{-13}$~kWh。
同一次审计还给出：问题 2 的盈余处置量（弃光与超购部分的处置量）为 $3\,106\,390.84$~kWh，
问题 3 为 $2\,638\,963.35$~kWh，均为正残差，与「供给不低于需求」的口径一致。

\paragraph{审计结论}

\begin{center}
\fbox{\ 全部通过，可进入写作\ }\qquad
\fbox{\ 存在失败项，已定位并修正\ }（未勾选）
\end{center}

\noindent A1--A10 十项全部通过；审计同时复核了交付物的结构与全天汇总列，
未发现任何越界、违反非负性或费用复算不一致。

\section{源程序清单}\label{sec:app-code}

以下为本论文建模、求解、绘图、审计与数值汇总所用的\textbf{全部源程序}（格式规范第五条），
与 \texttt{code/} 下的实际文件逐项一致；附录用 \texttt{lstinputlisting} 原样引入文件内容。
\texttt{[plain]} 选项表示该脚本不参与求解，只做绘图或文本处理。

\codefile{config.py}
\codefile{lp_engine.py}
\codefile{q1.py}
\codefile{q2.py}
\codefile{q3.py}
\codefile{q4.py}
\codefile{results_io.py}
\codefile{exp_scan.py}
\codefile{exp_sobol.py}
\codefile{exp_na.py}
\codefile{exp_oos.py}
\codefile{exp_benchmark.py}
\codefile{audit.py}
\codefile{paper_numbers.py}
\codefile{run_all.py}
\codefile{requirements.txt}
\codefile[plain]{plot_figures.py}
\codefile[plain]{plot_figures_v2.py}
\codefile{copula.py}
\codefile{scen_copula.py}
\codefile{exp_copula.py}
\codefile{exp_chance.py}
\codefile{exp_dro.py}
\codefile{decomp.py}

\begin{guideblock}
附录放\textbf{可复制的纯文本源码}，\textbf{不要用截图}，也不要放空壳函数。
每张数据图都要能追到生成脚本：图 3--图 10 由 \texttt{plot\_figures.py} 一次生成（每张图一个函数），
图 1、图 2 为 \texttt{figures/} 下的 TikZ 源文件、由 XeLaTeX 直接排版。
新增脚本必须同步补入本清单，否则「论文未列全部源程序」属格式规范第五条风险。
\end{guideblock}

\subsection{v2 新增实验的中间结果与复算入口}

v2 在原有主模型之外新增四组对照实验与两组拟合，全部数字同样只由脚本产生；
下表给出每个中间结果文件的复算入口与在正文中的引用位置。

\begin{table}[htbp]
  \centering
  \caption{v2 新增中间结果文件的复算入口（全部位于 \texttt{intermediate/} 目录）}
  \label{tab:v2map}
  \small
  \begin{tabular}{lll}
    \toprule
    中间结果文件 & 生成脚本 & 正文引用位置 \\
    \midrule
    \texttt{exp\_copula.json} & \texttt{exp\_copula.py}
      & \bref{sec:vine}（C 藤选族）、\ref{sec:q4-gen}（生成器对照）、\ref{sec:q4-risk}（分时段风险画像） \\
    \texttt{exp\_chance.json} & \texttt{exp\_chance.py} & \bref{sec:q2-chance}（机会约束对照与报童互证） \\
    \texttt{exp\_dro.json} & \texttt{exp\_dro.py} & \ref{sec:q4-dro}（盒式分布鲁棒的保守度—费用曲线） \\
    \texttt{decomp.json} & \texttt{decomp.py} & \ref{sec:lshape}（L 形分解与直接求解对拍） \\
    \bottomrule
  \end{tabular}
\end{table}

四组实验均复用 \texttt{config.py} 的口径（144 槽、$DT=1/6$~h、结果窗口 334 天、1 月预热、
终端 $S_{144}\ge S_0$、不允许售电），并各自包含独立自检：\texttt{exp\_dro.json} 的
\texttt{audit} 字段给出计划可行性、功率界、储电量界、递推残差、充放互斥与最坏情形覆盖的一致性残差；
\texttt{decomp.json} 的 \texttt{equivalence\_K60} 给出 334 天逐日的两算法目标值相对间隙；
\texttt{exp\_chance.json} 的 \texttt{cross\_check} 把两阶段随机规划的复算值与 \texttt{q2\_main.json}
对拍（相对差 $0.00$）；\texttt{exp\_copula.json} 的 \texttt{part3\_generators} 给出四种生成器的
样本外决策指标与边际保真度。

\section{数据字典}

\begin{table}[htbp]
  \centering
  \caption{附件字段、单位与时标约定}
  \label{tab:datadict}
  \footnotesize
  \begin{tabularx}{\linewidth}{C{1.4cm} L{3.0cm} C{1.6cm} C{2.4cm} Y}
    \toprule
    附件 & 字段 & 单位 & 时间范围 & 时标约定 \\
    \midrule
    附件 1 & 电价 & 元/kWh & 一天（144 时段） & 标签 = 时段末端；$0{:}00{+}1$ 表示次日 $0{:}00$（当日 $24{:}00$） \\
    附件 1 & 小区负载 & kW & 同上 & 同上，\textbf{实为附件 2 全年逐时段均值}（见下） \\
    附件 1 & 光伏发电「预测」功率 & kW & 同上 & \textbf{实为附件 2 全年逐时段均值，不是预报}（见下） \\
    附件 2 & 小区负载 & kW & 2025-01-01 至 12-31 & 365 行 $\times$ 144 列，标签 = 时段末端 \\
    附件 2 & 光伏发电实际功率 & kW & 同上 & 同上 \\
    附件 3 & 光伏发电预报 & kW & 同上 & 每天 $0/6/12/18$ 时发布，预报 $k$ 小时对应 $(h_0+k-1,\ h_0+k]$ \\
    附件 4 & 电价 & 元/kWh & 同上 & 365 行 $\times$ 144 列，标签 = 时段末端 \\
    附件 5 & 结果文件模板 & --- & 同上 & 工作表名、列名与列序即本文结果文件的格式约定 \\
    \bottomrule
  \end{tabularx}
\end{table}

\paragraph{附件 1 的身份核查（数据字典的关键一条）}

附件 1 的\textbf{电价、小区负载、光伏发电「预测」功率三列，实为附件 2 与附件 4 的全年逐时段均值，
并不是日前预报}：以 365 天逐时段均值逐槽比对，电价的最大偏差为 $5.15\times10^{-5}$~元/kWh
（\texttt{intermediate/q4\_price.json} 的 \texttt{skeleton\_is\_mean} 字段），
负载的最大偏差为 $1\times10^{-4}$~kW（\texttt{paper\_numbers.json}），光伏的最大偏差为 0.0431~kW。
因此附件 1 \textbf{只能用于问题 1 与全文基线口径}：问题 1 的典型日与无储能基线 $C^{base}$ 用它，
问题 2--4 的日前计划一律改用附件 3 的当次预报（问题 2 用 $0{:}00$ 预报，问题 3 在
$0{:}00$、$6{:}00$、$12{:}00$、$18{:}00$ 四次更新），结算一律用附件 2 与附件 4 的实际值。
若把附件 1 当作日前预报，计划就脱离了该时刻真正可得的信息，属信息结构的口径错误。
另有一个窗口口径需同时说明：只用结果窗口 2025-02-01 至 12-31 的 334 天求均值时，
电价的最大偏差约 $3.1\times10^{-2}$~元/kWh，即\textbf{窗口内均值并不完全等于附件 1}，
故问题 2--4 也不能把附件 1 当作窗口内的均值曲线使用。

（符号说明表已移入正文第 2.3 节，见正文表~\bref{tab:notation}。）


\subsection{相依结构：C 藤 pair-Copula 的逐层选族与 AIC 对照}\label{sec:vine-app}

正文 \S3.7 给出三点读数与结论；本小节给出逐层选族结果（\ref{tab:vine}）与五族候选的完整 AIC
对照（\ref{tab:vine-cand}），供复核。

\begin{table}[htbp]
  \centering
  \caption{C 藤 pair-Copula 逐层选族结果（样本为 2025-02-01 至 12-31 的 334 个日水平观测；
           AIC 越小越好，经验距离为与经验 Copula 的 \(L_2\) 距离）}
  \label{tab:vine}
  \footnotesize
  \setlength{\tabcolsep}{4pt}
  \begin{tabularx}{\linewidth}{L{2.4cm} C{1.2cm} L{2.55cm} C{1.15cm} C{1.15cm} C{1.15cm} Y}
    \toprule
    变量对 & 最优族 & 参数 & \(\tau\) & \(\lambda_U\) & \(\lambda_L\) & AIC / 经验距离 \\
    \midrule
    电价—负载 & Frank & \(\theta=10.4686\) & \(+0.6976\) & 0 & 0 & \(-453.54\) / \(0.0135\) \\
    电价—光伏 & \(t\) & \(\rho=-0.2850,\ \nu=3.2324\) & \(-0.2383\) & \(0.0479\) & \(0.0479\)
      & \(-44.50\) / \(0.0210\) \\
    负载—光伏 \(|\) 电价 & Frank & \(\theta=6.3107\) & \(+0.5542\) & 0 & 0 & \(-253.44\) / \(0.0281\) \\
    \bottomrule
  \end{tabularx}
\end{table}

\textbf{三点读数。} ① \textbf{电价与负载强正相依}（\(\tau=+0.698\)、Pearson \(+0.919\)），
但两端\textbf{无尾相依}（\(\lambda_U=\lambda_L=0\)）：贵的时候确实多用电，但这种关系在尾部并不额外
放大。② \textbf{电价与光伏为弱负相依并带轻微对称尾相依}：最优族是自由度仅 3.23 的 \(t\) Copula，
\(\tau=-0.238\)、\(\lambda_U=\lambda_L=0.0479\)。这正是“光伏极低的日子电价极高”这类\textbf{尾部
共现}的量化结果——它存在，但强度只有约 5\%，既不能忽略也不应夸大；若改用 Gaussian 族（AIC 高 14.89，
且强制 \(\lambda=0\)）就会把它抹平。③ \textbf{控制电价后负载与光伏仍强正相依}（\(\tau=+0.554\)），
说明两者有共同驱动（季节与天气），并非仅是电价的代理变量。\evid{C-Q4-02}

五族候选的 AIC 对照（表~\ref{tab:vine-cand}）显示：三对变量的 AIC 最优族与经验 Copula 距离最优族
排序一致，且最优与次优的 AIC 差距（51.37、14.89、37.29）远大于拟合噪声，选族结论不依赖单一判据。

\begin{table}[htbp]
  \centering
  \caption{五族候选 Copula 的 AIC 对照（同一 334 个日水平观测；加粗为各族中的最小值，
           “\(|\)”后为条件变量）}
  \label{tab:vine-cand}
  \small
  \begin{tabular}{lccccc}
    \toprule
    变量对 & Gaussian & \(t\) & Clayton & Gumbel & Frank \\
    \midrule
    电价—负载 & \(-402.17\) & \(-398.97\) & \(-357.13\) & \(-305.67\) & $\mathbf{-453.54}$ \\
    电价—光伏 & \(-29.61\) & $\mathbf{-44.50}$ & $2.16$ & $2.22$ & \(-22.75\) \\
    负载—光伏 \(|\) 电价 & \(-216.15\) & \(-212.55\) & \(-153.59\) & \(-173.89\) & $\mathbf{-253.44}$ \\
    \bottomrule
  \end{tabular}
\end{table}


五族 AIC 对照见表~\ref{tab:vine-cand}。


\subsection{相关性三尺度与预报误差分档明细}\label{sec:datadetail-app}

正文 \S3.4 与 \S3.5 给出结论；本小节给出完整表格，供复核。

\begin{table}[htbp]
  \centering
  \caption{电价与负载、净负荷、光伏的相关系数（三个尺度；池化与去形状口径的样本量为
           \(365\times144=52\,560\) 个时段观测，日级口径为 365 个日观测）}
  \label{tab:corr-app}
  \small
  \begin{tabular}{lccc}
    \toprule
    变量对 & 池化（全部时段） & 日级（逐日先算再合并） & 剔除日内形状后 \\
    \midrule
    电价 --- 负载 & \(+0.5077\) & \(+0.8977\) & \(+0.4554\) \\
    电价 --- 净负荷 \(L-PV\) & \(+0.2602\) & \(+0.9824\) & \(+0.1583\) \\
    电价 --- 光伏 & \(+0.0284\) & \(-0.3431\) & \(+0.0484\) \\
    \bottomrule
  \end{tabular}
\end{table}

\begin{table}[htbp]
  \centering
  \caption{光伏预报误差统计（样本为 2025 年每天 \(0{:}00\) 发布的 24~h 预报与附件 2 对应实际功率的
            逐槽配对，\(n=52\,560\)；MAE、RMSE、bias 单位均为 kW）}
  \label{tab:fcsterr-app}
  \footnotesize
  \setlength{\tabcolsep}{4pt}
  \begin{tabularx}{\linewidth}{L{3.1cm} C{1.5cm} C{1.5cm} C{1.5cm} C{1.5cm} Y}
    \toprule
    口径（目标整点范围） & 样本量 \(n\) & MAE & RMSE & bias & 结论 \\
    \midrule
    全时段（\(0\)--\(23\) 时，144 槽/天） & \(52\,560\) & \(361.58\) & \(636.37\) & \(-26.96\)
      & 夜间近零样本摊薄指标，只作对照 \\
    白天（\(\ge5\) 时，114 槽/天） & \(41\,610\) & \(456.72\) & \(715.22\) & \(-34.06\)
      & \textbf{建模依据}：MAE 为全时段的 1.26 倍 \\
    \midrule
    \multicolumn{6}{l}{\textbf{按目标整点聚合的偏差结构}（组内 bias 单位 kW）} \\
    \midrule
    夜间（光伏无有效出力） & \(0\)--\(4\) 时 & \multicolumn{2}{c}{\(0.00\) 至 \(+0.09\)}
      & \(4\) 时 \(+0.09\) & 无信息、无惩罚 \\
    上午（系统性\textbf{高估}） & \(5\)--\(10\) 时 & \multicolumn{2}{c}{\(+214.21\) 至 \(+733.00\)}
      & \(7\) 时 \(+733.00\) & 计划偏少 $\rightarrow$ 触发 \(5\) 倍紧急购电 \\
    正午过渡 & \(11\) 时 & \multicolumn{2}{c}{\(+7.21\)} & --- & 接近无偏 \\
    下午（系统性\textbf{低估}） & \(12\)--\(18\) 时 & \multicolumn{2}{c}{\(-801.87\) 至 \(-106.46\)}
      & \(15\) 时 \(-801.87\) & 计划偏多 $\rightarrow$ 白付电费或付违约电价 \\
    傍晚之后（预报空洞） & \(19\)--\(23\) 时 & \multicolumn{2}{c}{\(-0.15\) 至 \(0.00\)}
      & \(19\) 时 \(-0.15\) & 出力近零，不触发惩罚 \\
    \bottomrule
  \end{tabularx}
\end{table}

\subsection{联合情景保真度检验明细}\label{sec:scenfid-app}

正文 \S7.2 给出结论；本小节给出完整检验表。

\begin{table}[htbp]
  \centering
  \caption{联合情景的保真度检验（对象：64 个代表历史日与 64 个分块 Bootstrap 情景；参照：2025-02-01 至
           12-31 的 334 天实际数据；相对偏差单位 \%，相关系数为无量纲量）}
  \label{tab:scenfid}
  \footnotesize
  \setlength{\tabcolsep}{4pt}
  \begin{tabularx}{\linewidth}{L{4.35cm} C{2.35cm} C{2.35cm} Y}
    \toprule
    检验指标 & 代表历史日（主口径） & 分块 Bootstrap（对照） & 判读 \\
    \midrule
    电价均值相对偏差 & 0.08\% & 3.41\% & 主口径的均值几乎无偏 \\
    最大分位相对差（7 个分位点） & 6.99\%（1\% 分位） & 20.4\%（1\% 分位）
      & 两者都在低价尾部误差最大，对照口径的尾部偏差约为主口径的 3 倍 \\
    情景 \(\rho(p,\text{净负荷})\) & 0.2597（实际 0.2620） & 无法按日配对
      & 主口径 \textbf{保留}了联合结构 \\
    日净需求残差均值 & 保持 \(-715.85\)~kWh（全样本 \(-712.23\)） & 不适用（块级拼装）
      & 净需求侧边缘分布同样保真 \\
    \bottomrule
  \end{tabularx}
\end{table}
\section{对照实验明细}\label{sec:ctl-app}

正文 \S7.3 的三组对照实验只保留结论；本节给出其完整表格、逐条数据与判定阈值，便于复核。
三个小节沿用正文的小节标题与编号，便于互查。

\subsection{情景生成方式的样本外对照：结论对生成器的敏感性}
\label{sec:q4-gen}

\bref{sec:q4-price} 的联合情景用“二维分层取代表历史日”构造。这一做法是否只是“拟合了样本内的 334 天”？
换一种生成方式，结论会不会变？本节用\textbf{样本外对照实验}回答：情景集只用 \textbf{2025-02-01 至 06-30}
（150 天）拟合，在 \textbf{07-01 至 12-31}（184 天）上按真实轨迹结算，四种生成器共用同一决策模型、
同一价格信息结构（假设 A10）与同一终端条件：\textbf{（1）历史代表日}（主方案，二维分层 \(+\) 层内
medoid）；\textbf{（2）分块 Bootstrap}（6~h 块按时刻对齐重采样）；\textbf{（3）KDE \(+\) Copula}
（日水平核密度边缘 \(+\) 最优二元 Copula，抽样 600 次后削减到与（1）相同的情景数）；
\textbf{（4）按光伏规模缩放}（电价情景同（1），残差情景改为“训练窗相对残差形状 \(\times\) 当日光伏
预报电量”）。

\begin{table}[htbp]
  \centering
  \caption{四种情景生成器的样本外对照（训练窗 2025-02-01 至 06-30 共 150 天，测试窗 07-01 至 12-31 共 184 天；
           同一决策模型、同一实际轨迹；费用单位元，电量为 184 天累计 kWh，“边际偏差”为日总残差
           1\%/99\% 分位相对训练窗实际分布的最大绝对偏差）}
  \label{tab:gen-cmp}
  \small
  \begin{tabular}{lrrrrr}
    \toprule
    生成器 & 总费用 & 相对（1） & 紧急购电量 & 逐日 CVaR$_{0.95}$ & 边际偏差 \\
    \midrule
    （1）历史代表日 & $9\,046\,598.58$ & — & $76\,033.90$ & $79\,378.60$ & $345.7$ \\
    （2）分块 Bootstrap & $9\,038\,928.08$ & $-0.085\%$ & $72\,324.35$ & $79\,381.24$ & $2\,102.3$ \\
    （3）KDE $+$ Copula & $9\,072\,284.06$ & $+0.284\%$ & $107\,458.94$ & $80\,061.08$ & $1\,519.1$ \\
    （4）按光伏规模缩放 & $9\,068\,058.80$ & $+0.237\%$ & $92\,727.06$ & $81\,776.72$ & 不适用 \\
    \bottomrule
  \end{tabular}
\end{table}

\textbf{三条结论。} ① \textbf{费用结论对生成器不敏感}：四种生成器的样本外总费用极差
33\,355.98 元，仅为四者均值（9\,056\,467.38 元）的 0.37\%，相对无储能基线的节省率落在
11.10\%--11.43\%。这比“样本外劣化”检验更强：不仅主方案不过拟合，连换一种生成范式也不改变结论。
② \textbf{但分项指标并不等价}：紧急购电量在四种生成器下为 72\,324--107\,459~kWh（相差 48.5\%）、
逐日 CVaR\(_{0.95}\) 为 79\,379--81\,777~元（相差 3.0\%）——\textbf{费用稳健，而“靠多少紧急购电兜底”
可被生成方式改变}，引用该分项必须同时给出情景口径。③ \textbf{联合尾部的保真度上 Copula 更好}：
“高价 \(+\) 光伏高估”的联合尾共现率真值 0.0533，Copula 生成器 0.0500、历史代表日 0.0667，分块
Bootstrap 只有 0.0167；这与 \bref{sec:vine} 一致——尾部共现在数据中\textbf{本来就弱}（\(\lambda\approx0.048\)），
其保真度差异不足以传导到决策层。据此本文\textbf{主模型保留历史代表日生成器}，\textbf{Copula 用于相依
结构诊断}（\bref{sec:vine}）与\textbf{敏感性对照}（本小节），两者都不改变任何主结果。\evid{C-Q4-03}



\subsection{分时段 CVaR 风险画像}
\label{sec:q4-risk}

“全年总费用”只有一个数，掩盖了风险的时段结构（分时段风险度量在电力市场风险管理中已有应用~\cite{refwangtaiqiang}）。本节把费用\textbf{按槽位分解}：
第 $d$ 天第 $t$ 槽的费用为 $c_{d,t}=p_{d,t}E_{d,t}+5\,p_{d,t}g_{d,t}$（计划购电 $+$ 紧急购电），
对固定的 $t$ 取跨 334 天的 $\mathrm{CVaR}_{0.95}$，再按小时求和，得到\textbf{逐小时风险画像}
（图~\ref{fig:risk-profile}；问题 2 与问题 4-2 分别计算）。

\begin{figure}[htbp]
  \centering
  \includefig[width=0.68\linewidth]{fig13_risk_profile.pdf}
  \caption{分时段 CVaR 风险画像（柱：各槽位费用的 \(\mathrm{CVaR}_{0.95}\) 按小时求和；实线：逐小时均值；
            虚线：槽位缺电频率（右轴）；左：问题 2，右：问题 4-2；范围：334 天 \(\times\) 144 槽）}
  \label{fig:risk-profile}
\end{figure}

\textbf{风险的时段分化是结构性的。} ① \textbf{费用尾部集中在傍晚}：固定价下逐小时
\(\mathrm{CVaR}_{0.95}\) 的高点出现在 \(15\)--\(18\) 时（6\,075、5\,566、6\,604、6\,358 元/日），
波动价下同段升到 7\,492、6\,778、7\,550、8\,098 元/日。② \textbf{缺电尾部集中在上午}：槽位缺电频率
在 \(5\)--\(11\) 时高达 14.5\%--23.8\%（固定价）与 14.8\%--34.6\%（波动价），峰值分别在 \(11\)
时（23.8\%）与 \(5\) 时（34.6\%），而 \(0\)--\(3\) 时与 \(19\)--\(23\) 时几乎为 0——这与
\S3.5 的偏差剖面完全对应（上午光伏被系统性高估使 \(0{:}00\) 计划偏少）。③ \textbf{风险集中度}：
\(\mathrm{CVaR}\) 最高的 12 个槽位只占全部槽位之和的 16.73\%（问题 2）与 16.10\%（问题 4-2），
说明风险不是集中在个别槽位而是沿上午与傍晚两个时段带分布；逐槽 \(\mathrm{CVaR}\) 之和是逐槽均值之和的
2.01 倍（固定价）与 2.34 倍（波动价），即\textbf{电价波动把费用的尾部放大得比均值多}。\evid{C-Q4-04}

\textbf{决策含义。} 安全边际应当\textbf{按目标整点}配置（上午与傍晚需要，\(0\)--\(4\) 时不需要），
而不是全天统一加比例；滚动调整的价值同样集中在这两个时段带。


\subsection{分布鲁棒对照：保守度—费用曲线}
\label{sec:q4-dro}

\ref{sec:q4-gen} 的四种生成器都属\textbf{随机规划}：先估计一个分布，再在期望意义下优化。当样本有限或分布可能
误设时，另一种思路是\textbf{分布鲁棒优化}（DRO）——在模糊集上取最坏期望，其最简单的成员就是
\textbf{盒式不确定集}（矩约束模糊集下的多阶段建模见文献~\cite{refyuhaodong}）。本节把它作为
\textbf{保守度对照}：把“保守度”参数化后定量回答一个问题——\textbf{用多少实际费用，能买到多少最坏
情形下的改善？}

\textbf{模型。} 取逐槽盒式不确定集
\(p_t\in[\bar p_t-k\sigma^p_t,\ \bar p_t+k\sigma^p_t]\)、
\(\widetilde{ND}_t\in[ND^{\text{fcst}}_t-k\sigma^{nd}_t,\ ND^{\text{fcst}}_t+k\sigma^{nd}_t]\)，其中 \(\sigma^p_t\) 与 \(\sigma^{nd}_t\) 分别为附件 4 电价的逐槽样本
标准差与 \(0{:}00\) 预报残差的逐槽样本标准差（365 天，\(\mathrm{ddof}=1\)）。目标在盒的顶点取得，最坏情形目标与约束取
\(\min\ \sum_t p^{\text{hi}}_tE_t+5\sum_t p^{\text{hi}}_t z_t\)、\(z_t\ge ND^{\text{hi}}_t+c_t-d_t-E_t\)、\(z_t\ge0\)，外加计划可行性与储电量约束（同 \bref{sec:q4-price}）；\(k=0\) 时退化为\textbf{按预报的确定性购电}（零安全边际）。

\textbf{一个结构性推论（先于数值给出）。} 上述最坏情形目标中 $E_t$ 与 $z_t$ 在同槽位上的目标系数之比
恒为 $1:5$，而把 $E_t$ 提高 $\delta$ 可使 $z_t$ 下降同量、目标值下降 $4p^{\text{hi}}_t\delta>0$，
故只要 $z_t>0$ 就仍有改进空间：\textbf{最优解必取 $z_t\equiv0$}，即
$E_t=\max\{0,\ ND^{\text{hi}}_t+c_t-d_t\}$。也就是说，\textbf{盒式 DRO 在本问题设定下退化为
“按最坏参数做确定性优化”}——它逐槽覆盖 $k$ 倍标准差的不利偏差，而不依赖任何追索机制。
数值上 $z$ 的一致性残差为 $0.000$（表~\ref{tab:audit}），与推论一致。

\begin{table}[htbp]
  \centering
  \caption{盒式分布鲁棒对照：保守度 $k$ 与实际回测表现（334 天，2025-02-01 至 12-31；决策用盒式鲁棒模型，
           结算用附件 4 实际电价与附件 2 实际负载/光伏；$W_{k=3}$ 为各方案在\textbf{同一}
           $k=3$ 不确定集下的最坏情形目标值；费用单位元，电量为 kWh）}
  \label{tab:dro}
  \small
  \begin{tabular}{crrrrrr}
    \toprule
    $k$ & 回测总费用 & 节省率 & 紧急购电量 & 逐日 CVaR$_{0.95}$ & 最坏单日 & $W_{k=3}$（万元） \\
    \midrule
    $0.0$ & $18\,258\,473.08$ & $-7.23\%$ & $1\,156\,618.02$ & $91\,630.53$ & $101\,653.74$ & $5\,956.86$ \\
    $0.5$ & $17\,404\,407.87$ & $-2.22\%$ & $793\,522.04$ & $86\,473.81$ & $94\,613.68$ & $3\,557.47$ \\
    $1.0$ & $16\,936\,297.51$ & $+0.53\%$ & $521\,080.78$ & $84\,250.61$ & $89\,459.18$ & $3\,058.17$ \\
    $\mathbf{1.5}$ & $\mathbf{16\,819\,102.59}$ & $\mathbf{+1.22\%}$ & $328\,147.59$ & $84\,471.07$ & $87\,672.62$ & $2\,544.50$ \\
    $3.0$ & $18\,015\,766.37$ & $-5.81\%$ & $62\,776.53$ & $91\,718.61$ & $95\,526.83$ & $2\,512.55$ \\
    \midrule
    \textbf{随机规划}（\bref{sec:q4-price}、\bref{sec:q4-frontier}） & $\mathbf{15\,094\,255.84}$ & $\mathbf{+11.35\%}$ & $182\,759.86$ & $\mathbf{77\,697.98}$ & $\mathbf{81\,643.86}$ & — \\
    \bottomrule
  \end{tabular}
\end{table}

\textbf{四条结论。} ① \textbf{回测费用对 $k$ 呈 U 形}：$k$ 由 $0$ 增大到 $1.5$ 时，多买电的代价
小于少买电触发的 $5$ 倍紧急购电，费用由 $18\,258\,473.08$ 降到 $16\,819\,102.59$~元（省
$1\,439\,370.50$ 元）；此后计划购电量继续上升而紧急购电已接近清零，费用反弹，$k=3.0$ 时回到
$18\,015\,766.37$~元。实证最优点为 $k^{*}=1.5$。② \textbf{盒式鲁棒在最坏情形上确实有效}：在
\textbf{同一} $k=3$ 的不确定集下评估，$k=0$ 的方案最坏情形要花 $5\,956.86$ 万元，$k=3$ 的方案只要
$2\,512.55$ 万元（降 $57.8\%$）。③ \textbf{但它的“保守”买贵了}：从 $k^{*}$ 加到 $k=3.0$，实际回测
费用多付 $1\,196\,663.78$~元（$+7.12\%$），即最坏情形费用每降低 $1\%$、回测费用要多付 $3.98$ 万元
（约基线的 $0.234\%$）。④ \textbf{盒式鲁棒被随机规划全面占优}：与 \bref{sec:q4-frontier} 的随机规划相比，最优方案
（$k^{*}=1.5$）在\textbf{全部五项指标}上更差——总费用高 $1\,724\,846.75$~元（$+11.43\%$）、逐日
$\mathrm{CVaR}_{0.95}$ 高 $6\,773.09$~元/日（$+8.72\%$）、最坏单日高 $6\,028.76$~元（$+7.38\%$）、
紧急购电量多 $145\,387.73$~kWh（$+79.6\%$）、槽位缺电频率反而更高（$10.24\%$ 对 $9.04\%$）。
\evid{C-Q4-05}

\textbf{为什么会这样，以及正确的使用方式。} 盒式不确定集把 $144$ 个槽位\textbf{同时}取到极端
（电价与净需求全线 $+k\sigma$），而实测 334 天里槽间残差并不同步（\S7.1 的日内自相关与 $19$--$23$ 时
近乎零的偏差都说明这一点），因此这个“最坏情形”几乎不可能发生，为它买保险的钱花得不值。此外 $k$ 的
经验最优值恰好使\textbf{电价加权的缺货频率}降到 $1/5$ 附近，与 \bref{sec:q2-newsvendor} 报童模型的允许缺货概率
$\beta^{*}=1/5=0.20$ 一致——两条完全不同的建模路线给出同一个服务水平。故盒式鲁棒的价值是
\textbf{压力上界}与\textbf{轻度鲁棒中间方案}（$k=0.5$--$1.0$），而非“更好的策略”；真正改进的方向是把
\textbf{全盒}换成\textbf{带不确定性预算 $\Gamma$ 的不确定集}或\textbf{矩约束模糊集}
\cite{refyuhaodong}，使“每个时段都可能极端、但不允许全部时段同时极端”这一先验进入模型（\S10.4）。



\section{口径定义汇总}\label{sec:caliber-app}

正文各章在本题参数范围内用到时以~\calref{C-0x} 引用本节的对应条目。口径编号在建模前固定、未随结果调整；
每个口径含定义、公式、数据来源、时段、单位、不含项、对照与复算入口八项。

g

g

g

g

g

\section{长推导}

\subsection{对偶推导与阈值律}

记式~\bref{eq:q1-balance} 的乘子为 $\alpha_t$（电量平衡的影子价格，
在正常购电时段等于电价，见式~\eqref{eq:app-alpha-eq-p}），式~\bref{eq:q1-soc} 的乘子为
$\nu_t$（储电量的边际价值，即水价值），储电量上下界乘子为 $\lambda^{\pm}_t\ge 0$，
充放电功率界乘子为 $\varrho^{c}_t,\varrho^{d}_t\ge 0$。
把电量平衡写成 $(L_t-PV_t)\Delta t+c_t+s_t-E_t-d_t=0$，拉格朗日函数为

\begin{equation}
  \begin{aligned}
    \mathcal L=\ &\sum_{t}p_tE_t
    +\sum_{t}\alpha_t\Bigl[(L_t-PV_t)\Delta t+c_t+s_t-E_t-d_t\Bigr]\\
    &+\sum_{t}\nu_t\Bigl[S_t-S_{t-1}-\eta_cc_t+\frac{d_t}{\eta_d}\Bigr]
    +\sum_{t}\lambda^{+}_t\bigl[S_t-\overline S\bigr]
    +\sum_{t}\lambda^{-}_t\bigl[\underline S-S_t\bigr]\\
    &+\sum_{t}\varrho^{c}_t\bigl[c_t-\overline P\Delta t\bigr]
    +\sum_{t}\varrho^{d}_t\bigl[d_t-\overline P\Delta t\bigr]
  \end{aligned}
  \label{eq:app-lag}
\end{equation}

对各决策变量求平稳性条件（$E_t,s_t,c_t,d_t$ 对 $t=1,\dots,144$；$S_t$ 对 $t=1,\dots,143$，
因为日周期约束把 $S_{144}=S_0=6000$~kWh 固定）：

\begin{align}
  \frac{\partial\mathcal L}{\partial E_t}&=p_t-\alpha_t\ \ge 0,
  &E_t\bigl(p_t-\alpha_t\bigr)&=0 \label{eq:app-stat-E}\\
  \frac{\partial\mathcal L}{\partial s_t}&=\alpha_t\ \ge 0,
  &s_t\,\alpha_t&=0 \label{eq:app-stat-s}\\
  \frac{\partial\mathcal L}{\partial c_t}&=\alpha_t-\eta_c\nu_t+\varrho^{c}_t\ \ge 0,
  &c_t\bigl(\alpha_t-\eta_c\nu_t+\varrho^{c}_t\bigr)&=0 \label{eq:app-stat-c}\\
  \frac{\partial\mathcal L}{\partial d_t}&=-\alpha_t+\frac{\nu_t}{\eta_d}+\varrho^{d}_t\ \ge 0,
  &d_t\Bigl(-\alpha_t+\frac{\nu_t}{\eta_d}+\varrho^{d}_t\Bigr)&=0 \label{eq:app-stat-d}\\
  \frac{\partial\mathcal L}{\partial S_t}&=\nu_t-\nu_{t+1}+\lambda^{+}_t-\lambda^{-}_t=0
  \label{eq:app-stat-S}
\end{align}

由式~\eqref{eq:app-stat-E}，购电时段（$E_t>0$）恒有

\begin{equation}
  \alpha_t=p_t
  \label{eq:app-alpha-eq-p}
\end{equation}

即电量平衡的影子价格等于该时段的电价；数值上逐槽复核的最大偏差为 $0.000$~元/kWh。
若充放电功率界均不活跃（$\varrho^{c}_t=\varrho^{d}_t=0$），式~\eqref{eq:app-stat-c} 与
式~\eqref{eq:app-stat-d} 化为

\begin{equation}
  \alpha_t-\eta_c\nu_t\ \ge 0,\qquad
  \frac{\nu_t}{\eta_d}-\alpha_t\ \ge 0
  \label{eq:app-threshold-raw}
\end{equation}

且 $c_t>0$ 时前式取等、$d_t>0$ 时后式取等，即得正文命题~\bref{prop:threshold} 的阈值律：
\textbf{$\alpha_t\le\eta_c\nu_t$ 时充电、$\alpha_t\ge\nu_t/\eta_d$ 时放电}，
死区为 $\eta_c\nu_t<\alpha_t<\nu_t/\eta_d$，死区比为 $1/(\eta_c\eta_d)=1.2346$
（单边效率各 0.9、往返 0.81），与式~\bref{eq:q1-deadband} 一致。
由式~\eqref{eq:app-stat-S} 还可得水价值的跨期关系：在储电量上下界均不活跃的连续时段上
$\nu_t=\nu_{t+1}$（水价值恒定）；储电量触及上界时 $\lambda^{+}_t>0$、$\nu_t<\nu_{t+1}$；
触及下界时 $\lambda^{-}_t>0$、$\nu_t>\nu_{t+1}$。
这一关系是图~\bref{fig:dualprice} 中水价值曲线出现折点的原因。

数值复核（\texttt{intermediate/q1\_dual2.json}）：储电量双界均不活跃的槽位 97 个；
$\nu_t\in[0.4711,\ 1.1305]$~元/kWh，$\alpha_t\in[0.3713,\ 1.2561]$~元/kWh；
全部 144 槽的充放电阈值律违反次数为 $0$，互补松弛最大残差为 $0.000$~元/kWh，
式~\eqref{eq:app-alpha-eq-p} 的最大偏差为 $0.000$~元/kWh。

\subsection{报童临界分位数的推导}

设时段 $t$ 的需求 $D_t$ 为随机变量、分布函数为 $F_t$，计划购电量为 $E_t$，
过剩成本为 $c_o$（多买 1~kWh 的损失），缺货成本为 $c_u$（不足 1~kWh 的额外支出）。
期望总成本为

\begin{equation}
  \Pi(E)=c_o\,\mathbb E\bigl[(E-D)^{+}\bigr]+c_u\,\mathbb E\bigl[(D-E)^{+}\bigr]
  \label{eq:app-nv-cost}
\end{equation}

对 $E$ 求导并令其为零：

\begin{equation}
  \frac{\mathrm d\Pi}{\mathrm dE}
  =c_oF(E)-c_u\bigl[1-F(E)\bigr]=0
  \quad\Longrightarrow\quad
  F(E^{*})=\frac{c_u}{c_u+c_o}
  \label{eq:app-nv-foc}
\end{equation}

本题中不足 1~kWh 需按 $5p_t$ 紧急购电，而正常电价 $p_t$ 已随计划购电量支付，
故缺货的额外支出为 $c_u=5p_t-p_t=4p_t$；超购的电量若无处可用（超购即弃置），
则 $c_o=p_t$，代入式~\eqref{eq:app-nv-foc} 得

\begin{equation}
  F(E^{*})=\frac{4p_t}{4p_t+p_t}=\frac45=0.8
  \qquad\Longrightarrow\qquad
  \text{最优缺货概率}=1-F(E^{*})=\frac15=0.2
  \label{eq:app-nv-result}
\end{equation}

即「留安全边际」的临界分位数为 $4/5$、最优缺货概率为 $1/5$。
若超购电量能转入储能、留待后续更高价时段使用（单位回收价值 $r_t$），
则 $c_o=p_t-r_t$，临界分位数相应上移：取全年均价 $p=0.7662$~元/kWh 与全年最低价
$p^{low}=0.3713$~元/kWh，得 $c_o\approx0.395$~元/kWh，$F(E^{*})\approx0.886$，
最优缺货概率降至约 $1/10$。此即正文式~\bref{eq:q2-newsvendor} 与表~\bref{tab:newsvendor} 的来源。

实测对照（\texttt{paper\_numbers.json} 的「报童经验」）：主口径下逐槽缺货频率为 $8.55\%$，
缺口电量 $172\,015.82$~kWh，电量满足率 $99.04\%$，安全边际合计 $576\,299.67$~kWh
（$11.98$~kWh/槽），落在 $1/10$ 与 $1/5$ 两界\textbf{之下侧}——
因为储能同时承担了对冲功能，报童结论只作方向性近似与上下界使用，不作精确刻画。

\section{随机种子清单}

本工程\textbf{只在四处使用随机数}，四问的线性规划求解、主口径的场景削减、
审计与数值汇总均为确定性过程（同一输入必得同一输出）。逐项清单如表~\ref{tab:seeds}。

\begin{table}[htbp]
  \centering
  \caption{随机种子清单（代码中的全部随机源；其余过程确定）}
  \label{tab:seeds}
  \footnotesize
  \setlength{\tabcolsep}{4pt}
  \begin{tabularx}{\linewidth}{L{3.0cm} L{3.2cm} C{2.5cm} Y}
    \toprule
    用途 & 随机机制 & 种子 & 说明 \\
    \midrule
    问题 4 联合情景合成（条件对照） & 分块 Bootstrap，按时刻对齐 6~h 块重采样
      （\texttt{q4.build\_joint\_scenarios(mode=“boot”)}） & \texttt{numpy.random.default\_rng(seed=7)}
      & 电价块与光伏预报残差块成对抽取，保住日内形状与相关性；一次生成后固定使用 \\
    场景削减的对照路径 & $k$-means 初始化（\texttt{lp\_engine.reduce\_scenarios}）
      & \texttt{random\_state=0}（\texttt{n\_init=4}）
      & 仅在 \texttt{--scen kmeans} 对照下启用；\textbf{主口径不抽样}：按日总残差等频分层后在层内取 medoid，确定性 \\
    非预期性置换检验 & 对某次预报做日间随机置换（\texttt{exp\_na.py}）
      & \texttt{numpy.random.default\_rng(seed=0)}；对照用 \texttt{seed=1}
      & 检验「置换 $6{:}00$ 预报后槽 0--35 完全不变、槽 36--71 改变」，用于排除信息泄漏 \\
    全局灵敏度的采样 & Sobol 低差异序列置乱
      （\texttt{scipy.stats.qmc.Sobol(scramble=True)}，\texttt{exp\_sobol.py}）
      & 费用目标 \texttt{seed=42}；节省率目标 \texttt{seed=43}
      & 置乱只改变采样点的取法，不改变参数区间与模型本身 \\
    \bottomrule
  \end{tabularx}
\end{table}

\noindent\textbf{确定性的部分（无随机数）}：问题 1--4 的全部线性规划由 SciPy
\texttt{linprog(method=“highs”)} 精确求解；问题 2、3 的场景削减用「等频分层加层内 medoid」，
问题 4 的历史代表日情景用「日净需求残差与日均电价」二维分层取 medoid、按层样本占比赋权，
两者都是确定性算法；\texttt{audit.py}、\texttt{paper\_numbers.py} 与 \texttt{exp\_scan.py}
不含随机数。因此论文的每一个数字在同一环境下重跑即可复现。

\paragraph{诚实说明}：本工程\textbf{未开展「多种子重复生成情景」的稳健性实验}，
故不报告费用的均值、标准差与 $95\%$ 置信区间，也不列出 20 个种子。
本工程对这一缺口的处理是：把随机性限制在只影响未来情景、不影响已实现时段决策的合成与采样环节
（分块 Bootstrap 只作为问题 4 联合情景的条件对照，主口径用确定性的分层 medoid），
并以同源对照替代多种子重复——同一固定情景集下的信息结构对照（R1、R2、R3）、
场景数对照（$K=60$ 与 $K=120$）、终端条件对照与安全边际消融，四者的结论方向一致。
若需要多种子区间，只需改动 \texttt{build\_joint\_scenarios} 的 \texttt{seed} 参数并重跑
\texttt{run\_all.py}，其余环节无需修改。

\section{验证与灵敏度明细}

本节汇集正文第九章与第三章各表的完整明细，便于复核；正文已给出各表的结论与判据。

\subsection{非预期性（信息泄漏）置换检验的逐区明细}\label{sec:na-app}

把某一次发布的预报在 12 天之间随机置换后重解，观察四个控制区（槽 $0$--$36$、$36$--$72$、
$72$--$108$、$108$--$144$ 分别由 $0{:}00$、$6{:}00$、$12{:}00$、$18{:}00$ 预报控制）的最大变化量，
结果见表~\ref{tab:na-app}。

\begin{table}[htbp]
  \centering
  \caption{非预期性置换检验：置换某一个发布时刻的预报后，四个控制区的最大变化量
           （对象：单日两阶段模型；范围：四个控制区各 36 个 10~min 槽；单位 kWh；
           变化恰为 $0.000$ 表示该区决策不受该次预报影响）}
  \label{tab:na-app}
  \footnotesize
  \setlength{\tabcolsep}{5pt}
  \begin{tabular}{lcccc}
    \toprule
    被置换的预报 & 控制区 1（$0{:}00$--$6{:}00$） & 控制区 2（$6{:}00$--$12{:}00$）
                 & 控制区 3（$12{:}00$--$18{:}00$） & 控制区 4（$18{:}00$--$24{:}00$） \\
    \midrule
    \(0{:}00\) & 833.33 & 822.06 & 489.52 & 0 \\
    \(6{:}00\) & 0.000  & 858.80 & 833.33 & 0 \\
    \(12{:}00\)& 0.000  & 0.000  & 805.35 & 0 \\
    \(18{:}00\)& 0.000  & 0.000  & 0.000  & 0 \\
    \bottomrule
  \end{tabular}
  \\[1ex]
  {\footnotesize 判读：任何一次预报被置换，都不改变早于其发布时刻的决策（变化恰为 $0.000$），
   故不存在用未来信息决策的泄漏；而“置换 $0{:}00$ 预报会改变全天决策”说明检验本身有效力。
   $18{:}00$ 预报的零影响源于其覆盖的傍晚时段光伏实际与预报都接近 0、信息量极低。}
\end{table}


\subsection{电价过程的四项统计检验：完整统计量与判定}\label{sec:price-test-app}

正文 \S7.1 给出三条建模含义；本小节给出四项检验（T1--T4）的完整统计量、结果与逐项判定。

\begin{table}[htbp]
  \centering
  \caption{电价过程的四项统计检验（对象：附件 4 的逐时段电价与残差过程；范围：2025-01-01 至 12-31、
            逐 10~min 共 \(365\times144=52\,560\) 个观测；单位为量纲一的统计量）}
  \label{tab:price-test}
  \footnotesize
  \setlength{\tabcolsep}{3pt}
  \begin{tabularx}{\linewidth}{C{0.62cm} L{1.85cm} L{4.35cm} L{3.35cm} Y}
    \toprule
    编号 & 检验对象 & 检验方法与统计量 & 结果 & 对建模的含义 \\
    \midrule
    T1 & 确定性日内骨架是否存在
      & 逐时段均值与逐时段标准差的对比（变异系数）
      & 附件 1 剖面与附件 4 逐时段均值的最大偏差 \(5.15\times10^{-5}\)~元/kWh（全年 365 天口径）；
        残差逐时段标准差均值 0.1238~元/kWh，变异系数 16.94\%
      & 骨架稳定存在但残差不可忽略，采用“骨架 \(+\) 残差”两层结构 \\
    T2 & 残差的正态性与日间差异
      & Jarque--Bera 检验、偏度与峰度；同一槽位跨日自相关
      & \(\mathrm{JB}=656.70\)（\(p=2.51\times10^{-143}\)）拒绝正态；偏度 \(-0.169\)、峰度 \(-0.431\)；
        日间自相关 \(\mathrm{ACF}(1\text{--}3)=0.513/0.148/0.144\)
      & 残差非正态、日间相关性弱：不做正态参数化，改用非参数分层重采样；骨架不按月分段 \\
    T3 & 日内自相关与波动结构
      & 残差自相关 \(\mathrm{ACF}(1\text{--}6)\)；峰段与谷段残差标准差之比
      & \(\mathrm{ACF}(1\text{--}6)\) 为 0.8589、0.666、0.614、0.573、0.494、0.432；峰段（19:00--20:00）
        残差标准差 0.1634~元/kWh 为谷段（03:00--04:00）0.0580~元/kWh 的 \textbf{2.81 倍}
      & 不能按独立同分布抽样：情景须按时刻对齐分块，同时保留日内相关与按时段异方差 \\
    T4 & 与负载、光伏的相关性
      & Pearson 相关系数（池化与日聚合两个尺度）
      & 池化：与负载 \(+0.508\)、净负荷 \(+0.260\)、光伏 \(+0.028\)；日聚合：与负载 \(+0.898\)、
        光伏 \(-0.343\)
      & \textbf{决定必须生成电价与净需求的联合情景}，否则系统性低估“高价时段同时缺电”的尾部风险 \\
    \bottomrule
  \end{tabularx}
\end{table}
\subsection{单参数灵敏度与 Sobol 全局灵敏度的完整档位}\label{sec:sens-app}

正文第九章给出各参数的扰动档位与“结论是否改变”的判断；本小节补充 Sobol 全局灵敏度的完整
指数，并说明它为什么必须\textbf{区分输出量}。

\begin{table}[htbp]
  \centering
  \caption{Sobol 全局灵敏度指数（典型日模型，$N=256$；输出量 1：典型日全天购电费，
           均值 $34\,603.42$ 元；输出量 2：相对无储能基线的节省率，均值 $27.02\%$；
           指数为量纲一）}
  \label{tab:sobol-app}
  \footnotesize
  \setlength{\tabcolsep}{5pt}
  \begin{tabular}{lcc}
    \toprule
    参数 & 全天购电费的一阶指数 \(S_1\) & 节省率的一阶指数 \(S_1\) \\
    \midrule
    储电量上限 \(\overline S\)        & 0.2510 & \(\le0.0033\) \\
    储电量下限 \(\underline S\)      & 0.1099 & \(\le0.0033\) \\
    往返/单边效率 \(\eta_c,\eta_d\)  & 0.0695 & \(\le0.0033\) \\
    充放功率上限 \(\overline P\)     & 0.0263 & \(\le0.0033\) \\
    电价水平                         & 0 & 0.8464 \\
    负载水平                         & 0 & 0.4711 \\
    光伏出力水平                     & 0 & \(\le0.0033\) \\
    峰谷价差                         & 0 & \(\le0.0033\) \\
    \bottomrule
  \end{tabular}
  \\[1ex]
  {\footnotesize 读法：对\textbf{全天购电费}起主导作用的是储电量上下限与效率；
   对\textbf{节省率}起主导作用的却是电价水平与负载水平——在问题 1 的确定性模型中，
   电价、光伏与负载水平对前者只是整体缩放，不改变调度结构，故一阶指数为 0。
   结论是“哪一个参数决定哪一个输出量”，而非“所有参数都重要”。}
\end{table}

\subsection{容量与充放功率的双向扫描明细}\label{sec:scan-app}

正文第四章给出扫描结论（容量摆幅是主要瓶颈）；本小节给出完整档位与两个单自由度放开实验。

\begin{table}[htbp]
  \centering
  \caption{储能容量与充放功率上限的双向扫描（问题 1，典型日；费用单位 元；
           “相对主口径”以主口径最优费用 $35\,126.95$ 元为分母，正号表示更省；
           容量扫描同步移动储电量上下限以保持 10\%--90\% 放电深度）}
  \label{tab:scan-app}
  \footnotesize
  \setlength{\tabcolsep}{5pt}
  \begin{tabular}{rrr|rrr}
    \toprule
    容量（kWh） & 最优购电费 & 相对主口径 & 充放功率上限（kW） & 最优购电费 & 相对主口径 \\
    \midrule
    7\,200  & 38\,640.45 & \(-10.00\%\) & 1\,250  & 40\,166.17 & \(-14.35\%\) \\
    9\,600  & 36\,700.45 & \(-4.48\%\)  & 2\,500  & 36\,390.84 & \(-3.60\%\) \\
    12\,000 & \textbf{35\,126.95} & \(0.00\%\) & 5\,000 & \textbf{35\,126.95} & \(0.00\%\) \\
    14\,400 & 33\,875.57 & \(+3.56\%\)  & 7\,500  & 35\,027.13 & \(+0.28\%\) \\
    19\,200 & 31\,692.43 & \(+9.78\%\)  & 10\,000 & 34\,960.59 & \(+0.47\%\) \\
    28\,800 & 28\,854.53 & \(+17.86\%\) & ---     & ---        & --- \\
    48\,000 & 27\,601.80 & \(+21.42\%\) & ---     & ---        & --- \\
    \bottomrule
  \end{tabular}
  \\[1ex]
  {\footnotesize 另做两个单自由度放开实验：仅解除储电量上限（充放功率维持 $5000$~kW）为
   $27\,351.55$~元，相对主口径多省 $22.14\%$；仅解除充放功率上限（容量维持 $12\,000$~kWh）
   为 $34\,544.43$~元，多省 $1.66\%$。界活跃统计：充电功率界在 $19/144$ 个时段被打满、
   放电功率界 $0/144$，储电量上界活跃 $21/144$、下界活跃 $8/144$。}
\end{table}

\subsection{L 形（Benders）分解与直接求解的独立对拍}\label{sec:lshape}

问题 2--4 的模型都是\textbf{两阶段结构}：第一阶段是全场共享的计划量 $(E_t,c_t,d_t,S_t)$，
第二阶段是逐情景的追索量 $g_{\omega,t}$。本文主结果全部用 HiGHS 直接求解（单日 $4n+Kn$ 个变量，
$K=60$ 时 $9216$ 个）；但“求解器返回最优”本身不构成模型正确性的证据，因此本节\textbf{另写一套
分解算法}作为对照：把问题改写为 L 形（Benders）形式，用与直接求解\textbf{完全不同的数值路径}
再算一遍。给定第一阶段 $(E,c,d,S)$ 后，情景 $\omega$ 的追索问题逐槽完全可分：
\begin{equation}
  Q_\omega(E,c,d)=\sum_t 5p_t\max\bigl\{0,\ \widetilde{ND}_{\omega,t}+c_t-d_t-E_t\bigr\},
  \label{eq:lshape-sub}
\end{equation}
每个槽的最优值就是一个正部函数，对 $E_t$ 的次梯度为 $-5p_t$（缺口为正时）或 $0$（否则），
故\textbf{每个情景每次迭代可解析地给出 144 条最优性割}（多割），无需调用求解器求子问题；
主问题只保留 $4n+1=577$ 个变量。

\begin{table}[htbp]
  \centering
  \caption{L 形分解与直接求解的对照（对象：单日两阶段随机规划，题面指定四个代表日
           （2025-03-20、06-21、09-23、12-21）的结果合计；范围：$s_{init}=6000$~kWh、
           终端模式 $S_n\ge s_{init}$、两法共用同一份分层 medoid 情景；
           时间列单位为 s（四日合计），变量数单位为个；
           相对间隙 $=|z_{\text{L 形}}-z_{\text{直接}}|/|z_{\text{直接}}|$；
           数据源：\texttt{intermediate/decomp.json} 的 \texttt{scenario\_scan} 字段）}
  \label{tab:lshape}
  \small
  \setlength{\tabcolsep}{4pt}
  \begin{tabular}{crrrrrc}
    \toprule
    情景数 $K$ & 直接求解变量 & L 形变量 & 直接求解 & L 形分解 & 加速比 & 最大相对间隙 \\
    \midrule
    $20$  & $3\,456$ & $720$ & $0.142$ & $0.281$ & $0.50\times$ & $6.65\times10^{-16}$ \\
    $60$  & $9\,216$ & $720$ & $0.356$ & $0.358$ & $1.00\times$ & $2.09\times10^{-15}$ \\
    $120$ & $17\,856$ & $720$ & $1.291$ & $0.401$ & $3.22\times$ & $3.06\times10^{-15}$ \\
    $300$ & $43\,776$ & $720$ & $3.721$ & $0.635$ & $5.86\times$ & $2.60\times10^{-15}$ \\
    $500$ & $72\,576$ & $720$ & $8.962$ & $0.875$ & $10.25\times$ & $4.60\times10^{-15}$ \\
    \bottomrule
  \end{tabular}
\end{table}

\textbf{三条结论。} ① \textbf{两条数值路径给出的最优值一致到机器精度}：结果窗口全部 334 天上
（$K=60$）逐日对比最大相对间隙 $3.76\times10^{-15}$、全部收敛（平均迭代 9.75 次、
单日最多 1097 条割），$K=300$ 的 42 天子样本上为 $4.82\times10^{-15}$——两个独立实现同时给出
同一最优值，比“求解器说 optimal”强得多。② \textbf{加速比随情景数增长}：$K=20$ 时分解反而慢约
一倍，$K\approx60$ 附近是分界点，$K=120$、300、500 时分别加速 3.22、5.86、10.25 倍；
主问题变量数始终为 720、不随 $K$ 增长，这是加速的来源。③ \textbf{工程取舍}：本题主口径 $K=60$、
单日直接求解 0.106~s（L 形 0.092~s，两者同量级），分解法在该规模上\textbf{没有}速度优势；
把解析割聚合成单割后 200 次迭代仍未收敛、相对间隙 0.65\%--3.42\%，而多割 9--14 次即收敛，
故\textbf{主结果仍用直接求解}；分解法的价值在于给出独立最优性对拍，并在情景数放大到数百时
把单日求解稳定在 0.8~s 以内。\evid{C-Q5-01}







\begin{table}[htbp]
  \centering
  \caption{光伏预报偏差的逐整点结构（样本为 2025 年每天 \(0{:}00\) 发布的 24~h 预报与
           附件 2 对应实际功率的逐槽配对，\(n=52\,560\)；bias 单位 kW，
           “极值整点”给出该组内绝对偏差最大的整点及其取值）}
  \label{tab:biasprofile}
  \footnotesize
  \setlength{\tabcolsep}{4pt}
  \begin{tabular}{lccc}
    \toprule
    整点分组 & 整点范围 & 组内 bias 区间 & 极值整点 \\
    \midrule
    夜间（光伏无有效出力） & \(0\)--\(4\) 时 & \(0.00\) 至 \(+0.09\) & \(4\) 时 \(+0.09\) \\
    上午（系统性\textbf{高估}） & \(5\)--\(10\) 时 & \(+214.21\) 至 \(+733.00\) & \(7\) 时 \(+733.00\) \\
    正午过渡 & \(11\) 时 & \(+7.21\) & --- \\
    下午（系统性\textbf{低估}） & \(12\)--\(18\) 时 & \(-801.87\) 至 \(-106.46\) & \(15\) 时 \(-801.87\) \\
    傍晚之后（预报空洞） & \(19\)--\(23\) 时 & \(-0.15\) 至 \(0.00\) & \(19\) 时 \(-0.15\) \\
    \bottomrule
  \end{tabular}
\end{table}





\section{场景保真度与电价检验明细}

本节给出正文第五章场景集与第七章电价过程检验、CVaR 风险—费用前沿的完整明细表。

（联合情景保真度检验表见 \ref{tab:scenfid}。）



\begin{table}[htbp]
  \centering
  \caption{风险前沿三个代表点的情景内指标与真实轨迹回测
           （对象：联合情景集上按 \(\lambda\) 求解的三组解；范围：2025-02-01 至 12-31 中每 6 天取 1 天的 56 个代表日；
            单位：费用 元/日、电量 kWh/日；样本：64 个代表情景，\(\alpha=0.95\)）}
  \label{tab:frontier-backtest}
  \footnotesize
  \setlength{\tabcolsep}{4pt}
  \begin{tabularx}{\linewidth}{L{3.0cm} C{2.3cm} C{2.3cm} C{2.5cm} Y}
    \toprule
    代表点 & 情景内期望费用 & 情景内 95\% CVaR & 真实轨迹回测费用 & 紧急购电量 \\
    \midrule
    \(\lambda=0\)（期望最优） & 41\,995.59 & 52\,176.36 & 45\,694.27 & 480.19 \\
    \(\lambda=0.3\)（推荐） & 42\,168.43 & 50\,814.83 & 45\,778.03 & 324.60 \\
    \(\lambda=1\)（近似最坏） & 42\,411.56 & 50\,510.46 & 45\,952.85 & 313.83 \\
    \bottomrule
  \end{tabularx}
\end{table}

\section{验证映射、风险结果与局限明细}

本节给出正文第九章验证层级映射、第八章问题 4 逐日期结果与第十章局限清单的完整表格。





\section{附件结构与问题 4 逐日期结果明细}

本节给出四个附件的字段结构与问题 4 在四个指定日期的逐项结果。



【问题 4 逐日期结果】




\section{明细表与口径框（正文摘要引用）}
本节集中给出正文以摘要形式引用的明细表与全部核心指标口径框（\calref{C-01}--\calref{C-11}），便于评审核对；正文对应处均给出结论、口径编号与本节的表号。













\begin{caliber}{C-02}{全天购电量}
\textbf{定义}\quad \(E^{\text{day}}=\sum_{t=1}^{144}E_t\)，即该日 144 个 10~min 槽的\textbf{计划购电量}之和；%
\textbf{公式}\quad 结果为逐槽求和，无闭式；结果文件每行末列即该日汇总；%
\textbf{数据}\quad 结果文件「计划购电量」工作表的 144 个时段列；%
\textbf{时段}\quad 单日；问题 1 为典型日，问题 2--4 为 2025-02-01 至 12-31 的 334 天；%
\textbf{单位}\quad kWh（结果文件保留 4 位小数，正文保留 2--4 位）；%
\textbf{含}\quad 问题 2、4-2 中即第一阶段计划量；问题 3、4-3 中调整后的量一律称“全天调整购电量”并单独命名；%
\textbf{不含}\quad 不含储能充电量、不含紧急购电量 \calref{C-05}、不含盈余处置量 \calref{C-06}；%
\textbf{对照}\quad 问题 2、3 的计划购电量模型内值为 \(22\,220\,333.2033\)~kWh，结果文件回读值为 \(22\,220\,333.1391\)~kWh，
相对差 \(2.9\times10^{-9}\)，来自写入时的 4 位小数取整；%
\textbf{复算}\quad \texttt{results\_io.py::plan\_sheet\_frame()}（写出行汇总列）与 \texttt{audit.py::audit\_result\_files()}（回读结果文件后重算 \(\sum_{144}\) 槽）
\end{caliber}

\begin{caliber}{C-03}{总购电费用（计划 \(+\) 紧急 \(+\) 调整相关）}
\textbf{定义}\quad \(C^{\text{total}}=C^{\text{plan}}+C^{\text{breach}}+C^{\text{excess}}+C^{\text{emg}}\)；
问题 2、4-2 无调整机制，后两项恒为 0；%
\textbf{公式}\quad 问题 2、4-2：\(C=\sum_tp_tE_t+5\sum_tp_tg_t\)；
问题 3、4-3：\(C^{\text{plan}},C^{\text{breach}},C^{\text{excess}}\) 按 \calref{C-09} 的偏差量定价口径计算，
\(C^{\text{emg}}=\sum_t5p_tg_t\)；%
\textbf{数据}\quad 电价取附件 1（问题 1--3）或附件 4（问题 4）；电量取本队结果文件；%
\textbf{时段}\quad 问题 1 为典型日；问题 2--4 为 334 天；%
\textbf{单位}\quad 元（保留 2 位小数，内部计算不取整）；%
\textbf{不含}\quad 不含储能投资与运维成本、容量补偿、辅助服务收益、绿电环境价值、循环衰减与安全成本；%
\textbf{对照}\quad 若只报计划购电费而不计紧急购电，问题 2 为 \(13\,637\,552.90\)~元，
比 \(C^{\text{total}}\) 低 \(627\,239.24\)~元（\(4.40\%\)）；%
\textbf{复算}\quad \texttt{q2.py::run\_q2()}、\texttt{q3.py::summary()}、\texttt{q4.py::run\_q4\_2()} 与 \texttt{audit.py::audit\_result\_files()}，
四项之和与总费用的偏差为 0
\end{caliber}

\begin{caliber}{C-04}{节省率}
\textbf{定义}\quad \(\eta=\bigl(C^{base}-C^{\text{total}}\bigr)/C^{base}\times100\%\)，分母与分子必须同窗口、同电价口径；%
\textbf{公式}\quad \(C^{base}\) 取 \calref{C-01} 的对应变体（V1 典型日 / V2 固定价窗口 / V3 波动价窗口，
变体定义见正文表~\bref{tab:cbases}）；%
\textbf{数据}\quad \(C^{base}\) 与 \(C^{\text{total}}\) 均取自本队结果；%
\textbf{时段}\quad 与 \calref{C-01} 变体一致；%
\textbf{单位}\quad \%（保留 2 位小数）；%
\textbf{不含}\quad 不得相对贪心策略或完全信息下界报同一个数而不加限定；相对其他基线的结果必须\textbf{单独命名}；%
\textbf{对照}\quad 问题 1 相对 V3 为 \(26.90\%\)；问题 2、3 相对 V1 为 \(13.06\%\)、\(15.08\%\)；
问题 4-2、4-3 相对 V2 为 \(11.35\%\)、\(13.54\%\)；%
\textbf{复算}\quad \texttt{paper\_numbers.py::main()}（由 \calref{C-01} 与 \calref{C-03} 导出）
\end{caliber}

\begin{caliber}{C-06}{弃光/盈余处置电量}
\textbf{定义}\quad 光伏盈余与超购部分的处置量：典型日按
\(s_t=PV_t\Delta t+E_t+d_t-L_t\Delta t-c_t\ge0\) 逐槽统计（式~\bref{eq:q1-balance} 的松弛量）；
窗口口径按实际数据复算为 \(\sum_t\max\{PV_t\Delta t+E_t+d_t+g_t-L_t\Delta t-c_t,0\}\)；%
\textbf{公式}\quad 上式为定义式；该量非负，其余额为超购部分被弃置的电量；%
\textbf{数据}\quad \(L,PV\) 取附件 2 实际值，其余取本队结果文件；%
\textbf{时段}\quad 逐槽统计、按窗口（334 天）汇总；%
\textbf{单位}\quad kWh（按电量统计，\textbf{不得}按功率报）；%
\textbf{不含}\quad 不含“因充电功率受限而未发生的充电”，不含未购电部分；%
\textbf{对照}\quad 问题 1 典型日最优弃光 \(=0.0000\)~kWh（数值残差 \(1.37\times10^{-12}\)~kWh）；
问题 2、3 的窗口盈余处置量为 \(3\,106\,390.84\)~kWh 与 \(2\,638\,963.35\)~kWh；
问题 4-2、4-3 的同一量为 \(3\,100\,463.79\)~kWh 与 \(2\,645\,427.81\)~kWh（均见正文表~\bref{tab:summary}）；%
\textbf{复算}\quad \texttt{config.py::curtailment(price, load, pv, E, c, d)}（典型日）与 \texttt{audit.py::audit\_internal()}（窗口口径）
\end{caliber}

\begin{caliber}{C-07}{充放电量与储电量（交流侧）}
\textbf{定义}\quad 结果文件中的 \(c_t,d_t\) 为\textbf{并网点（AC 侧）}电量；储电量为\textbf{时段末}值 \(S_t\)；%
\textbf{公式}\quad 递推见式~\bref{eq:q1-soc}，\(\eta_c=\eta_d=0.9\)（单边口径，往返 \(0.81\)）；
充放功率上限 \(5000\)~kW 对应单槽上限 \(\overline P\Delta t=833.3333\)~kWh；%
\textbf{数据}\quad 本队结果文件「充放电量」工作表；%
\textbf{时段}\quad 逐槽；表 2 格式按 6 个 4~h 块（每块 24 槽）汇总，并给出 \(0{:}00\) 与 \(24{:}00\) 两个储电量点；%
\textbf{单位}\quad kWh；运行区间 \([1200,10800]\)~kWh；%
\textbf{不含}\quad 不含电池侧（DC）电量、不含自放电与循环衰减（假设 A4）；%
\textbf{对照}\quad 恒等式校验：问题 1 的充放比 \(20\,740.67/16\,799.94=1.2346=1/(\eta_c\eta_d)\)，
与 \(\eta_c\eta_d=0.81\) 互为倒数，偏差 \(4.44\times10^{-16}\)；问题 2 四个指定日期的合计充放比同为 \(1.2346\)；%
\textbf{复算}\quad \texttt{results\_io.py::soc\_sheet\_frame()}（块级汇总）与 \texttt{audit.py::audit\_result\_files()}（块级净递推）
\end{caliber}

\begin{caliber}{C-08}{储电量递推与终端条件}
\textbf{定义}\quad 储电量\textbf{跨日连续递推}（\(S_{d,0}=S_{d-1,144}\)，假设 A5），
窗口末端施加终端条件 \(S_T\ge S_0\)（假设 A6），问题 1 另加严格日周期 \(S_0=S_{144}\)；%
\textbf{公式}\quad \(S_t=S_{t-1}+\eta_cc_t-d_t/\eta_d\)，\(S_t\in[1200,10800]\)~kWh；
2025 年 1 月为预热期（假设 A8），其终值即 2 月 1 日初值；%
\textbf{数据}\quad 结果文件「充放电量」工作表的储电量列与中间结果 \texttt{q2\_main.json} 等；%
\textbf{时段}\quad 逐槽；窗口 2025-02-01 至 12-31（334 天）；%
\textbf{单位}\quad kWh；%
\textbf{不含}\quad 不含“每日清零”“每日重置为 6000”两种读法——前者会引入末日掏空型短视解，后者与 1 月预热期矛盾；%
\textbf{对照}\quad 问题 1：日周期口径 \(35\,126.9486\)~元（\(S_0=S_{144}=6000.00\)~kWh）；
仅要求两端相等的自由环状口径 \(35\,118.5986\)~元（最优 \(S=8550.00\)~kWh），相差 \(8.35\)~元（\(0.02\%\)）；
窗口口径下“末值低于初值的天数”为 0；%
\textbf{复算}\quad \texttt{lp\_engine.py::solve\_day()}（终端约束）与 \texttt{audit.py::audit\_result\_files()}（\texttt{A6\_末值低于初值天数}）
\end{caliber}



\section{问题 2--4 的逐时段结果明细表}
第八章已按题面表 1／表 2／表 3 的格式给出四个指定日期的结果；本节补充\textbf{指定 10~min 时段}的购电量明细与逐时段充放电明细，便于与结果文件逐格对照。

\begin{table}[htbp]
  \centering
  \caption{微网在指定时段与指定日期的购电量、全天购电量和购电费（问题 2；电量单位：\unit{kWh}，费用单位：\unit{元}）}
  \label{tab:q2-buy}
  \small
  \begin{tabular}{lrrrr}
    \toprule
    时段／指标 & 2025-03-20 & 2025-06-21 & 2025-09-23 & 2025-12-21 \\
    \midrule
    \(10{:}00\)--\(10{:}10\) & 0.0000 & 0.0000 & 0.0000 & 0.0000 \\
    \(12{:}00\)--\(12{:}10\) & 804.3029 & 0.0000 & 536.2415 & 997.0335 \\
    \(14{:}00\)--\(14{:}10\) & 0.0000 & 0.0000 & 0.0000 & 0.0000 \\
    \(16{:}00\)--\(16{:}10\) & 817.3061 & 255.2254 & 902.7969 & 1001.2240 \\
    \(18{:}00\)--\(18{:}10\) & 697.9973 & 0.0000 & 764.3298 & 715.4144 \\
    \(20{:}00\)--\(20{:}10\) & 0.0000 & 0.0000 & 0.0000 & 0.0000 \\
    \midrule
    全天购电量 & 74\,178.1298 & 35\,701.5391 & 74\,744.7963 & 99\,692.6601 \\
    全天购电费 & 45\,229.90 & 20\,324.26 & 47\,063.39 & 64\,376.98 \\
    紧急购电量 & 1025.8805 & 187.0843 & 193.3669 & 22.5250 \\
    全天总费用 & 49\,654.02 & 21\,255.45 & 47\,717.70 & 64\,432.28 \\
    \bottomrule
  \end{tabular}
  \\[1ex]
  {\footnotesize 表注：时段按题面表述取区间左端时刻；
  全天购电费为计划口径（\(\sum_t p_tE_t\)），全天总费用为计划购电费与紧急购电费之和；
  四个日期的无储能基线分别为 52\,290.65、29\,668.18、54\,741.30、71\,176.92~\unit{元}。}
\end{table}

\begin{table}[htbp]
  \centering
  \caption{储能设备在指定日期的充放电量及 \(0{:}00\) 与 \(24{:}00\) 的储电量（问题 2，单位：\unit{kWh}）}
  \label{tab:q2-soc}
  \footnotesize
  \begin{tabular}{L{2.0cm}L{2.1cm}rr@{\quad}L{2.0cm}L{2.1cm}rr}
    \toprule
    日期 & 时段 & 充电量 & 放电量 & 日期 & 时段 & 充电量 & 放电量 \\
    \midrule
    2025-03-20 & \(0{:}00\)--\(4{:}00\)   & 4\,500.0000 & 0.0000    & 2025-09-23 & \(0{:}00\)--\(4{:}00\)   & 4\,500.0000 & 0.0000 \\
    2025-03-20 & \(4{:}00\)--\(8{:}00\)   & 833.3333   & 6\,224.1108 & 2025-09-23 & \(4{:}00\)--\(8{:}00\)   & 833.3333   & 5\,547.8264 \\
    2025-03-20 & \(8{:}00\)--\(12{:}00\)  & 6\,500.0000 & 2\,415.8892 & 2025-09-23 & \(8{:}00\)--\(12{:}00\)  & 5\,833.3333 & 3\,020.9227 \\
    2025-03-20 & \(12{:}00\)--\(16{:}00\) & 6\,905.5931 & 2\,218.5304 & 2025-09-23 & \(12{:}00\)--\(16{:}00\) & 6\,728.4081 & 1\,606.2615 \\
    2025-03-20 & \(16{:}00\)--\(20{:}00\) & 0.0000    & 5\,709.7835 & 2025-09-23 & \(16{:}00\)--\(20{:}00\) & 0.0000    & 5\,602.9091 \\
    2025-03-20 & \(20{:}00\)--\(24{:}00\) & 5\,333.3333 & 2\,930.2165 & 2025-09-23 & \(20{:}00\)--\(24{:}00\) & 5\,333.3333 & 3\,037.0909 \\
    \midrule
    2025-06-21 & \(0{:}00\)--\(4{:}00\)   & 0.0000    & 0.0000    & 2025-12-21 & \(0{:}00\)--\(4{:}00\)   & 4\,500.0000 & 0.0000 \\
    2025-06-21 & \(4{:}00\)--\(8{:}00\)   & 140.7851  & 4\,303.4721 & 2025-12-21 & \(4{:}00\)--\(8{:}00\)   & 1\,666.6667 & 1\,508.3333 \\
    2025-06-21 & \(8{:}00\)--\(12{:}00\)  & 5\,207.8672 & 130.5638 & 2025-12-21 & \(8{:}00\)--\(12{:}00\)  & 5\,833.3333 & 7\,806.6667 \\
    2025-06-21 & \(12{:}00\)--\(16{:}00\) & 5\,458.7995 & 0.0000    & 2025-12-21 & \(12{:}00\)--\(16{:}00\) & 8\,333.3333 & 2\,835.0000 \\
    2025-06-21 & \(16{:}00\)--\(20{:}00\) & 0.0000    & 6\,154.1006 & 2025-12-21 & \(16{:}00\)--\(20{:}00\) & 0.0000    & 5\,797.4742 \\
    2025-06-21 & \(20{:}00\)--\(24{:}00\) & 5\,333.3333 & 2\,485.8994 & 2025-12-21 & \(20{:}00\)--\(24{:}00\) & 5\,333.3333 & 2\,842.5258 \\
    \bottomrule
  \end{tabular}
  \\[1ex]
  {\footnotesize 表注：四个指定日期的 \(0{:}00\) 与 \(24{:}00\) 储电量均为 6000.0000~\unit{kWh}；
  每段含 24 个 10~\unit{min} 槽，充放电量为交流侧电量（\calref{C-07}）；
  四个日期的放电量／充电量之比均等于往返效率 \(\eta_c\eta_d=0.8100\)，
  其倒数即充电量／放电量之比 \(1/(\eta_c\eta_d)=1.2346\)。
  本表与正文表~\bref{tab:q2-t2} 取自同一数据源 \texttt{results/result2.xlsx}，逐格一致。}
\end{table}

\begin{table}[htbp]
  \centering
  \caption{微网在指定日期的紧急购电量（问题 2；购电量单位：\unit{kWh}，电价单位：\unit{元/kWh}）}
  \label{tab:q2-emergency}
  \footnotesize
  \begin{tabular}{llrr}
    \toprule
    日期 & 购电时间段 & 紧急购电量 & 交易时刻电价（加权平均） \\
    \midrule
    2025-03-20 & \(6{:}00\)--\(8{:}00\)   & 228.9914 & 0.8625 \\
    2025-03-20 & \(8{:}10\)--\(8{:}30\)   & 49.1586  & 0.8625 \\
    2025-03-20 & \(8{:}50\)--\(10{:}30\)  & 531.2739 & 0.8625 \\
    2025-03-20 & \(10{:}40\)--\(11{:}40\) & 216.4567 & 0.8625 \\
    \midrule
    2025-06-21 & \(4{:}00\)--\(4{:}50\)   & 5.2416   & 0.9955 \\
    2025-06-21 & \(14{:}00\)--\(14{:}30\) & 181.8427 & 0.9955 \\
    \midrule
    2025-09-23 & \(6{:}50\)--\(7{:}40\)   & 132.5948 & 0.6768 \\
    2025-09-23 & \(11{:}00\)--\(11{:}40\) & 60.7721  & 0.6768 \\
    \midrule
    2025-12-21 & \(11{:}00\)--\(11{:}30\) & 22.5250 & 0.4910 \\
    \midrule
    \multicolumn{2}{l}{四个日期合计} & 1428.8567 & --- \\
    \bottomrule
  \end{tabular}
  \\[1ex]
  {\footnotesize 表注：本表按“日期 \(\times\) 紧急购电事件”逐行展开，与题面表 3 的横向版式一一对应；
  同一日期可有多行，连续触发的相邻槽按事件聚合，时间段为该事件涉及的槽位区间；
  交易时刻电价为该事件内按电量加权的平均电价（事件电价按 \(5p_t\) 计价，本列为 \(p_t\)）。
  无紧急购电的时段不列入本表（全窗口 334 天中 296 天有触发）。}
\end{table}

\section{口径声明与时段对照明细表}




\section{伪验证自查清单（逐条证据）}\label{sec:fakeverify-app}

\noindent
\textbf{① 自证循环}：审计脚本不共用建模费用函数，费用双路径复算偏差 \(4.32\times10^{-7}\)；\\
\textbf{② 以求解器状态代替物理校验}：求解器返回 optimal 只作必要条件，同时报告供给缺口、储电量递推、
越界与互斥等物理残差；\\
\textbf{③ 比较不同样本}：R1／R2／R3、\(n=1/2/4\)、\(K=60/120\)、四种情景生成器均在同一场景集或同一
实际轨迹上配对比较；\\
\textbf{④ 训练集当测试集}：主结果明确声明为样本内评估，另做样本外回测（\bref{sec:oos}）与生成器样本外对照（\ref{sec:q4-gen}）；\\
\textbf{⑤ 只报点估计}：关键量均给多档对照——效率 4 档、容量 7 档、功率 5 档、\(\kappa\) 5 档、
场景数 2 档、CVaR 的 \(\lambda\) 11 档、机会约束 \(\beta\) 6 档、盒式鲁棒 \(k\) 7 档；\\
\textbf{⑥ 事后挑选指标}：本节指标在建模前由口径 C-01 至 C-11 固定，未随结果调整；\\
\textbf{⑦ 把“结果好看”当“模型正确”}：问题 2 的 \(13.06\%\) 与问题 1 的 \(26.90\%\) 同量级且更小，
差额可由预报误差成本 \(853\,231.23\)~元解释，未把“节省率更高”当作模型更好的证据；\\
\textbf{⑧ 扰动无效}：灵敏度扰动作用于运行区间与功率上限本身，且每档重新求解，不复用已有解。

\section{附图}

\begin{figure}[htbp]
  \centering
  \includefig[width=0.86\linewidth]{fig12_scenario_generators.pdf}
  \caption{四种情景生成器的样本外对照（左：相对“历史代表日”的总费用差；
           中：184 天累计紧急购电量；右：日总残差的 1\%/99\% 分位最大绝对偏差；
           对象：训练窗 2025-02-01 至 06-30、测试窗 07-01 至 12-31 共 184 天；
           单位：百分比、kWh；样本：情景数由二维分层固定为 48 个，Copula 生成器抽样 600 条后削减）}
  \label{fig:gen-cmp}
\end{figure}
；需要合并版请编译 full.tex。

\end{document}
